\documentclass[11pt,reqno]{amsart}
\usepackage[UTF8]{ctex}
\usepackage[T1]{fontenc}
\usepackage{amsmath,amssymb,amsthm,mathtools}
\usepackage{dsfont}
\usepackage[margin=1in]{geometry}
\usepackage{microtype}
\usepackage[colorlinks=true,allcolors=blue]{hyperref}
\numberwithin{equation}{section}
\newtheorem{theorem}{定理}[section]
\newtheorem{lemma}[theorem]{引理}
\newtheorem{proposition}[theorem]{命题}
\newtheorem{corollary}[theorem]{推论}
\theoremstyle{definition}\newtheorem{definition}[theorem]{定义}
\theoremstyle{remark}\newtheorem{remark}[theorem]{注}
\renewcommand{\proofname}{证明}
\renewcommand{\abstractname}{摘要}
\renewcommand{\refname}{参考文献}
\renewcommand{\datename}{\textit{日期}：}
\newcommand{\poly}{\operatorname{poly}}
\newcommand{\Prob}{\mathbb P}
\DeclareMathOperator{\Bin}{Bin}
\allowdisplaybreaks[2]
\renewcommand{\le}{\leqslant}
\renewcommand{\ge}{\geqslant}
\setlength{\emergencystretch}{1.5em}
\setcounter{tocdepth}{1}
\title{指数时间假设不成立}
\author{OpenAI}
\date{2026年10月17日，2026年10月25日修订}
\hypersetup{pdftitle={ETH 不成立},
 pdfauthor={OpenAI},
 pdfsubject={精确计数、稀疏奇偶检验与多项式插值}}
\begin{document}
\begin{abstract}
\normalsize
我们给出一个确定性算法，用于精确计算任意 $k$-CNF 公式的满足赋值个数。
对每个固定的 $k\ge2$，在含 $n$ 个变量与 $m$ 个子句的公式上，
我们的算法使用 $\poly_k(n,m)$ 空间，运行时间为
$(n+m)^{O_k(1)}2^{\,n/(\log n)^{\frac12+o(1)}}$。
我们的算法否证了指数时间假设（ETH），该假设断言
在 $n$ 个变量上判定 $3$-SAT 在最坏情况下需要
$2^{\Omega(n)}$ 时间。
\end{abstract}
\maketitle

\tableofcontents
\vspace{-\baselineskip}

\section{引言}
对每个固定的 $k\ge3$，$k$-SAT 是典型的 NP 完全
问题~\cite{Cook71}，其指数时间复杂性构成了指数时间
假设（ETH）与强指数时间假设
（SETH）~\cite{IP01,IPZ} 的基础。
本文中，我们通过为每个固定的 $k\ge2$ 给出一个
确定性算法来否证 ETH，该算法在亚指数时间与多项式空间内
计算 $k$-CNF 的满足赋值个数。

将输入 $k$-CNF 写在 $n$ 个声明变量上、含 $m$ 个输入子句，记为
\[
 F(x_1,\ldots,x_n)=\bigwedge_{e=1}^m
 \bigl[(x_{i_{e,1}}\oplus b_{e,1})\vee\cdots\vee
       (x_{i_{e,k}}\oplus b_{e,k})\bigr],
\]
其中 $i_{e,j}\in\{1,\ldots,n\}$ 与 $b_{e,j}\in\{0,1\}$ 分别是
第 $e$ 个子句中第 $j$ 个字面量的变量下标与否定位，
$\oplus$ 表示模二加法。重复字面量允许
子句的字面量少于 $k$ 个。
下标 $i_{e,j}$ 与位 $b_{e,j}$ 由 $F$ 固定；
只有 $x_1,\ldots,x_n$ 在赋值上变化。
我们的目标是计算量
\begin{equation}\label{eq:solution-count}
 \#F=\bigl|\{x\in\{0,1\}^n:F(x)=1\}\bigr|.
\end{equation}
检验 $\#F>0$ 同样判定可满足性。
我们让 $k$ 在 $n$ 增长时保持固定。
全文中 $\log$ 表示自然对数。
\begin{theorem}\label{thm:main}
对每个固定的 $k\ge2$，存在一个确定性算法：给定 $n\ge20$ 个声明变量上、含 $m$ 个子句的 $k$-CNF $F$，
它以
\[
 (n+m)^{O_k(1)}2^{O_k(n\log\log n/\sqrt{\log n})}
\]
位运算与 $(n+m)^{O_k(1)}$ 位空间计算 \eqref{eq:solution-count} 中定义的 $\#F$。
对每个固定的 $D>0$，它在满足 $m\le Dn$ 的输入上的运行时间为
\[
 (n+m)^{O_{k,D}(1)}2^{O_{k,D}(n/\sqrt{\log n})}.
\]
\end{theorem}

第~\ref{subsec:general-reduction} 节中的有理插值论证
给出较弱的指数 $O_k(\rho(n))$，其中
\[
 \rho(n)=\frac{n(\log\log n)^2}{\sqrt{\log n}}
          =\frac{n}{(\log n)^{\frac12+o(1)}}.
\]
附录~\ref{app:quantitative} 通过张量收缩证明
定理~\ref{thm:main} 中的两个界。

当 $\#F>0$ 时，标准的自归约通过每次固定一个变量并检验
剩余计数是否为正来找到一个满足赋值。若改为依据两个
条件计数来选择每一位，则得到恰好均匀随机的满足
赋值，其期望额外开销为 $\poly_k(n,m)$~\cite[定理~3.3(1)]{JVV86}。

该定理还给出 Max-$k$-SAT 的任意精度亚指数时间
近似算法。
Unique Games 定理~\cite[第~8.1 节]{OpenAIUGC26} 蕴含
Max-$k$-SAT 的 NP 困难性界，这些界在近似比上以任意小的误差
匹配~\cite{Rag08} 的半正定规划算法的保证。
记 $\operatorname{OPT}(F)$ 为 $F$ 中被一个赋值同时满足的
子句的最大数目。
\begin{corollary}\label{cor:max-sat-approximation}
对每个固定的 $k\ge2$ 与 $\varepsilon>0$，存在一个随机化
算法：给定 $n\ge20$ 个变量上、含 $m$ 个子句的 $k$-CNF $F$，
它以概率 $1-2^{-\Omega(n)}$ 找到满足至少
$\operatorname{OPT}(F)-\varepsilon m$ 个子句的赋值，运行时间为
\[
 (n+m)^{O_{k,\varepsilon}(1)}2^{O_{k,\varepsilon}(\rho(n))}
\]
空间为 $(n+m)^{O_{k,\varepsilon}(1)}$。
\end{corollary}
\begin{proof}
不妨设 $m>0$。
令 $f(x)$ 为 $x$ 所满足的子句比例，$\widehat f(x)$ 为在独立均匀
有放回选取的 $O(n/\varepsilon^2)$ 个子句样本中的该比例。
Hoeffding 不等式以及对所有 $2^n$ 个赋值的联合界
给出，以概率 $1-2^{-\Omega(n)}$，
\[
 |\widehat f(x)-f(x)|\le\varepsilon/2
 \qquad\text{对每个 }x\in\{0,1\}^n.
\]
因此，若 $\widehat x$ 使 $\widehat f$ 最大，则
$f(\widehat x)\ge\max_x f(x)-\varepsilon$。
给定整数阈值 $b$，一个每门至多两个输入、规模为 $O_{k,\varepsilon}(n)$ 的
布尔电路检验一个赋值是否满足至少 $b$ 个抽样得到的子句。
标准的电路到 $3$-CNF 归约~\cite[引理~15.9--15.10]{Bar23}
产生一个含 $O_{k,\varepsilon}(n)$ 个变量与子句的公式，它可满足
当且仅当该阈值可以达到。
定理~\ref{thm:main}、对阈值的二分搜索与自归约随后在所述界内
找到精确的样本最优解。
将拒绝采样在每个样本处截断于 $O(n+\log(n/\varepsilon^2))$
次试验，得到最坏情形 $\poly(n,m,\varepsilon^{-1})$
的采样时间，并使失败概率至多增加 $2^{-\Omega(n)}$。
\end{proof}

注意形如 $2^{o(n)}$ 的运行时间界在一般多项式时间归约下不会传递（为此需要 $2^{n^{o(1)}}$），因此我们的结果并不意味着一般 NP 问题允许亚指数时间算法。不过，我们确实得到了计数图着色与完美匹配的亚指数时间算法。

\begin{corollary}\label{cor:graph-counting}
在含 $v$ 个顶点与 $e$ 条边的简单图上，独立集的数目，以及对每个固定的 $q\ge2$，正常
$q$-着色的数目，都可以确定性地在 $2^{o(v)}$ 时间与
$\poly_q(v,e)$ 空间内计算。完美匹配的数目可以
确定性地在 $2^{o(v+e)}$ 时间与 $\poly(v,e)$ 空间内计算。
\end{corollary}
\begin{proof}
对 $G=(V,E)$，为每个 $u\in V$ 引入一个布尔变量 $x_u$。
\[
 \bigwedge_{\{u,w\}\in E}(\neg x_u\vee\neg x_w)
\]
的满足赋值与独立集 $\{u\in V:x_u=1\}$ 一一对应。
对正常 $q$-着色，为每个顶点与颜色的组合使用一个布尔变量，
要求每个顶点恰有一种颜色，并禁止相邻顶点取相同
颜色。这给出一个保持计数的
$q$-CNF 编码，含 $qv$ 个变量与 $O_q(v+e)$ 个子句。

对完美匹配，为每条边使用一个布尔变量，并用一个
线性规模的电路检验每个顶点恰与一条被选中的边相接。
标准的电路到 $3$-CNF
归约~\cite[引理~15.9--15.10]{Bar23} 给出 $O(v+e)$ 个变量
与子句。一旦边上的位固定，门变量便具有唯一的值，
因此该编码也保持计数。
定理~\ref{thm:main} 给出所述时间与空间界。
\end{proof}

定理~\ref{thm:main} 的固定密度情形给出
$\mathrm E^{\mathrm{NP}}$ 的电路下界，该类是
在确定性时间 $2^{O(n)}$ 内借助 NP 谕示可判定的语言类。

\begin{corollary}\label{cor:quantitative-circuits}
令 $C_L(n)$ 为在 $\{0,1\}^n$ 上计算 $L$ 的电路中非输入门的最少数目，
允许任意扇入至多为二的布尔门与无限制的扇出。存在语言
$L\in\mathrm E^{\mathrm{NP}}$ 与常数 $c>0$，使得
\begin{equation}\label{eq:quantitative-circuit-lower-bound}
 C_L(n)>c\,n\sqrt{\log n}\qquad\text{对无穷多个 }n.
\end{equation}
\end{corollary}

附录~\ref{subsec:circuit-lower-bounds} 证明该推论。
附录~\ref{subsec:local-inversion} 给出具有常数输出局部性的布尔映射的
反演与精确原像计数算法。

\subsection{相关工作}\label{sec:related}
我们首先讨论早期的 SAT 与计数算法，以及
定理~\ref{thm:main} 所否证的复杂性假设。
随后我们将自己的方法与早期的代数算法作比较，
并指出我们证明中使用的既有构造，它们的角色
在第~\ref{sec:technical-overview} 节的路线图中展开。

\subsubsection{SAT 算法与指数时间假设}\label{subsec:related-sat}
Sch\"oning 的局部搜索算法与 PPSZ 是 $k$-SAT 指数时间算法
发展中的里程碑~\cite{Schon99,PPSZ05}。
Sch\"oning 的随机化算法从随机赋值出发，
反复翻转某个未满足子句中均匀选取的变量。其
有界错误判定时间为
\[
 (n+m)^{O(1)}(2-2/k)^n
\]
对每个固定的 $k\ge3$~\cite[第~2 节]{Schon99}。
特别地，这给出 $3$-SAT 的 $(n+m)^{O(1)}(4/3)^n$ 时间。
PPSZ 算法~\cite{PPSZ05} 使用不同的机制：有界消解导出
额外的子句，随后按随机顺序为变量赋值，在可能时取被迫值，
否则进行猜测。对 $k=3,4$，其原始分析对唯一可满足公式
给出了比一般公式更强的界。
唯一 SAT 的保证后来被推广到一般 $3$-SAT 与
$4$-SAT；对 $3$-SAT，这给出随机化
$O(1.308^n)$ 界~\cite{Hertli14}。
Hansen、Kaplan、Zamir 与 Zwick 通过对算法的猜测引入偏置，
改进了每个固定 $k\ge3$ 的 PPSZ 界~\cite{HKZZ19}。
Scheder 通过细化原始 PPSZ 算法的分析，在同一 $k$ 范围内
获得改进~\cite[定理~1.5]{Scheder24}。
对一般 $3$-SAT，Jiang 与 Cai 的一篇预印本给出随机化
$O^*(1.307031578^n)$ 界~\cite[推论~1.2]{JC26}，其中 $O^*$
省略输入规模的多项式因子。Liu 的确定性
算法将分支与去随机化局部搜索相结合，在 $O(1.32793^n)$ 时间内
求解 $3$-SAT~\cite[定理~21]{Liu18}。

随机局部搜索也产生了有效的启发式方法。
WalkSAT 启发式选取一个未满足子句并翻转其一个变量，
将随机选择与贪心移动相结合~\cite{SKC94}。在来自电路综合与
诊断的公式上的实验展示了该方法的实用
价值~\cite[第~340--342 页，表~4--5]{SKC94}。

在均匀随机的 $k$-CNF 中，可满足性阈值处的实例
曾被猜测对判定可满足性是最难的~\cite[第~1 节]{VW21}；
对固定 $k\ge3$ 与较低的子句密度 $m/n\le 2^k/k^C$，
其中 $C$ 为通用常数，
已知在公式上以高概率存在多项式时间的几乎均匀的满足赋值
采样~\cite[定理~1.2]{CLWYY25}。
对含多项式多个子句的均匀随机公式，
在可满足的条件下，Vyas 与 Williams~\cite[定理~3]{VW21} 给出
基于 PPZ 的算法，它在公式与算法随机选择两方面均以高
概率在 $2^{n(1-\Omega(\log k/k))}$ 时间内找到解。

指数时间假设（ETH）断言 $3$-SAT 在最坏情况下需要
$2^{\Omega(n)}$ 确定性时间，因此对每个固定的 $k\ge3$，
$k$-SAT 也是如此~\cite[定理~1]{IP01}。
强指数时间假设（SETH）作出更尖锐的
预测：当 $k$ 增长时，$k$-SAT 需要 $2^{n(1-o_k(1))}$ 时间。
确切地说，对每个 $\varepsilon\in(0,1)$，存在固定的 $k$，
使得没有确定性算法能在
$2^{(1-\varepsilon)n}(n+m)^{O(1)}$ 时间内求解 $k$-SAT。
Sch\"oning 算法与 PPSZ 的运行时间界有助于
推动这一预测：它们相对穷举搜索的改进
随 $k$ 增长而减小~\cite[开放问题，第~2 项]{IP01}。

Gap 指数时间假设（Gap-ETH）通过假定常数 $\delta,\gamma>0$
来加强 ETH，使得没有 $O(2^{\gamma n})$ 时间
算法能区分一个可满足的 $3$-CNF 与一个其中
每个赋值至少违反 $\delta$ 比例子句的 $3$-CNF，
即使 $m=O(n)$~\cite{Din16,MR17}。其随机化版本排除了
具有相同运行时间、且在每个承诺输入上以至少 $2/3$ 的概率
区分这两种情形的算法~\cite[猜想~2.5]{CCK20}。
它已被用于证明 Label Cover~\cite{Din16} 与稠密约束满足
问题~\cite{MR17} 的近似困难性。它还排除了 Clique 与 Dominating Set 的
非平凡固定参数近似。对 Clique，被排除的
比值为 $o(\mathrm{OPT})$，其中 $\mathrm{OPT}$ 是最大
团的大小~\cite[推论~5.3 与~5.6]{CCK20}。

非确定性强指数时间假设（NSETH）
断言，对每个 $\varepsilon>0$，某个固定的 $k$ 不存在
非确定性 $2^{(1-\varepsilon)n}(n+m)^{O(1)}$ 时间算法
来判定一个 $k$-CNF 是否\emph{不可满足}，或等价地
判定一个 $k$-DNF 是否为重言式~\cite{CGIMPS16}。
NSETH 已被用于建立从 SAT 到 3SUM 与 All-Pairs Shortest
Paths 等问题的确定性细粒度归约的障碍~\cite{CGIMPS16}，
并用于刻画近似有向图直径的基于 SETH 的下界~\cite{Li21}。

定理~\ref{thm:main} 否证了 ETH、SETH、Gap-ETH 与 NSETH。
事实上，检验精确计数为正还是为零，对每个固定的 $k$ 都在确定性
$2^{o(n)}(n+m)^{O_k(1)}$ 时间内判定可满足性与不可满足性；
判定可满足性也解决了 Gap-ETH 的承诺问题。
该确定性算法还否证了 ETH~\cite[第~1 节]{DHMTW}、SETH~\cite[定义~2.1]{CDL16}
以及 Gap-ETH 的有界错误随机化版本。
它与正则消解~\cite{BeckImp13} 等受限证明系统的无条件指数下界
并不冲突。
2016 年，Ryan Williams 通过给出一个在 $2^{n/2}\poly(n,m)$ 时间内
认证 $k$-CNF 不可满足性的协议，否证了 SETH 的
Merlin--Arthur 版本（MA-SETH）~\cite[定理~3.3]{Wil16}。

\subsubsection{精确计数与计数复杂性}\label{subsec:related-counting}
在 Valiant 的框架中，$\#\mathrm{SAT}$ 与 $\#3$-SAT 在节俭
归约下是 $\#\mathrm P$ 完全的~\cite[第~3 节，事实~1--2]{Val79}。
对固定的 $k\ge3$，此前的精确计数算法达到
$\poly_k(n,m)2^{n-\Omega(n/k)}$ 时间，其中~\cite[引理~3.5]{IMP12} 为 Las Vegas
期望时间界，~\cite[定理~1.6]{CW21} 为确定性
界。对 $\#2$-SAT，Wahlstr\"om 的算法在
$O^*(1.2377^n)$ 时间内计数满足赋值~\cite{Wahlstrom08}。对 $\#3$-SAT，Kutzkov
使用互补分支在 $O^*(1.637^n)$ 时间内计数满足赋值~\cite{Kutzkov26}。两个算法都是确定性的。

计数指数时间假设（$\#\mathrm{ETH}$）
断言精确计数一个 $3$-CNF 的满足赋值
需要 $2^{\Omega(n)}$ 确定性时间~\cite[第~1 节]{DHMTW}。
在 $\#\mathrm{ETH}$ 下，Curticapean~\cite{Cur18} 排除了
即使在含 $v$ 个顶点与 $O(v)$ 条边的无权二部图中计数完美匹配的
确定性 $2^{o(v)}$ 时间算法。
推论~\ref{cor:graph-counting} 中的完美匹配算法
特别适用于这些图。
我们的算法（定理~\ref{thm:main}）在 $2^{o(n)}\poly(n,m)$ 时间内确定性地
求解 $\#3$-SAT，因此否证了 $\#\mathrm{ETH}$。

\subsubsection{对细粒度复杂性的影响}\label{subsec:related-consequences}
ETH 与 SETH 成为参数化与细粒度复杂性中的标准
假设~\cite{IPZ,VW19}。
基于 ETH 的归约给出了 Clique 中依赖解规模的下界~\cite{CHKX06}，
以及约束满足与划分式子图同构中依赖树宽的下界~\cite{Marx10}。
SETH 支持了 Hitting Set 与 Set
Splitting 的尖锐指数界~\cite{CDL16}，以及以树宽为参数的
Independent Set 与 Dominating Set 的界~\cite{LMS18}。
它还给出了 Subset Sum 与 Bicriteria Path 的伪多项式算法
可能加速幅度的界~\cite{ABHS19}。

对 $\mathrm P$ 中的问题，SETH 被用于排除
编辑距离~\cite{BackursIndyk18}、最长公共子序列与动态
时间规整~\cite{BK15}、Fr\'echet 距离~\cite{Bri14} 以及
正则表达式匹配~\cite{BI16} 的相对于二次运行时间的多项式因子加速。
其他应用涉及近似图
直径~\cite{RV13}，以及在图更新下维护可达计数或
最大强连通分量的大小~\cite{AW14}。
这些结果合在一起使 SETH 成为解释
相当不同问题上算法表面最优性的共同基础。

我们对 ETH 与 SETH 的否证移除了上述困难性主张的
一个基础。其中一些同样的下界可由可能仍然成立的
更弱假设推出。例如，最长公共
子序列与正则表达式匹配的下界可由正交
向量猜想（OVC）推出~\cite{BK15,BI16}。OVC 断言对每个
固定的 $\varepsilon>0$，没有 $O(N^{2-\varepsilon}\poly(d))$ 时间算法
能判定两组各含 $N$ 个 $d$ 维布尔向量、每组取一个向量的
集合是否包含一对正交向量。从 SAT 到正交向量的归约
表明 SETH 蕴含 OVC~\cite{Wil05,VW19}。
Fr\'echet 距离、图直径近似
以及 Independent Set 与 Dominating Set 的树宽依赖性的下界也
可由如下假设推出：无限制 CNF-SAT 对任何常数
$\delta>0$ 都不存在 $2^{(1-\delta)n}\poly(n,m)$ 时间算法~\cite{Bri14,RV13,LMS18}。
定理~\ref{thm:main} 对每个固定子句宽度 $k$ 给出
运行时间保证，而该假设允许宽度随 $n$ 增长。

要由我们的 SAT 算法得到目标问题的更快算法，
我们需要反方向的归约，即从该问题到
固定宽度 SAT 的归约，并控制变量数目。集合大小至多为 $k$ 的
Set Splitting 就有这样的
归约~\cite[观察~3.6]{CDL16}：为每个元素使用一个布尔变量
来指定它在划分中的一侧。对每个集合，一个子句
要求其变量中至少有一个为真，第二个子句要求
至少有一个为假。所得的 $k$-CNF 对每个
全集元素有一个变量，因此对固定 $k$，定理~\ref{thm:main} 给出
在全集规模上的亚指数时间，直至多项式输入规模
因子。

\subsubsection{组合问题的代数算法研究}\label{subsec:related-algebraic}
我们的算法是代数性的，基于对精心选取的辅助变量作多项式插值。下面我们讨论用代数方法设计组合问题算法的既有工作。更广泛的背景可参见 Koutis 与
Williams~\cite{KW16} 关于代数指纹的综述、Williams~\cite{Wil14}
关于算法设计中多项式方法的综述，以及
Husfeldt~\cite{Hus11} 关于算法容斥的综述。

Williams 的分裂列举算法~\cite{Wil05} 将
$2$-CSP 的变量划分为三组，枚举部分赋值，
并用快速矩阵乘法将它们组合。
Koivisto~\cite{Koi06} 给出直接的和积表述。
第~\ref{sec:counting}--\ref{sec:hashes} 节的算法
同样在组内枚举赋值。固定
共享的测量 $\vec y$ 使各组贡献独立相乘。
Scott 与 Sorkin~\cite[第~6 节]{SS09} 将
分裂列举推广到多项式值权重与任意环。
我们的求和例程同样使用截断多项式环
在一次计算中保留若干系数。

Bj\"orklund~\cite[第~2--3 节]{Bjo16} 使用容斥与
多项式插值，将积和式与有向哈密顿回路计数归约为对更小实例的加权求和。
他把相同实例的权重累加进一张表，然后对每个
总权重非零的不同实例只求解一次。
Bj\"orklund、Kaski 与 Williams~\cite{BKW19} 使用
低次曲线与广义 Kakeya 集，由在小的结构集上的求值恢复多项式
值，从而得到更快的 fermionant 算法，包括元素绝对值以固定常数为界的整数矩阵的积和式。
我们的多项式替换服务于类似目的：它们将
所需计数与可由其重构的结构化求值联系起来。这里的
有用结构是在固定少数共享测量后分裂成组。

Bhargava、Ghosh、Kumar 与
Mohapatra~\cite[第~3.1--3.2 节]{BGKM23} 沿低次曲线在有限域上求值多元
多项式，在仅有值不够时使用
Hasse 导数与 Hermite 插值。对一元限制 $g$，$t$ 处的 $j$ 阶
Hasse 导数是 $g(t+z)$ 中 $z^j$ 的系数。我们使用同样的恢复原理：
在每个允许快速求和的点处计算若干泰勒
系数，然后通过 Hermite 插值恢复
原始计数。

Woodruff 与 Yekhanin~\cite{WY07} 在私有信息检索中
沿一条直线使用值与一阶导数来恢复目标值。
Kopparty、Saraf 与 Yekhanin~\cite{KSY14} 使用高阶
Hasse 导数对多重码作局部纠正，推广了
Reed--Muller 码的基于直线的恢复。在我们的算法中，
矩阵构造提供低次替换形式，使我们能在固定少数共享测量后
通过快速求和计算插值点处的泰勒系数。
\subsubsection{矩阵构造的源流}\label{subsec:related-spectral}
我们的计数算法计算多项式乘积在
布尔赋值上的和。这些多项式共享变量，因此这些和
不分解为独立的更小的和。我们引入线性
替换，它们以多项式方式依赖于一个插值参数，
使得在固定少数整数
线性测量之后这种分解成为可能。一个辅助方阵有理矩阵 $M_0$
是这些替换的基础：$\operatorname{rank}(I-M_0)$ 的界
限制了测量的数量，而对 $M_0$ 的非零
元素的限制使替换次数保持较小。
第~\ref{subsec:preview-auxiliary} 节解释这一联系。

相关的插值方法曾被用于获得对结构化矩阵的
低秩逼近。在~\cite[第~3 节]{AW17} 中，
一个低次多项式在整数
内积 $j$ 的区间上与 $(-1)^j$ 相符，使用的是~\cite[引理~3.1]{AW15}
的插值引理。这给出在大多数元素上与
Walsh--Hadamard 矩阵一致的低秩矩阵，这是
矩阵刚性研究~\cite{Val77} 中的一个结果。相关的插值在~\cite[定理~6.1]{AGP23} 中
给出了 Kronecker 幂的刚性界。
Dvir 和 Liu~\cite[断言~2.20 与引理~3.3--3.4]{DL20} 使用谱方法：
对一个在其坐标的置换下不变的函数做稀疏修改会使许多傅里叶系数消失，
从而降低其群循环矩阵的秩。

我们的构造也在一个短区间上插值 $(-1)^j$，现在
是为了得到一个关于两两交换的置换矩阵的多项式。同一个
多项式既给出 $I-M_0$ 的秩界，也给出对 $M_0$
非零元素的限制。我们在第~\ref{sec:construction} 节中
直接证明这些性质及其定量界。

\subsubsection{稀疏化与奇偶检验}\label{subsec:related-reductions}
我们使用稀疏化引理~\cite{IPZ} 的保持计数版本~\cite[附录~A]{DHMTW}，
其定量界来自~\cite[第~6 节]{CIP} 和
\cite[引理~5.14 与等式~(5.4)]{CalabroThesis}。
这将满足赋值划分到同一组布尔变量上的稀疏公式之间。
附录~\ref{app:sparsify} 记录了不相交分支论证
以及当其参数随输入规模变化时所需的界。

初始的奇偶检验是 Tanner 码的约束
\cite[第~II 节]{Tanner81}，作用于记录一个赋值违反了哪些子句的
二进制向量。我们使用~\cite[第~IV 与~VI 节]{SS96} 中
基于边的扩展码构造：
每条图边携带一个比特，每个顶点要求其
关联比特属于一个固定的二进制线性码。局部
码的最小距离连同图扩展性，迫使
满足所有这些约束的每个非零向量具有
大支撑集。这一构造与距离论证是标准的。
引理~\ref{lem:sparse-checks} 提供了此处所需的参数选择，
包括检验的数量与支撑集大小的界，
并通过确定性穷举搜索构造图和局部码。该图是一个 Cayley 多重图的二分二重覆盖。
定义该多重图的移位在基群上几乎平衡每个非常数
奇偶性：这就是小偏差条件~\cite{NN93,AR94}，在
\eqref{eq:hash-walsh} 中明确陈述。将未使用的边比特固定为零，得到
恰好 $m$ 个坐标上的码。
额外的奇偶检验排除其余非零违反
向量。我们依据条件期望方法~\cite{Raghavan88} 所依赖的
平均原理来选择这些检验，使用精确计数来评估每个
候选。第~\ref{sec:hashes} 节给出此处所需的构造与
定量界。

定理~\ref{thm:main} 促使人们重新寻找那些已知困难性结果依赖于 ETH 或 SETH 的
问题的算法。第~\ref{subsec:related-consequences} 节中的
例子既提示了直接的推论，也提示了需要新算法思想的
问题。

\subsection{记号}\label{subsec:notation}
对整数 $s\ge0$，记 $[s]=\{1,\ldots,s\}$，其中
$[0]=\varnothing$。对条件 $E$，标量
$\mathds{1}_{\{E\}}$ 在 $E$ 成立时等于一，否则为零。
我们使用自然对数，当底数为二时写作 $\log_2$。
对复杂度界，$\poly(t_1,\ldots,t_s)$ 表示形如
$(1+t_1+\cdots+t_s)^{O(1)}$ 的界，其隐含常数为绝对常数。
如下标 $\poly_{k,c}(t_1,\ldots,t_s)$ 那样，下标允许这些常数
依赖于所指明的固定参数。整数以带符号二进制记法存储，因此一个整数 $z$ 使用 $O(\log(|z|+2))$ 比特。
有理数存储为一个整数分子和一个正
整数分母；所述的比特长度界适用于它们中的每一个。
空间以比特计量。算术运算次数与
比特运算次数是区分开的。

对有限集 $S$，$\mathbb Q^S$ 中的向量以 $S$ 为坐标
索引，且除非另有说明均为列向量。
对 $U\subseteq S$，指示向量 $\mathbf1_U\in\{0,1\}^S$
的分量为 $(\mathbf1_U)_v=\mathds{1}_{\{v\in U\}}$，其中 $v\in S$。
我们用 $\vec{0}$ 表示零向量。对 $v\in S$，行
$\delta_v\in\mathbb Q^{1\times S}$ 选出坐标 $v$，因此
$\delta_vx=x_v$。以 $U\subseteq S$ 为索引的向量通过
在 $S\setminus U$ 的每个坐标处置零而延拓到 $S$。

我们用 $\vec{X}=(X_1,\ldots,X_n)$ 表示形式变量向量，
并用箭头表示诸如
$\vec{P}_{f,v}(T)$ 的多项式系数向量。对系数行 $\vec{p}=(p_i)_{i\in S}$
和向量 $z=(z_i)_{i\in S}$，我们将其标量配对写作
$\langle\vec{p},z\rangle=\sum_{i\in S}p_i z_i$。
特别地，$\langle\vec{P}_{f,v}(T),\vec{X}\rangle$ 是关于变量 $X_i$ 的
标量线性形式。

对关于变量 $\vec{X}$ 和变量 $z$ 的多项式 $P(\vec{X},z)$，
$[z^j]P$ 表示 $z^j$ 的系数，视为关于 $\vec{X}$ 的多项式，
而 $\deg_{\vec{X}} P$ 表示在 $\vec{X}$ 中的总次数。
一行多项式的次数是任一分量的最大
次数。除非另有说明，多项式
系数是有理的，矩阵核在
$\mathbb Q$ 上取。

\section{技术概览}\label{sec:technical-overview}
我们首先对一个 $3$-CNF 的满足赋值计数，其中每个
变量至多出现在三个子句中。其子句
多项式的乘积次数至多 $3n$，因此我们可以使用插值
从变量的小组上的和恢复精确计数。
第~\ref{subsec:preview-general-sat} 节随后使用不相交
稀疏化和变量复制，对每个
固定 $k\ge2$ 在 $(n+m)^{O_k(1)}2^{o(n)}$ 时间内计算任意 $k$-CNF 的满足赋值。
取 $k=3$ 已经反驳 ETH。
第~\ref{subsec:preview-full-case} 节解释在
原变量上计数如何给出命题~\ref{prop:interpolation-bound} 的指数 $O_k(\rho(n))$。

\subsection{有界出现次数情形}\label{subsec:preview-setting}
设 $F$ 有 $n$ 个声明的布尔变量。每个子句包含
至多三个文字，且每个变量至多出现在
三个子句中，正出现与负出现合并计数。
我们的任务是精确计算 $F$ 的满足赋值数。

我们删除重言子句以及子句内重复的
文字。空子句立即给出计数零。
这些操作保持出现次数界。对每个剩余
子句 $e$，设 $I_e\subseteq[n]$ 为其变量集。
对每个不出现在任何子句中的变量，添加一个等于
一、作用域为 $\{v\}$ 的因子。这样的因子不改变计数，
并确保每个变量在构造中都有指定的位置。我们称这些常数一项为\emph{哑子句}。
设 $\mathcal F$ 索引所有子句，包括哑子句。则
\begin{equation}\label{eq:overview-incidences}
 |I_e|\le3,\qquad
 \#\{e\in\mathcal F:v\in I_e\}\le3,\qquad
 n\le M:=\sum_{e\in\mathcal F}|I_e|\le3n.
\end{equation}
下面所有的和仍遍历原来的 $n$ 个布尔变量。

\subsection{计数的直接多项式}\label{subsec:preview-encoding}
对正文字写 $L_{e,v}(X_v)=X_v$，对负文字
写 $L_{e,v}(X_v)=1-X_v$。子句 $e$ 的满足
多项式为
\[
 \Phi_e((X_v)_{v\in I_e})
 =1-\prod_{v\in I_e}(1-L_{e,v}(X_v)).
\]
例如，$x_a\vee\neg x_b\vee x_c$ 给出
\[
 \Phi_e(X_a,X_b,X_c)=1-(1-X_a)X_b(1-X_c).
\]
在布尔输入上，该多项式在子句被
满足时为一，否则为零。对哑子句 $e$，令
$\Phi_e=1$。因此
\begin{equation}\label{eq:overview-direct-count}
 \#F=\sum_{x\in\{0,1\}^n}\prod_{e\in\mathcal F}
                   \Phi_e((x_v)_{v\in I_e}).
\end{equation}
该乘积的总次数至多 $M\le3n$。
这一线性次数界是三次
出现假设的有用推论。

等式~\eqref{eq:overview-direct-count} 是计数的精确表达式，
但计算其 $2^n$ 个求和项将是穷举搜索。
为引出算法，假设给定变量到 $r$ 个组 $S_1,\ldots,S_r$ 的划分，
满足以下两个条件：
\begin{enumerate}
\item 每组至多包含 $2n/r$ 个变量。
\item 对每个子句，其所有变量都属于单一组。
\end{enumerate}
这是一个假设性的、非常严格的假定；我们并不
假设输入公式允许这样的划分。

在此假定下，我们可以独立地计算
每组的满足赋值。设 $N_g$ 为 $S_g$ 中变量上满足该组所有子句的
赋值数。任何这样的赋值都可以与
其他组中的任何满足赋值组合，因为没有子句连接不同组。
每个全局满足赋值恰好以这样一种方式出现。
因此，
\[
 \#F=\prod_{g=1}^r N_g.
\]
在其组内通过穷举搜索计算每个 $N_g$ 共检查
至多
\[
 \sum_{g=1}^r 2^{|S_g|}\le r\,2^{2n/r}
\]
个局部赋值，每个赋值的多项式工作量。对大的
$r$，这比枚举所有 $2^n$ 个全局赋值快得多。

我们将替换子句输入，使得不同组仅通过
少量共享线性摘要相互作用。一旦
这些摘要的值固定，我们就可以再次对
每组分别求和。

\subsection{求值点更容易的多项式族}
\label{subsec:preview-deformation}
我们引入参数 $T$。在 $T=0$ 时，替换必须
恢复原来的子句输入。在选定的非零点
$T=1,\ldots,Q$ 处，它们应让我们通过固定
共享线性摘要来分离各组。这里 $Q$ 是选定点的数量，
$R$ 将界定替换系数的次数。
我们稍后选择这两个参数。

用 $\vec{X}=(X_1,\ldots,X_n)$ 表示形式变量向量。
对每个出现 $v\in I_e$，使用多项式系数向量
$\vec{P}_{e,v}(T)=(p_{e,v,1}(T),\ldots,p_{e,v,n}(T))$ 将
子句 $e$ 中的 $X_v$ 替换为标量内积
\[
 \langle \vec{P}_{e,v}(T),\vec{X}\rangle
 =\sum_{i=1}^n p_{e,v,i}(T)X_i.
\]
我们称此表达式为变量 $v$ 在子句 $e$ 中出现的\emph{替换线性形式}。
它替换子句多项式 $\Phi_e$ 中的
输入 $X_v$。
我们要求
\[
 \langle \vec{P}_{e,v}(0),\vec{X}\rangle=X_v,
 \qquad \deg \vec{P}_{e,v}\le R.
\]
不同子句可以对同一变量使用不同的替换线性形式。
定义
\begin{equation}\label{eq:simple-counting-polynomial}
 H(T)=\sum_{x\in\{0,1\}^n}\prod_{e\in\mathcal F}
       \Phi_e\bigl((\langle \vec{P}_{e,v}(T), x\rangle)_{v\in I_e}\bigr)
       \in\mathbb Q[T].
\end{equation}
在零点处每个替换恢复原变量，而
每个子句多项式的次数至多增长 $R$ 倍。因此
\begin{equation}\label{eq:overview-value-degree}
 H(0)=\#F,\qquad \deg H\le MR\le3nR.
\end{equation}
我们将从 $t=1,\ldots,Q$ 处的泰勒系数恢复 $H(0)$。
次数界意味着每个点 $O(1+nR/Q)$ 个系数
就足够，如第~\ref{subsec:preview-interpolation} 节所述。
因此我们的目标是使 $R$ 相对于 $Q$ 较小：许多
点配合低次替换，在每个点需要更少的系数。

在非零 $t$ 处，替换线性形式的值 $\langle\vec{P}_{e,v}(t),x\rangle$
不必属于 $\{0,1\}$，即使
$x$ 是布尔的。因此变换后的子句多项式
不必充当满足指示函数。

\subsection{通过固定测量分离求和}
\label{subsec:preview-separation}
固定 $t\in[Q]$。给定系数向量
$\vec{P}_{e,v}(t)$，其分量现在为常数。我们解释
什么结构让我们快速计算 $H(t)$，以及如何使用它。
在本小节中，记 $S_g=S_{t,g}$ 和 $C_g=C_{t,g}$。

\begin{definition}[$(r,\tau,D)$-设计]
\label{def:overview-measurement-structure}
对给定的系数向量 $\vec{P}_{e,v}(t)$，在 $t$ 处的
\emph{$(r,\tau,D)$-设计}
由组 $S_g$、子句指派 $\sigma_e(t)$
和测量矩阵 $C_g$ 组成，满足以下四个
条件：
\begin{enumerate}
\item \emph{小组。} 将变量划分为不相交的
组 $[n]=S_1\sqcup\cdots\sqcup S_r$，且 $|S_g|\le2n/r$。
\item \emph{每个子句一个指派组。} 将每个子句
$e$ 指派给一个组 $\sigma_e(t)\in[r]$。条件~4 必须对子句 $e$ 的每个替换线性形式
用同一组成立。
\item \emph{少量有界测量。} 每组 $g$ 有一个
整数矩阵 $C_g$，其列以 $S_g$ 为索引，有 $c_g$ 行。
每行指定对 $S_g$ 中变量的一个线性测量。
总行数至多 $2\tau n$，且每行中
系数绝对值之和至多 $D$。
对赋值 $x$，写
\[
 x_g=x|_{S_g},\qquad y_g=C_gx_g,\qquad
 \vec{y}=(y_1,\ldots,y_r).
\]
\item \emph{仅通过测量依赖于其他组。}
假设子句 $e$ 被指派给组 $g=\sigma_e(t)$。
该子句的每个替换线性形式可以自由依赖于
$S_g$ 中的变量。它对任何其他组 $S_h$ 的依赖必须是
该组测量 $y_h$ 的线性组合。
确切地说，对每个 $v\in I_e$，必须存在系数向量
$\vec{\lambda}_{e,v,h}\in\mathbb Q^{c_h}$，$h\ne g$，
使得对每个 $x\in\mathbb Q^n$，
\begin{equation}\label{eq:outline-separation}
 \langle\vec{P}_{e,v}(t),x\rangle
 =\langle\vec{P}_{e,v}(t)|_{S_g},x_g\rangle
   +\sum_{h\ne g}\langle\vec{\lambda}_{e,v,h},y_h\rangle.
\end{equation}
等价地，对每个 $h\ne g$，$\vec{P}_{e,v}(t)|_{S_h}$ 位于
$C_h$ 的行空间中。
\end{enumerate}
\end{definition}
矩阵 $C_h$ 由所有子句共享。一个子句的指派
组不必包含其原始变量。然而，一旦 $\vec{y}$
固定，整个变换后的子句多项式就只依赖于其指派组中的变量。

\smallskip\noindent
\textbf{算法。}
给定 $t$ 处的 $(r,\tau,D)$-设计，
首先通过求解条件~4 所保证的线性系统
找出系数向量 $\vec{\lambda}_{e,v,h}$。
对固定的 $\vec{y}$，设 $\mathcal X_{t,g}(y_g)$ 为
$S_g$ 中测量值等于 $y_g$ 的变量赋值：
\[
 \mathcal X_{t,g}(y_g)
 =\{u\in\{0,1\}^{S_g}:C_gu=y_g\}.
\]
对这种赋值 $u$，如下计算组 $g$ 的贡献。对指派给 $g$ 的每个子句 $e$，首先使用
\eqref{eq:outline-separation} 从 $u$ 和固定测量值计算其
替换线性形式的值。
接着在这些值上计算子句多项式 $\Phi_e$。
最后，将指派给 $g$ 的所有子句得到的值
相乘。将此乘积
记为 $A_{t,g}(u;\vec{y})$。此计算不需要
其他组的赋值。

布尔赋值的每个测量都是
$[-D,D]$ 中的整数。初始化 $Z=0$ 并执行以下步骤：
\begin{enumerate}
\item 尝试每个测量元组
$\vec{y}\in\mathcal Y:=\prod_{g=1}^r\{-D,\ldots,D\}^{c_g}$。
对每个元组，执行步骤~2 和~3。
\item 对每组 $g$，尝试其变量的每个布尔赋值 $u$。仅当 $C_gu=y_g$ 时保留它。将相应的
乘积 $A_{t,g}(u;\vec{y})$ 相加得到组和
\[
 B_{t,g}(\vec{y})=
 \sum_{u\in\mathcal X_{t,g}(y_g)}A_{t,g}(u;\vec{y}).
\]
若无赋值满足 $C_gu=y_g$，置 $B_{t,g}(\vec{y})=0$。
\item 将组和相乘，并把结果
$\prod_{g=1}^r B_{t,g}(\vec{y})$ 加到 $Z$。
\end{enumerate}
处理完所有测量元组后，返回 $Z$。

为看出为何这返回 $H(t)$，固定 $\vec{y}$ 并从每个集合 $\mathcal X_{t,g}(y_g)$ 中选一个
赋值。由于各组包含
不相交的变量集，这些选择组合成
一个完整赋值。每个测量值为
$\vec{y}$ 的完整赋值恰好以这样一种方式出现。其子句
多项式之积是组贡献
$A_{t,g}$ 之积。因此，将组和相乘再对
$\vec{y}$ 求和给出
\begin{equation}\label{eq:outline-separated-sum}
 H(t)=\sum_{\vec{y}\in\mathcal Y}\prod_{g=1}^r
       \left(\sum_{u\in\mathcal X_{t,g}(y_g)}
                         A_{t,g}(u;\vec{y})\right).
\end{equation}
每个布尔赋值恰好有一个测量元组，因此
此和对每个赋值的贡献恰好计数一次。

测量元组至多 $(2D+1)^{2\tau n}$ 个。对每个
元组，算法至多尝试
$\sum_g2^{|S_g|}\le r2^{2n/r}$ 个对各个组的赋值。
因此，运行时间为
\[
 (2D+1)^{2\tau n}2^{2n/r}
\]
至多相差多项式因子。第一个因子界定测量
元组的数量；第二个界定一组内
赋值的数量。算法一次扫描一个组，
仅保留当前测量元组、当前赋值
以及运行中的和与积；见
引理~\ref{lem:measurement-separation}。

对第~\ref{subsec:preview-parameters} 节中的参数选择，
我们可以构造次数至多 $R$ 的替换形式，
它们在零点恢复原变量，并在每个 $t\in[Q]$ 处允许
$(r,\tau,D)$-设计。
第~\ref{subsec:preview-auxiliary} 节给出该构造。
我们首先解释这些性质如何让我们
高效地恢复计数。

\subsection{从泰勒系数恢复计数}
\label{subsec:preview-interpolation}
我们的目标是从点
$t=1,\ldots,Q$ 处的信息恢复 $H(0)=\#F$。由于 $\deg H\le3nR$，仅 $Q$ 个值 $H(t)$
可能无法确定 $H$。我们改为在每个点计算若干泰勒
系数。选择
\[
 q=\left\lfloor\frac{3nR}{Q}\right\rfloor,\qquad J=q+1.
\]
假设我们知道对每个 $t\in[Q]$，$H(t+z)$ 中 $z^0,\ldots,z^q$ 的
系数。这些数据唯一确定 $H$：若两个
次数至多 $3nR$ 的多项式有相同数据，其
差对每个 $t\in[Q]$ 可被 $(T-t)^J$ 整除，从而可被
$\prod_{t=1}^Q(T-t)^J$ 整除。该乘积的次数为 $QJ>3nR$，因此
差必须为零。每个泰勒系数给出关于 $H$ 的系数的一个线性
方程。求解此有理线性
系统恢复 $H$，其常数项即为所需的
计数 $H(0)$。

剩下的是计算这些泰勒系数。固定 $t\in[Q]$；
我们想要 $H(t+z)$ 关于 $z$ 直到次数 $q$ 的部分。设计控制在 $t$ 处的替换形式，
但不必控制它们从 $t$ 到 $t+z$ 的变化。
我们首先将这些变化作为单独的输入，
使得固定此输入保持各组的分离。

\smallskip\noindent
\textbf{保持移位后的输入分离。}
写
\[
 \Delta\vec P_{e,v}(z)=\vec P_{e,v}(t+z)-\vec P_{e,v}(t).
\]
保留 $x$ 作为被求和的布尔赋值，并引入
单独的形式向量 $\vec X$ 表示变化。定义
\begin{equation}\label{eq:outline-taylor-product}
 \mathcal G_x(\vec X,z)=
 \prod_{e\in\mathcal F}\Phi_e\bigl(
       (\langle\vec P_{e,v}(t),x\rangle+
        \langle\Delta\vec P_{e,v}(z),\vec X\rangle)_{v\in I_e}\bigr).
\end{equation}
这里 $x$ 索引该多项式，而 $\vec X$ 和 $z$ 是其
形式变量。令 $\vec X=x$ 恢复每个移位后的输入
$\langle\vec P_{e,v}(t+z),x\rangle$，因此
\[
 H(t+z)=\sum_{x\in\{0,1\}^n}\mathcal G_x(x,z).
\]
另一方面，对某个
$U\subseteq[n]$ 固定 $\vec X=\mathbf1_U$，使每个修正
$\langle\Delta\vec P_{e,v}(z),\mathbf1_U\rangle$ 与 $x$ 无关。
所有对 $x$ 的依赖于是仍通过
$\langle\vec P_{e,v}(t),x\rangle$ 保留，而设计在此适用。
因此我们的任务是从 $\vec X$ 固定的和
恢复 $\vec X=x$ 的和。

次数界使得这只需少量固定替换即可实现。
$\Delta\vec P_{e,v}(z)$ 的每个条目都能被 $z$ 整除，因此展开项中每出现一次 $X_i$ 至少贡献一个 $z$ 因子。于是
\[
 \deg_{\vec X}([z^j]\mathcal G_x(\vec X,z))\le j.
\]
因此，在丢弃 $q$ 次以上的 $z$ 的幂之后，每个单项式至多涉及 $\vec X$ 中的 $q$ 个不同变量。

\smallskip\noindent
\textbf{用容斥恢复平移后的和。}
在模 $z^J$ 下工作，从而只保留到 $z^q$ 的系数。将 $\mathcal G_x$ 中 $X_i$ 的每个正幂替换为 $X_i$。这在布尔输入上保持其值不变。上述次数界给出，对布尔 $\vec X$，
\[
 \mathcal G_x(\vec X,z)\equiv
 \sum_{\substack{E\subseteq[n]\\|E|\le q}}
       a_E(x,z)\prod_{i\in E}X_i\pmod{z^J},
\]
其中 $a_E(x,z)$ 是此替换后所指单项式的系数。在 $\mathbf1_U$（恰在 $U$ 中为 1 的向量）处求值，会对以 $U$ 的子集为下标的系数求和。容斥提取出每一个：
\[
 a_E(x,z)\equiv
 \sum_{U\subseteq E}(-1)^{|E|-|U|}
                 \mathcal G_x(\mathbf1_U,z)\pmod{z^J}.
\]
当令 $\vec X=x$ 时，单项式 $\prod_{i\in E}x_i$ 恰好当 $E$ 中所有位均为 1 时等于 1。因此它对关于 $x$ 的求和的贡献为 $\sum_{x:\,x_i=1\ (i\in E)}a_E(x,z)$，其中 $x$ 为布尔值。代入上面的系数公式后，这归结为形式输入固定时的求和。

\smallskip\noindent
\textbf{计算受限和。}
对满足 $|E|\le q$ 的 $U\subseteq E\subseteq[n]$，定义
\begin{equation}\label{eq:outline-restricted-taylor-sum}
 \mathcal Z_{t,E,U}(z)=
 \sum_{\substack{x\in\{0,1\}^n\\x_i=1\ (i\in E)}}
                    \mathcal G_x(\mathbf1_U,z)
       \quad\text{in }\mathbb Q[z]/(z^J).
\end{equation}
这里 $U$ 固定形式输入 $\vec X=\mathbf1_U$，而 $E$ 将求和变量 $x$ 的某些位固定为 1。对单项式贡献求和得到
\begin{equation}\label{eq:outline-taylor-recovery}
 H(t+z)\equiv
 \sum_{\substack{E\subseteq[n]\\|E|\le q}}
       \sum_{U\subseteq E}(-1)^{|E|-|U|}
                         \mathcal Z_{t,E,U}(z)\pmod{z^J}.
\end{equation}
这同时恢复全部 $J=q+1$ 个泰勒系数。引理~\ref{lem:boolean-differences} 陈述了此处使用的一般恒等式。

每个和 $\mathcal Z_{t,E,U}$ 都可用第~\ref{subsec:preview-separation} 节的算法计算。在 $\vec X=\mathbf1_U$ 固定时，每个修正 $\langle\Delta\vec P_{e,v}(z),\mathbf1_U\rangle$ 与 $x$ 无关。对 $x$ 的全部依赖都通过形式 $\langle\vec P_{e,v}(t),x\rangle$ 出现，因此固定 $\vec y$ 仍可分离子句多项式。限制 $x_i=1$ 对每个组是局部的。分离恒等式只使用加法、乘法和分配律，因此也适用于这些截断的多项式值。每个子句在代入 $\vec X=\mathbf1_U$ 后求值，只使用次数小于 $J$ 的关于 $z$ 的一元多项式，并在每次乘法后丢弃更高次幂。多元展开及其系数 $a_E(x,z)$ 仅用于说明该恒等式的依据；算法从不构造它们。

在 $t=1,\ldots,Q$ 处重复此计算，可提供上述插值所需的全部泰勒系数，从而恢复 $H(0)=\#F$。

\subsection{三项开销与运行时间}
\label{subsec:preview-parameters}
提前说明，第~\ref{subsec:preview-auxiliary} 节将构造替换形式，它们在零点恢复原始变量、次数至多为 $R$，并且在每个 $t\in[Q]$ 处容许一个 $(r,\tau,D)$-设计。对充分大的 $n$，该构造在多项式时间内完成，并达到以下参数界：
\begin{equation}\label{eq:overview-parameter-bounds}
 \begin{aligned}
 r&=\Theta\!\left(\frac{\sqrt{\log n}}{(\log\log n)^2}\right),
 &Q&=\Theta\!\left(\frac{\log n}{\log\log n}\right),\\
 R&=\Theta(\sqrt{\log n}),
 &\tau&=\Theta\!\left(\frac1{r\log n}\right),\qquad
 D=r^{2Q}\le(2n)^{1/128}.
 \end{aligned}
\end{equation}
我们首先说明这些界如何给出快速计数算法；下一小节将论证它们可以同时达到。

回顾 $q=J-1=\lfloor3nR/Q\rfloor$。对每个求值点 $t$，第~\ref{subsec:preview-interpolation} 节对每个满足 $|E|\le q$ 的对 $U\subseteq E$ 计算一个受限和。大小为 $j$ 的 $E$ 有 $\binom nj$ 种选择，随后 $U$ 有 $2^j$ 种选择。每个受限和枚举测量元组以及各个组内的赋值。因此三项开销相乘，至多相差多项式因子，得到
\begin{equation}\label{eq:overview-three-costs}
 \underbrace{\sum_{j=0}^{\min(q,n)}2^j\binom nj}_{\text{Taylor data}}
 \ \cdot\
 \underbrace{(2D+1)^{2\tau n}}_{\text{measurements}}
 \ \cdot\
 \underbrace{2^{2n/r}}_{\text{one group}}.
\end{equation}
为估计第一项开销，注意对充分大的 $n$ 有 $1\le q\le n/2$，因为 $R/Q\to0$。利用 $\binom nj\le(en/j)^j$ 以及 $(2en/j)^j$ 在 $1\le j\le n$ 上递增这一事实，得到
\[
 \sum_{j=0}^q2^j\binom nj
 \le(q+1)\left(\frac{2en}{q}\right)^q.
\]
由于 $q=\Theta(nR/Q)$，此开销的对数为 $O(n(R/Q)\log(2Q/R))$。因此对 \eqref{eq:overview-three-costs} 中三项开销取对数，得到
\[
 O\!\left(n\left[
       \frac RQ\log\frac{2Q}{R}
       +\tau\log(2D+1)+\frac1r\right]\right).
\]
在 \eqref{eq:overview-parameter-bounds} 下每一项都很小：
\[
 \begin{aligned}
 \frac RQ\log\frac{2Q}{R}
   &=O\!\left(\frac{\log\log n}{\sqrt{\log n}}
                      \cdot\log\log n\right),\\
 \tau\log(2D+1)&=O\!\left(\frac1{r\log n}\cdot\log n\right)
                  =O(1/r),\\
 \frac1r&=O\!\left(\frac{(\log\log n)^2}{\sqrt{\log n}}\right).
 \end{aligned}
\]
因此，有界出现次数的实例可在
\begin{equation}\label{eq:overview-bounded-time}
 n^{O(1)}2^{O(n(\log\log n)^2/\sqrt{\log n})}
\end{equation}
位运算和多项式空间内计数。有限多个更小的输入由穷举搜索处理。

点数 $Q$、组数 $r$ 和截断阶 $J=O(1+nR/Q)$ 都是 $n$ 的多项式。因此，对点和组的循环、次数小于 $J$ 的多项式上的算术以及最终插值，在上述三个枚举的每次迭代中只增加多项式多次运算。

关于位复杂度和空间界，见附录~\ref{subsec:overview-arithmetic}。

\subsection{构造低次替换线性形式}
\label{subsec:preview-auxiliary}
考虑某个子句中 $X_v$ 的一次出现。在 $T=0$ 时，其替换必须等于 $X_v$。在每个 $t=1,\ldots,Q$ 处，该子句必须能够仅使用其被分配的组和少量共享测量来计算同一替换。在每个点选择不同的线性形式并进行插值，通常要付出次数 $Q$ 的代价。为使泰勒恢复的开销低，我们需要次数 $R$ 远小于 $Q$。

假设替换中的某一项仅在 $t=2$ 处使用子句被分配组之外的变量。将该项乘以 $1-T/2$ 使它在零点保持不变，并在 $t=2$ 处将其消去。在每个其他求值点，子句都可以从自己的组中读取该变量。同样的技巧对若干这样的点都有效：每个点一个因子，每处代价为一次。因此，挑战在于把每个输入写成这样的和：共享测量确定其中一部分，而每个剩余项仅在少数几个点使用被分配组之外的变量。我们首先在记录组选择的辅助坐标上构造这样的分解，然后代入公式的变量。

\medskip\noindent
\textbf{第 1 步：建立辅助坐标。}\par\nobreak
\smallskip\noindent
在每个求值点，一次出现 $(e,v)$ 有两种组选择：包含变量 $v$ 的组，以及分配给子句 $e$ 的组。在 $Q$ 个点上，这些选择构成 $Q$ 个对的序列。因此我们使用坐标集
\[
 \mathcal B=([r]\times[r])^Q,\qquad D=|\mathcal B|=r^{2Q}.
\]
将一个下标写成
\[
 u=\bigl((a_u(1),d_u(1)),\ldots,(a_u(Q),d_u(Q))\bigr).
\]
我们将为每个这样的下标构造一个替换行，因此无论之后为公式作出怎样的组选择，它都可用。

令 $\vec\zeta=(\zeta_w)_{w\in\mathcal B}$ 为原始变量的形式占位符，我们将在第~4 步中代入它们。在点 $t$，坐标 $\zeta_w$ 属于组 $a_w(t)$，而 $\zeta_u$ 的替换被分配到组 $d_u(t)$。因此，该替换可以自由使用的坐标为
\[
 \mathcal S_{t,d_u(t)},\qquad
 \mathcal S_{t,g}:=\{w\in\mathcal B:a_w(t)=g\}.
\]
特别地，$\zeta_w$ 属于被分配组当且仅当 $a_w(t)=d_u(t)$。集合 $\mathcal S_{t,g}$ 划分 $\mathcal B$，每组有 $D/r$ 个坐标。

我们的辅助目标是构造系数行 $\vec P_u(T)\in\mathbb Q[T]^D$ 和一个矩阵
\[
 C\in\{0,1,-1\}^{c\times D},\qquad c\le\frac{D\tau}{3r},
\]
使得
\[
 \langle\vec P_u(0),\vec\zeta\rangle=\zeta_u,
 \qquad \deg\vec P_u\le R-1,
\]
并且在每个 $t\in[Q]$ 处，替换线性形式 $\langle\vec P_u(t),\vec\zeta\rangle$ 是可用数据
\begin{equation}\label{eq:overview-auxiliary-requirement}
 \vec\zeta|_{\mathcal S_{t,d_u(t)}}\quad\text{and}\quad C\vec\zeta.
\end{equation}
的线性组合。相同的测量 $C\vec\zeta$ 服务于每一行和每个求值点。这是第~\ref{subsec:preview-separation} 节设计中依赖条件的辅助对应物。

\medskip\noindent
\textbf{第 2 步：构造常数矩阵。}\par\nobreak
\smallskip\noindent
我们从恒等式出发
\[
 \zeta_u=((I_D-M_0)\vec\zeta)_u+
                   \sum_{w\in\mathcal B}M_0[u,w]\zeta_w.
\]
我们希望第一项由 $C\vec\zeta$ 确定。在和式中，每个非零系数应只在少数几个点连接不相容的组，从而低次多项式可以消除这些依赖。用下式定义这些点
\[
 E_{u,w}=\{t\in[Q]:a_w(t)\ne d_u(t)\}.
\]
对 $M_0\in\mathbb Q^{D\times D}$ 的两个要求是
\begin{equation}\label{eq:outline-matrix-conditions}
 \begin{gathered}
 M_0\vec\xi=\vec\xi\quad\text{for every }\vec\xi\in\ker C,\\
 M_0[u,w]\ne0\quad\Longrightarrow\quad |E_{u,w}|\le R-1.
 \end{gathered}
\end{equation}
第一条说明 $I_D-M_0$ 在 $\ker C$ 上为零。等价地，它的每一行都位于 $C$ 的行空间中，因此 $(I_D-M_0)\vec\zeta$ 是测量 $C\vec\zeta$ 的线性函数。第二条将界定第~3 步所用多项式的次数。

\smallskip\noindent
\emph{通过保持 $f$ 低次来限制支撑集。}
令 $\Pi_s$ 为置换矩阵，它在每个下标中交换第 $s$ 对的两个条目，保持所有其他对不变。这些矩阵可交换并满足 $\Pi_s^2=I_D$。令
\[
 V_0=\prod_{s=1}^Q\Pi_s,\qquad
 W_0=\frac12\sum_{s=1}^Q(I_D-\Pi_s),\qquad
 M_0=V_0f(W_0),
\]
其中我们将选择次数至多为 $R-1$ 的多项式 $f$。完全交换 $V_0$ 已经尊重每个组：其第 $u$ 行的非零元位于满足
\[
 (a_w(s),d_w(s))=(d_u(s),a_u(s))\quad\text{for every }s.
\]
的列 $w$。特别地，$E_{u,w}$ 为空。

乘以 $f(W_0)$ 可以撤销交换，但 $f$ 的次数限制了可撤销的数量。展开 $f(W_0)$ 得到至多 $R-1$ 个交换的乘积。重复的交换会抵消，因为 $\Pi_s^2=I_D$。因此乘以 $V_0$ 得到
\begin{equation}\label{eq:overview-partial-swaps}
 M_0=\sum_{\substack{S\subseteq[Q]\\|S|\le R-1}}
             \gamma_S\prod_{s\notin S}\Pi_s
 \qquad\text{for some }\gamma_S\in\mathbb Q.
\end{equation}
以下标 $S$ 索引的求和项中的置换交换 $S$ 之外的每一对。因此其非零元满足：对每个 $s\notin S$ 有 $a_w(s)=d_u(s)$，所以 $E_{u,w}\subseteq S$。每当 $|E_{u,w}|>R-1$ 时每个求和项为零，因此它们的和也为零。这证明了对任意次数至多为 $R-1$ 的多项式 $f$，\eqref{eq:outline-matrix-conditions} 中的第二个条件成立。

\smallskip\noindent
\emph{选择 $f$ 以固定大多数基方向。}
我们现在选择 $f$，使 $M_0$ 固定大多数基方向。首先考虑一个位置，其坐标以 $[r]\times[r]$ 为下标。令 $\vec e_{ab}$ 为在 $(a,b)$ 处为 1、其他处为 0 的向量。交换该对具有正交特征基
\[
 \begin{array}{ll}
 \vec e_{aa}\quad(a\in[r]),\quad
 \vec e_{ab}+\vec e_{ba}\quad(a<b),
     &\text{eigenvalue }+1,\\[2pt]
 \vec e_{ab}-\vec e_{ba}\quad(a<b),
     &\text{eigenvalue }-1.
 \end{array}
\]
有 $r(r+1)/2$ 个特征值为 $+1$ 的向量和 $r(r-1)/2$ 个特征值为 $-1$ 的向量。所有条目都在 $\{0,1,-1\}$ 中。

对每个位置 $s\in[Q]$，选择这样一个向量 $\vec\psi_s$，并逐坐标定义 $\vec\psi\in\mathbb Q^{\mathcal B}$：
\[
 \psi_u=\prod_{s=1}^Q\psi_s(a_u(s),d_u(s)).
\]
这 $D=(r^2)^Q$ 种选择给出一个正交基：内积分解为 $Q$ 个位置处内积的乘积。交换 $\Pi_s$ 通过 $\vec\psi_s$ 的特征值作用于 $\vec\psi$。若恰有 $j$ 个选择的特征值为 $-1$，则
\[
 W_0\vec\psi=j\vec\psi,\qquad
 V_0\vec\psi=(-1)^j\vec\psi,\qquad
 M_0\vec\psi=(-1)^jf(j)\vec\psi.
\]
因此我们通过令 $f(j)=(-1)^j$ 使 $M_0$ 固定这些向量。

对均匀选择的基向量，数 $j$ 服从分布
\[
 W\sim\operatorname{Bin}(Q,p),\qquad p=\frac{r-1}{2r},
\]
因为在每个位置，$r^2$ 种选择中有比例 $p$ 的特征值为 $-1$。因此大多数基向量的 $j$ 接近 $Qp$。在 $Qp$ 附近选择由 $R$ 个连续整数构成的区间 $\mathcal I$，并对 $j\in\mathcal I$ 插值 $f(j)=(-1)^j$。这给出 $\deg f\le R-1$ 并使 $M_0$ 固定每个满足 $j\in\mathcal I$ 的基向量。第~5 步指定端点并检查该区间位于 $\{0,\ldots,Q\}$ 中。

\smallskip\noindent
\emph{测量例外方向。}
将其余基向量的转置用作 $C$ 的行。它们的条目在 $\{0,1,-1\}$ 中，且由正交性，$\ker C$ 由上面固定的基向量张成。这证明了 \eqref{eq:outline-matrix-conditions} 中的第一个条件。

现在可以看出如何选择 $R$。$W$ 的标准差为 $\sqrt{Qp(1-p)}=\Theta(\sqrt Q)$。取 $R$ 远大于 $\sqrt Q$ 会使大多数基向量落入区间内，因此只需测量一小部分。我们仍可取 $R$ 远小于 $Q$，使得泰勒恢复每个点只需 $1+\lfloor 3nR/Q\rfloor=o(n)$ 个系数。对第~5 步中的选择，二项式尾界给出
\[
 c=D\Pr[W\notin\mathcal I]
      \le 2D\exp\!\left(-\frac{R^2}{8Q}\right)
      \le\frac{D\tau}{3r},
\]
符合要求。因此，低次给出受限支撑集，而插值使修正 $I_D-M_0$ 只依赖于这少量测量。

\medskip\noindent
\textbf{第 3 步：定义辅助替换行。}\par\nobreak
\smallskip\noindent
我们现在把第~2 步中 $\zeta_u$ 的分解转化为其替换。保持由测量确定的项不变。对每个 $M_0[u,w]\ne0$，将其在另一项中的系数乘以
\[
 \varphi_{u,w}(T)=\prod_{s\in E_{u,w}}(1-T/s).
\]
该多项式在零点等于 1，并在 $\zeta_w$ 位于被分配组之外的每个点处消失。通过恒等式定义系数行 $\vec P_u(T)$
\begin{equation}\label{eq:outline-auxiliary-row}
 \langle\vec P_u(T),\vec\zeta\rangle
 =((I_D-M_0)\vec\zeta)_u+
   \sum_{\substack{w\in\mathcal B\\M_0[u,w]\ne0}}
                M_0[u,w]\varphi_{u,w}(T)\zeta_w.
\end{equation}
这直接给出全部三个所需性质：
\begin{itemize}
\item 在 $T=0$ 时，每个乘数均为 1，因此两项之和为 $\zeta_u$。
\item 每个乘数的次数为 $|E_{u,w}|\le R-1$，因此 $\deg\vec P_u\le R-1$。
\item 在 $T=t$ 时，第一项是 $C\vec\zeta$ 的线性函数。在和式中，系数只有在 $a_w(t)=d_u(t)$ 时才能保留，因此只剩下 $\mathcal S_{t,d_u(t)}$ 中的坐标。
\end{itemize}
因此，该替换恰好使用 \eqref{eq:overview-auxiliary-requirement} 中允许的那类数据。

\medskip\noindent
\textbf{第 4 步：代入原始变量。}\par\nobreak
\smallskip\noindent
辅助行现已构造完成。为将其用于公式，选择调度
\[
 a_v\in[r]^Q\quad(v\in[n]),\qquad
 d_e\in[r]^Q\quad(e\in\mathcal F).
\]
在点 $t$，将变量 $v$ 放入组 $a_v(t)$，并将子句 $e$ 分配到组 $d_e(t)$：
\[
 S_{t,g}=\{v:a_v(t)=g\},\qquad \sigma_e(t)=d_e(t).
\]
出现 $(e,v)$ 使用辅助标签
\[
 u(e,v)=\bigl((a_v(s),d_e(s))\bigr)_{s\in[Q]}\in\mathcal B.
\]
令 $h$ 为任一标签所接收的出现次数的最大值。我们需要变量组小且 $h$ 也小。对于我们的三次出现、宽度为三的公式，引理~\ref{lem:incidence-labels} 在多项式时间内确定性地选择调度，使得
\[
 |S_{t,g}|\le2n/r,\qquad h\le2M/D,
\]
只要
\begin{equation}\label{eq:overview-balance-requirements}
 12D^2\le M,\qquad 2Qr^2\le n.
\end{equation}
其思路是对调度的随机选择去随机化。若所有 $a_v$ 和 $d_e$ 都从 $[r]^Q$ 中独立均匀地选取，则每个标签的期望重数为 $M/D$，每个组 $S_{t,g}$ 的期望大小为 $n/r$。在 \eqref{eq:overview-balance-requirements} 的假设下，条件期望的标准方法在多项式时间内对该选择去随机化，确保每个标签的重数至多为 $2M/D$，每个组的大小至多为 $2n/r$。第~5 步对我们的参数选择验证这些假设。
不同的出现可能获得相同标签，却在该坐标处需要不同的原始变量。使用 $\mathcal B$ 的 $h$ 个副本，把具有相同标签的出现放在不同的副本中。令 $j(e,v)\in[h]$ 表示 $(e,v)$ 所用的副本；把 $X_v$ 放在其坐标 $u(e,v)$ 处，并在未使用的位置置零。用 $\vec\zeta^{(j)}(\vec X)$ 表示副本 $j$ 中得到的填充向量。每个条目都是零或一个原始变量。$X_v$ 的重复出现使用相同的值，因此这些副本不引入新的布尔变量。

对包含该出现的副本应用辅助行：
\begin{equation}\label{eq:overview-transferred-row}
 \langle\vec P_{e,v}(T),\vec X\rangle
 =\left\langle\vec P_{u(e,v)}(T),
                  \vec\zeta^{(j(e,v))}(\vec X)\right\rangle.
\end{equation}
在零点处，右侧读取包含 $X_v$ 的坐标，且代入变量保持次数上界 $R-1$。在点 $t$ 处，其自由使用的坐标只包含满足 $a_{v'}(t)=d_e(t)$ 的变量 $v'$。这些正是该子句被指派的组 $S_{t,d_e(t)}$ 中的变量。子句 $e$ 的所有出现都使用这同一个被指派的组。

剩余的依赖通过 $h$ 个副本上的 $hc$ 个线性测量 $C\vec\zeta^{(j)}(\vec X)$ 体现。这些测量可能涉及若干原始组。固定诸如 $x_1+x_2$ 这样的总和仍会把包含 $x_1$ 与 $x_2$ 的两个组耦合起来。因此，我们把每个测量拆分为来自各个组的贡献。对每个测量 $\ell(\vec X)$，写
\[
 \ell(\vec X)=\sum_{g=1}^r
        \ell|_{S_{t,g}}(\vec X|_{S_{t,g}}),
\]
并把各限制的系数行收集到 $C_{t,g}$ 中。现在每个辅助测量都是组测量之和。对于指派到组 $g$ 的子句，来自 $S_{t,g}$ 的部分是局部线性形式，而其他每一部分都是相应组测量的线性组合。这就是 \eqref{eq:outline-separation} 中的条件。

总共有 $rhc$ 行。由于 $hD\le2M\le6n$，
\[
 \sum_{g=1}^r c_{t,g}
 =rhc\le rh\frac{D\tau}{3r}\le2\tau n.
\]
$C$ 的每一行在 $\{0,1,-1\}$ 中有 $D$ 个条目。代入变量、合并重复副本以及限制到一个组，都不能增加绝对系数之和。因此 $C_{t,g}$ 的每一行都是整值的，绝对系数和至多为 $D$。结合组大小界，这证明了替换形式在每一点都容许一个 $(r,\tau,D)$-设计。

\medskip\noindent
\textbf{第 5 步：验证所承诺的参数界。}\par\nobreak
\smallskip\noindent
该构造需要两个平衡不等式 \eqref{eq:overview-balance-requirements} 以及尾界 $2e^{-R^2/(8Q)}\le\tau/(3r)$。我们现在按 \eqref{eq:overview-parameter-bounds} 中所述的阶选择参数，并验证这些要求。

令 $\lambda=\lceil\log_2 n\rceil$，$\nu=\lceil\log_2\lambda\rceil$，并选择
\begin{equation}\label{eq:overview-explicit-parameters}
 \begin{aligned}
 r&=\left\lfloor\frac{\sqrt\lambda}{\nu^2}\right\rfloor,
 &Q&=\left\lfloor\frac{\lambda}{128\nu}\right\rfloor,
 &R&=4\lceil\sqrt\lambda\rceil,\\
 \tau&=\frac1{r\lambda},
 &D&=r^{2Q}.
 \end{aligned}
\end{equation}
它们具有 \eqref{eq:overview-parameter-bounds} 中所述的阶。对充分大的 $n$，有 $r\ge2$、$R\ge32$ 以及 $6R\le Q$，因为 $R/Q=O(\nu/\sqrt\lambda)\to0$。我们依次检验其余要求：
\begin{itemize}
\item \emph{插值区间。}
取 $p=(r-1)/(2r)$，如第~2 步中那样，并取
\[
 A_0=\left\lfloor Qp-\frac{R-1}{2}\right\rfloor,\qquad
 \mathcal I=\{A_0,\ldots,A_0+R-1\}.
\]
由于 $Q/4\le Qp\le Q/2$ 且 $6R\le Q$，该区间落在 $\{0,\ldots,Q\}$ 中。此外，$R\ge32$ 确保它包含与 $Qp$ 相距 $R/4$ 以内的每个整数。因此二项尾界给出第~2 步中使用的估计 $\Pr[W\notin\mathcal I]\le2e^{-R^2/(8Q)}$。
\item \emph{多项式维数与平衡调度。}
由于 $\log_2 r\le\nu/2$，
\[
 \log_2 D=2Q\log_2 r\le\lambda/128,\qquad
 D\le(2n)^{1/128}.
\]
对充分大的 $n$，这给出
\[
 12D^2\le12(2n)^{1/64}\le n\le M,\qquad
 2Qr^2=O(\lambda^2)\le n,
\]
正如第~4 步所要求。
\item \emph{测量数量少。}
我们有 $R^2/(8Q)\ge256\nu$。由于 $\lambda\le2^\nu$ 且 $r^2\le\lambda$，
\[
 2e^{-R^2/(8Q)}\le2e^{-256\nu}
 \le\frac1{3\lambda^2}
 \le\frac1{3r^2\lambda}=\frac{\tau}{3r}.
\]
因此第~2 步提供了第~4 步中使用的辅助测量界。
\end{itemize}

最后，所有构造维数都由 $n$ 的多项式界定。引理~\ref{lem:row-coefficients} 给出第~2--3 步中矩阵与行的多项式构造时间，引理~\ref{lem:incidence-labels} 对第~4 步中的调度给出同样的结论。该转移只使用 $hD\le6n$ 个坐标位置。因此我们在多项式时间内构造了替换形式，它们在零点处恢复每个原始输入，次数至多为 $R-1\le R$，并且在每个 $t\in[Q]$ 处容许所承诺的 $(r,\tau,D)$-设计。这提供了第~\ref{subsec:preview-parameters} 节中使用的界；引理~\ref{lem:incidence-parameters} 记录了完整证明中更一般情形的参数检验。

\subsection{针对一般 \texorpdfstring{$\#k$-SAT}{\#k-SAT} 的较弱算法}
\label{subsec:preview-general-sat}
上述算法已经给出了一般 $\#k$-SAT 的次指数算法：在不相交稀疏化之后，我们把每个实例直接归约为一个 $3$-CNF，其中每个变量至多出现三次。两步都保持计数，如下所述。

固定 $k\ge2$。我们使用稀疏化~\cite{IPZ} 的保计数版本~\cite[附录~A]{DHMTW}，其定量界来自~\cite[第~6 节]{CIP} 和 \cite[引理~5.14 与方程~(5.4)]{CalabroThesis}。如引理~\ref{lem:sparsify} 所记录，这把输入计数表示为至多 $\operatorname{poly}_k(n,\sigma^{-1})2^{\sigma n}$ 个公式 $G$ 的计数之和。每个 $G$ 有相同的 $n$ 个声明变量，每个变量的出现次数至多
\[
 B=O_k\!\left((\sigma^{-1}\log(2/\sigma))^{k-1}\right)
\]
因而非空子句至多 $Bn$ 个。空子句立即给出计数为零。

首先把 $G$ 变成一个不改变其计数的 $3$-CNF。保留大小至多三的子句。对每个更长的子句，用二元 OR 门计算其析取，并要求最终输出为一。对于门 $z=u\vee v$，施加
\[
 (\neg u\vee z)\ \wedge\ (\neg v\vee z)\ \wedge\
 (u\vee v\vee\neg z).
\]
这里 $u,v$ 是输入文字或更早的门输出，$z$ 是新的。这些子句强制与 OR 相等，因此门值由原始赋值唯一确定。于是，$G$ 的一个满足赋值恰有一个到新公式的扩展，而不满足 $G$ 的赋值没有扩展。每个子句使用 $O(k)$ 个门。

接下来减少出现次数。如果变量 $x$ 在这个 $3$-CNF 中出现 $s\ge2$ 次，则把它的出现替换为不同的副本 $x_1,\ldots,x_s$，保持它们的符号，并添加 Tovey 的蕴含环~\cite[第~86 页]{Tov84}：
\[
 (\neg x_1\vee x_2)\ \wedge\cdots\wedge\
 (\neg x_{s-1}\vee x_s)\ \wedge\ (\neg x_s\vee x_1).
\]
该环强制所有副本一致。每个副本出现在一个继承自 $3$-CNF 的子句和两个环子句中，因此恰好出现在三个子句中。出现次数至多一次的变量保持不变，包括未使用的声明变量。因此这个标准的变量复制归约~\cite[命题~4.7]{Gol08} 也在满足赋值之间给出双射。最终的公式 $\widehat G$ 满足 $\#\widehat G=\#G$ 且
\[
 n\le N:=\#\{\text{$\widehat G$ 的变量}\}=O_k(Bn).
\]
它的所有子句大小至多三，包括双子句环子句，并且每个变量至多出现三次。因此它恰好是本概览所讨论情形的实例。

对每个 $\widehat G$ 应用 \eqref{eq:overview-bounded-time} 并把计数相加。当 $\sigma^{-1}$ 被 $\log n$ 的固定幂界定时，$B$ 也是如此，所得时间为
\[
 (n+m)^{O_k(1)}
 2^{O_k\!\left(\sigma n+
          Bn(\log\log n)^2/\sqrt{\log n}\right)}.
\]
例如，选择 $\sigma$ 在 $(\log n)^{-1/(2k)}$ 的常数倍之内，得到
\[
 (n+m)^{O_k(1)}
 2^{O_k\!\left(n(\log\log n)^{k+1}/(\log n)^{1/(2k)}\right)}
   = (n+m)^{O_k(1)}2^{o(n)}
\]
只使用多项式空间，逐个处理稀疏化树与归约后的实例。参数 $\sigma$ 可以取有理数。

\subsection{通过计数 \texorpdfstring{$\mathbb F_2$}{F\_2} 上稀疏低次多项式方程的解得到更快的 \texorpdfstring{$\#k$-SAT}{\#k-SAT} 算法}
\label{subsec:preview-full-case}

第~\ref{subsec:preview-general-sat} 节中的运行时间指数是 $O_k(\sigma n+B\rho(n))$，其中 $\rho(n)=n(\log\log n)^2/\sqrt{\log n}$。如果稀疏化在 $\sigma=\Theta(\rho(n)/n)$ 处给出 $B=O_k(1)$，我们就已经有一个至多相差多项式因子的、用于一般 $\#k$-SAT 的 $2^{O_k(\rho(n))}$ 时间算法：复制在每个叶上只会使用 $O_k(n)$ 个变量，而叶数也会落在同一界内。然而，在这个 $\sigma$ 选择下，保证仅为 $B=(\log n)^{O_k(1)}$。

我们分两步获得 $2^{O_k(\rho(n))}$ 界。第一，本概览的计数算法可推广到 $\mathbb F_2$ 上 $O(n)$ 个多项式方程构成的方程组，每个方程次数至多 $k$，且只有对数多项式多个单项式（定理~\ref{thm:overview-sparse-polynomials}）。第二，即使叶 $G$ 有 $Bn$ 个子句，我们也能在原始 $n$ 个变量上构造这样的方程组，其公共零点恰为 $G$ 的满足赋值。我们先用一个自包含的随机化构造说明该归约，然后解释计数算法。

\medskip\noindent
\textbf{到稀疏多项式方程的随机化归约。}
若 $x$ 违反子句 $e$，则写 $v_e(x)=1$，否则为 $0$。对子句集合 $S$ 的\emph{奇偶检验}是 $\mathbb F_2$ 上的多项式方程 $\sum_{e\in S}v_e(x)=0$：若 $x$ 违反 $S$ 中偶数个子句，它接受。每个满足赋值都通过，所以我们只需捕获每个不满足的赋值。如果 $G$ 有 $m_G\le n$ 个子句，每个子句一个检验就足够。假设 $m_G>n$。

短检验很容易漏掉单个被违反的子句，但所有只违反该子句的赋值给出相同答案。因此只有 $m_G$ 种单点子句模式需要捕获。使用 $8n$ 个独立检验。对每个检验，有放回地均匀抽取 $L=2B\lceil\log_2(2m_G)\rceil$ 个子句。对每次抽取，以概率 $1/2$ 独立地把其违反位加到奇偶和上；两个被纳入的副本相消。由于 $m_G\le Bn$，总体而言每个子句在期望上至少被对数次抽到。对于固定的 $w\ge1$ 个被违反子句的集合，令 $u=Lw/m_G$ 为每个检验的期望命中数。如果样本命中该集合，则第一次命中上的硬币以概率 $1/2$ 使奇偶为奇，即使所有其他硬币固定也是如此。因此拒绝概率为
\[
 p_w=\tfrac12\bigl(1-(1-w/m_G)^L\bigr)\ge\tfrac12(1-e^{-u}).
\]

对少量违反我们计数\emph{集合}，对大量违反我们计数\emph{赋值}：
\begin{itemize}
\item 若 $u\le1$，则 $p_w\ge u/4$。大小 $w$ 的集合至多 $m_G^w$ 个，所以至少有一个逃过所有检验的概率至多 $m_G^w e^{-8np_w}\le m_G^w e^{-2nLw/m_G}\le16^{-w}$。对 $w\ge1$ 求和，得到失败概率至多 $1/15$。
\item 若 $u>1$，则 $p_w>1/4$。对 $2^n$ 个赋值作联合界，把此情形下任何赋值逃过的概率界为 $2^n(3/4)^{8n}<1/4$。
\end{itemize}
因此以大于 $2/3$ 的概率，公共零点恰为 $G$ 的满足赋值。每个 $v_e$ 是至多 $k$ 个形如 $x_j$ 或 $1+x_j$ 的因子的乘积，所以每个检验的次数至多 $k$，单项式至多 $2^kL$ 个，关于 $n$ 是对数多项式的。第~\ref{sec:hashes} 节确定性地构造具有相同解集的短检验。我们首先使用来自扩展码的奇偶校验，以确保任何仍然通过的不满足赋值必定违反许多子句。然后一个随机短检验以至少 $1/4$ 的概率捕获每个这样的赋值。对每个候选检验，我们计数同时通过该检验与迄今所选所有检验的赋值数。我们追加使该计数最小的候选：每个满足赋值都通过每个候选，因此这至少消去剩余不满足赋值中的四分之一。

\medskip\noindent
\textbf{计数公共零点。}
回顾第~\ref{subsec:preview-encoding} 节，若 $x$ 满足子句 $e$，则 $\Phi_e(x)$ 为 $1$，否则为 $0$。当我们一次检验整个子句集合 $S_i$ 时，应该用什么替换这个因子？首先用通常的整数加法计数\emph{这个检验}中有多少子句被违反：
\[
 L_i(x)=\sum_{e\in S_i}v_e(x)=\sum_{e\in S_i}\bigl(1-\Phi_e(x)\bigr).
\]
检验恰好当 $L_i(x)$ 为偶数时通过。选择 $\ell=n+1$，使至多 $2^n$ 的计数由其模 $2^\ell$ 的余数确定。补 Beigel--Tarui 多项式 $p_\ell$（引理~\ref{lem:rational-idempotent}）在偶数整数上模 $2^\ell$ 为 $1$，在奇数整数上为 $0$。例如，违反 $S_i$ 中两个子句仍然通过检验：$p_\ell(2)\equiv1\pmod{2^\ell}$。记 $m_{\mathrm{test}}$ 为检验个数，我们的计数表达式变为
\[
 \underbrace{\sum_x\prod_{e\in G}\Phi_e(x)}_{\text{每个子句一个因子}}
 \qquad\longrightarrow\qquad
 \underbrace{\sum_x\prod_{i=1}^{m_{\mathrm{test}}} p_\ell(L_i(x))}_{\text{每个检验一个因子}}.
\]
左边是 $\#G$。右边模 $2^\ell$ 的最小非负余数是通过每个检验的赋值数。当检验捕获每个不满足赋值时，这恰好是 $\#G$。同样的编码对 $\mathbb F_2$ 上任一稀疏多项式方程 $P_i(x)=0$ 都有效：取 $L_i(x)$ 为其所列单项式的整数和。同样，方程恰好当该和为偶数时成立。

为了复用泰勒步骤，我们需要输入次数为 $O(n)$ 的乘积。直接相乘 $m_{\mathrm{test}}$ 个检验因子只给出界 $O(km_{\mathrm{test}}\ell)$，可能是 $n$ 的二次。关键在于\emph{跨越所有检验一起}截断乘积。第~\ref{subsec:rational-lift} 节把 $p_\ell(y)$ 写成项 $e_j(y)h(y)^j$ 之和，其中 $e_j$ 次数至多一，$h(y)=y(y-1)$ 在每个整数处为偶。因此，包含 $\ell$ 个或更多 $h$ 副本的乘积项在布尔赋值上模 $2^\ell$ 为零，即使这些副本来自不同的检验。我们只保留总副本数少于 $\ell$ 的项，并在相乘时跟踪该数目。每个保留项在 $L_i$ 中至多有 $m_{\mathrm{test}}$ 个线性因子和 $\ell-1$ 个二次因子。由于 $\deg L_i\le k$，其输入次数至多为 $k(m_{\mathrm{test}}+2(\ell-1))=O(n)$。

对于替换形式，重要的是每个因子读取哪些变量。第~\ref{subsec:preview-deformation} 节为子句中的每个变量使用一个形式，而这里我们为检验中的每个变量使用一个形式，并在整个检验中复用它。每个稀疏方程只读取对数多项式多个变量。对 $O(n/\sqrt{\log n})$ 个最频繁变量的值进行分支，使每个剩余变量只出现在对数多项式多个检验中，正如第~\ref{sec:construction} 节所要求。在每个分支上，固定共享测量再次让我们一次求和一个组。由于截断乘积的输入次数为 $O(n)$，次数为 $R$ 的替换给出 $\deg_T H=O(nR)$：每点同样的 $O(1+nR/Q)$ 个泰勒系数恢复 $T=0$ 处具有 $0,1,\ldots,\ell-1$ 个 $h$ 副本的贡献。它们模 $2^\ell$ 之和，取最小非负余数，就是精确的分支计数。各分支计数相加即得公共零点数。

\begin{theorem}[计数稀疏多项式的公共零点]
\label{thm:overview-sparse-polynomials}
固定整数 $k\ge2$ 与 $c\ge1$。设 $P_1,\ldots,P_{m_{\mathrm{test}}}$ 是 $\mathbb F_2$ 上关于 $n\ge20$ 个布尔变量的多项式，且 $m_{\mathrm{test}}\le cn$。每个 $P_i$ 以次数至多 $k$ 的至多 $(\log n)^c$ 个单项式的列表给出。$\{0,1\}^n$ 中公共零点的精确数目可在 $2^{O_{k,c}(n(\log\log n)^2/\sqrt{\log n})}$ 次位运算和 $n^{O_{k,c}(1)}$ 空间内确定性计算。
\end{theorem}
第~\ref{sec:counting} 与~\ref{sec:construction} 节提供计数例程与替换形式；第~\ref{sec:setup} 节用它们对多项式方程直接证明该定理。

对于为 $G$ 构造的最终方程组，公共零点数为 $\#G$。第~\ref{sec:hashes} 节对每个叶只使用 $2^{(\log n)^{O_k(1)}}$ 次计数例程调用，这保持了定理~\ref{thm:overview-sparse-polynomials} 中的运行时间指数。最后，在第~\ref{subsec:general-reduction} 节中选择阶为 $(\log\log n)^2/\sqrt{\log n}$ 的 $\sigma$，使 $B$ 保持对数多项式，并让叶数保持在 $2^{O_k(n(\log\log n)^2/\sqrt{\log n})}$ 之内。它们的满足赋值互不相交，因此相加其计数给出命题~\ref{prop:interpolation-bound} 中的界。

\section{通过共享线性摘要的快速计数算法}\label{sec:counting}
第~\ref{sec:counting}--\ref{sec:setup} 节证明定理~\ref{thm:overview-sparse-polynomials}。我们从第~\ref{subsec:preview-encoding} 节的和开始：对每个布尔赋值，相乘各因子，然后对所有赋值求乘积之和。这里因子可以是任意有理多项式。我们展示在给定替换形式与测量时如何计算该和。第~\ref{sec:construction} 节只使用每个因子所读取的变量来构造它们。第~\ref{sec:setup} 节为每个稀疏多项式方程提供一个因子。

\subsection{多项式因子与待恢复的值}
\label{subsec:product-interface}
令 $\mathcal F$ 索引各因子，并令 $I_f\subseteq[n]$ 为 $\Phi_f$ 所读取的变量。对于有理系数，我们的目标是
\begin{equation}\label{eq:general-product-sum}
 Z=\sum_{x\in\{0,1\}^n}\prod_{f\in\mathcal F}\Phi_f((x_v)_{v\in I_f}).
\end{equation}
对于子句满足指示，$Z$ 是满足赋值的数目。该证明对任意交换 $\mathbb Q$-代数 $\mathcal A$ 中的系数也成立：此时和与积在 $\mathcal A$ 中进行，而插值只用到有理标量。读者可在本节始终取 $\mathcal A=\mathbb Q$。

不同因子可以对同一变量使用不同的替换。给每个关联 $(f,v)$ 各自的未定元 $Y_{f,v}$，并假设
\begin{equation}\label{eq:general-product-degree}
 \deg_{\{Y_{f,v}\}}\!
    \left(\prod_{f\in\mathcal F}\Phi_f((Y_{f,v})_{v\in I_f})\right)\le d .
\end{equation}
这是整个乘积在 $\mathcal A$ 上计算所得次数。选择次数至多 $R$ 且 $\vec P_{f,v}(0)=\delta_v$ 的替换行 $\vec P_{f,v}(T)$，并构造
\begin{equation}\label{eq:general-count-polynomial}
 H(T)=\sum_{x\in\{0,1\}^n}\prod_{f\in\mathcal F}
   \Phi_f\bigl((\langle\vec P_{f,v}(T),x\rangle)_{v\in I_f}\bigr)
   \in\mathcal A[T].
\end{equation}
替换在零点恢复输入，并把次数至多乘以 $R$。于是 $H(0)=Z$ 且 $\deg_T H\le d_H:=dR$。我们的应用将有 $d=O(n)$，如有界出现次数例子中那样。

\subsection{所需行、划分与测量}
\label{subsec:layout-interface}
我们使用与第~\ref{subsec:preview-separation} 节相同的设计，用每个因子的组指派代替每个子句的组指派。下面的最终条件记录了精确算术所需的系数界。这里 $\widetilde\xi$ 表示从组到全部 $n$ 个坐标的零延拓。

\begin{definition}\label{def:layout}
对于因子作用域 $(I_f)_{f\in\mathcal F}$，令 $r\ge2$、$D\ge1$ 与 $1\le R<Q$ 为整数，并令 $0<\tau\le1/2$。我们需要以下数据与性质。
\begin{enumerate}
\item 对每个 $t\in[Q]$，一个划分
$[n]=S_{t,1}\sqcup\cdots\sqcup S_{t,r}$，允许空的部分，且 $|S_{t,g}|\le2n/r$。
\item 对每个因子 $f\in\mathcal F$，一个序列
$\sigma_f\in[r]^Q$。值 $\sigma_f(t)$ 把因子 $f$ 指派到点 $t$ 处的一个组。
\item 对每个 $t\in[Q]$ 与 $g\in[r]$，一个整数矩阵
$C_{t,g}\in\mathbb Z^{[c_{t,g}]\times S_{t,g}}$，其中 $c_{t,g}$ 是非负整数。每行的绝对值之和至多 $D$，且
\[
 \sum_{g=1}^r c_{t,g}\le2\tau n.
\]
对赋值 $x$，$C_{t,g}(x|_{S_{t,g}})$ 的条目是组 $g$ 上的测量。
\item 对每个 $f\in\mathcal F$ 与 $v\in I_f$，一行
$\vec{P}_{f,v}(T)\in\mathbb Q[T]^{1\times n}$，满足
$\vec{P}_{f,v}(0)=\delta_v$ 且 $\deg \vec{P}_{f,v}\le R$。对每个 $t\in[Q]$、$g\in[r]\setminus\{\sigma_f(t)\}$ 与
$\xi\in\mathbb Q^{S_{t,g}}$，要求
\[
 C_{t,g}\xi=\vec{0}\quad\Longrightarrow\quad
 \langle \vec{P}_{f,v}(t), \widetilde\xi\rangle=0,
\]
其中 $\widetilde\xi\in\mathbb Q^n$ 是 $\xi$ 的零延拓。
\item 令 $d_{\mathrm{den}}=(R-1)!Q!$，$d_{\mathrm{den}} \vec{P}_{f,v}$ 的每个系数都是具有 $O(Q\log Q+\log D)$ 位的整数，隐含常数为绝对常数。
\end{enumerate}
\end{definition}
为理解核条件，固定 $t$、$f$ 与 $v\in I_f$，并令 $g=\sigma_f(t)$。记 $x_h=x|_{S_{t,h}}$。
若 $x,x'\in\mathbb Q^n$ 在某个 $h\ne g$ 的 $S_{t,h}$ 外相同，且满足 $C_{t,h}x_h=C_{t,h}x'_h$，则
$x'_h-x_h\in\ker C_{t,h}$。该条件给出
\[
 \langle \vec{P}_{f,v}(t), x'\rangle=\langle \vec{P}_{f,v}(t), x\rangle.
\]
逐个群组应用这一点表明，元组
$(\langle \vec{P}_{f,v}(t), x\rangle)_{v\in I_f}$ 由 $x_g$ 与
$(C_{t,h}x_h)_{h\ne g}$ 决定。因此，固定所有测量向量
使整个因子成为 $x_g$ 的函数。这适用于每个
多项式因子 $\Phi_f$，包括系数在
$\mathcal A$ 中的因子。

布尔向量的每个测量值都落在 $[-D,D]\cap\mathbb Z$ 中。
我们限制矩阵元素绝对值之和，因为
构造可能把同一个布尔坐标放入多个
辅助位置。识别这些位置后，
$C_{t,g}$ 的元素不必属于 $\{0,1,-1\}$。仅靠秩界
无法控制可能测量值的数量：单独一行
$(1,2,4,\ldots,2^{n-1})$ 就能区分全部 $2^n$ 个布尔向量。

\subsection{通过固定测量分离求和}
\label{subsec:measurement-factorization}
固定测量使每个因子只依赖于其被分配的
群组。下面的引理给出在交换环上得到的计算，因此它同样适用于有理值因子与
$\mathcal A$ 中的因子。限制 $x_i=1$ 将用于
泰勒系数，并且是各个群组局部的。

\begin{lemma}\label{lem:measurement-separation}
设 $[n]=S_1\sqcup\cdots\sqcup S_r$ 是一个划分，允许空的部分。
对每个 $g\in[r]$，设 $C_g$ 是一个整数矩阵，具有 $c_g$ 行，
列由 $S_g$ 索引。假设每一行的绝对值
之和至多为整数 $D\ge0$，并记 $c=\sum_g c_g$。
设 $\mathcal A$ 是带单位的交换环，并设
$A_g:\mathbb Q^n\to\mathcal A$ 满足
\[
 A_g(x+\xi)=A_g(x)
 \qquad
 (x\in\mathbb Q^n,\ \xi\in\ker C_h,\ h\ne g),
\]
其中 $\xi$ 在 $S_h$ 外延拓为零。
对任意 $E\subseteq[n]$，和
\[
 Z_E=\sum_{\substack{x\in\{0,1\}^n\\x_i=1\ (i\in E)}}
           \prod_{g=1}^r A_g(x)
\]
可以通过检查至多
\[
 2(2D+1)^c\sum_{g=1}^r2^{|S_g|}
\]
个布尔向量来计算，至多进行一次 $A_g$ 求值，
$\poly(n,r,c,\log(D+2))$ 次额外位运算，以及每个向量 $O(1)$ 次环运算。除求值器自身的存储与
$O(1)$ 个环元素外，该过程使用 $\poly(n,r,c,\log(D+2))$
位空间。
\end{lemma}
\begin{proof}
对 $y_g\in\mathbb Z^{c_g}$，定义
\[
 \mathcal X_g(y_g)=\{u\in\{0,1\}^{S_g}:C_gu=y_g\},
 \qquad
 \mathcal Y=\prod_{g=1}^r\{-D,\ldots,D\}^{c_g}.
\]
布尔赋值的每个测量元组都属于
$\mathcal Y$，且 $|\mathcal Y|=(2D+1)^c$。
我们首先确定固定测量元组
$\vec y=(y_1,\ldots,y_r)$ 的贡献。若每个 $\mathcal X_g(y_g)$ 都非空，
对每个 $g$ 取代表元 $\bar x_g\in\mathcal X_g(y_g)$，并令
\[
 \mathcal X_{g,E}(y_g)=
 \{u\in\mathcal X_g(y_g):u_i=1\text{ 对所有 }i\in E\cap S_g\}.
\]
若对每个 $h$ 都有 $x_h\in\mathcal X_h(y_h)$，则
$x_h-\bar x_h\in\ker C_h$。对 $h\ne g$ 依次应用不变性假设可得
\[
 A_g\Bigl(\sum_{h=1}^r x_h\Bigr)
 =A_g\Bigl(x_g+\sum_{h\ne g}\bar x_h\Bigr).
\]
具有测量元组 $\vec y$ 且由
$E$ 索引的固定位的赋值构成 $\prod_g\mathcal X_{g,E}(y_g)$。定义
\begin{equation}\label{eq:group-sum}
 \Sigma_g(\vec y)=\sum_{x_g\in\mathcal X_{g,E}(y_g)}
       A_g\Bigl(x_g+\sum_{h\ne g}\bar x_h\Bigr).
\end{equation}
分配律表明这些赋值对 $Z_E$ 的贡献为
$\prod_g\Sigma_g(\vec y)$。若 $\mathcal X_{g,E}(y_g)$ 为空，
则 $\Sigma_g(\vec y)=0$。代表元 $\bar x_g$ 本身
不必满足由 $E$ 索引的固定位条件。
每个布尔赋值都有唯一的测量元组，因此
\[
 Z_E=\sum_{\vec y\in\mathcal Y}\prod_{g=1}^r\Sigma_g(\vec y),
\]
其中若任一 $\mathcal X_g(y_g)$
为空，则乘积取为零。

算法一次处理一个测量元组来求值此恒等式，
从 $Z=0_{\mathcal A}$ 开始。对每个 $\vec y\in\mathcal Y$，
对每个 $\{0,1\}^{S_g}$ 的首次扫描找到一个代表元
$\bar x_g\in\mathcal X_g(y_g)$；若任一这样的集合为空，则
算法跳过该元组。否则初始化
$p=1_{\mathcal A}$。为计算 $\Sigma_g(\vec y)$，初始化
$s=0_{\mathcal A}$ 并再次扫描 $u\in\{0,1\}^{S_g}$。每当
\[
 C_gu=y_g,\qquad u_i=1\quad(i\in E\cap S_g),
\]
算法执行更新
\[
 s\leftarrow s+A_g\Bigl(u+\sum_{h\ne g}\bar x_h\Bigr).
\]
扫描结束时，$s=\Sigma_g(\vec y)$，算法令
$p\leftarrow ps$。处理完所有群组后，令 $Z\leftarrow Z+p$。
处理完所有元组 $\vec y$ 后，$Z=Z_E$。

对每个 $\vec y$，寻找代表元的扫描与
计算 $\Sigma_g(\vec y)$ 的扫描各自检查至多
$\sum_g2^{|S_g|}$ 个布尔向量。这给出所述界。
检查一个向量的测量值与固定位使用
$\poly(n,r,c,\log(D+2))$ 次位运算：每个矩阵行的
绝对值之和至多为 $D$，因此其部分和有
$O(\log(D+2))$ 位。每个向量至多进行一次因子求值。
保留的状态为 $\vec y$、$(\bar x_1,\ldots,\bar x_r)$、$g$、$u$、
$s$、$p$ 与 $Z$。元组 $\vec y$ 包含 $c$ 个位于 $[-D,D]$ 中的整数，
代表元总共包含 $n$ 位，且 $s,p,Z\in\mathcal A$。
\end{proof}

为在固定 $t$ 处计算 \eqref{eq:general-count-polynomial}，按 $\sigma_f(t)$
将变换后的因子分组，并在
$\mathcal A$ 中取 $E=\varnothing$ 应用该引理。定义~\ref{def:layout}
给出其不变性假设，因此和为 $H(t)$。如概述中那样，两个指数
代价分别是测量选择与一个群组内的布尔赋值。

\subsection{由固定布尔替换得到泰勒系数}
\label{subsec:fixed-boolean-substitutions}
$Q$ 个值 $H(1),\ldots,H(Q)$ 不必确定 $H$。
令
\begin{equation}\label{eq:hasse-order}
 q=\left\lfloor\frac{d_H}{Q}\right\rfloor,\qquad J=q+1.
\end{equation}
由于 $QJ>d_H$，只需计算
\begin{equation}\label{eq:taylor-truncation}
 \mathcal H_t(z)=\sum_{j=0}^{J-1}([z^j]H(t+z))z^j
 \qquad(t\in[Q]).
\end{equation}
Hermite 插值从 $\mathcal A$ 中的这些
$QJ$ 个系数恢复 $H$，从而恢复 $Z=H(0)$。事实上，次数小于 $QJ$ 的多项式的
有理 Hermite 矩阵在
$\mathbb Q$ 上可逆。其逆通过有理
线性组合作用于 $\mathcal A$ 的元素。

固定 $t$。$t+z$ 处的替换形式不必遵循
$t$ 处的设计。记
\[
 \Delta\vec P_{f,v}(z)=\vec P_{f,v}(t+z)-\vec P_{f,v}(t).
\]
每个元素都能被 $z$ 整除。对固定的 $x\in\mathbb Q^n$，
引入独立不定元 $\vec X$，并将每个输入替换为
\begin{equation}\label{eq:jet-substitution}
 \langle\vec P_{f,v}(t),x\rangle
       +\langle\Delta\vec P_{f,v}(z),\vec X\rangle.
\end{equation}
令
\begin{equation}\label{eq:auxiliary-polynomial}
 \mathcal G_x(\vec X,z)=\prod_{f\in\mathcal F}
 \Phi_f\!\left(
  (\langle\vec P_{f,v}(t),x\rangle+
   \langle\Delta\vec P_{f,v}(z),\vec X\rangle)_{v\in I_f}\right).
\end{equation}
令 $\vec X=x$ 可得
$\sum_x\mathcal G_x(x,z)=H(t+z)$。
反过来固定 $\vec X$ 使每个修正独立于
$x$，因此 $t$ 处的设计仍能分离因子。
此外，$\vec X$ 的每次出现都带有一个 $z$ 因子，从而
\begin{equation}\label{eq:degree-cost}
 \deg_{\vec X}([z^j]\mathcal G_x(\vec X,z))\le j.
\end{equation}
下一个恒等式从固定布尔替换与
对被求和赋值的限制中恢复替换 $\vec X=x$。

\begin{lemma}\label{lem:boolean-differences}
设 $q\ge0$ 为整数，并设 $\mathcal A$ 是
带单位的交换环。对每个 $x\in\{0,1\}^n$，设
$p_x(X_1,\ldots,X_n)$ 是 $\mathcal A$ 上的多项式，次数
至多为 $q$。则
\begin{align}
 \sum_{x\in\{0,1\}^n}p_x(x)
 &=\sum_{\substack{E\subseteq[n]\\|E|\le q}}
   \sum_{U\subseteq E}(-1)^{|E|-|U|}
   \sum_{\substack{x\in\{0,1\}^n\\x_i=1\ (i\in E)}}
                  p_x(\mathbf1_U).
\label{eq:boolean-differences}
\end{align}
这里 $\mathbf1_U$ 是 $U$ 的指示向量。
\end{lemma}
\begin{proof}
对固定的 $x$，将 $X_i$ 的每个正幂替换为 $X_i$。
这不改变布尔赋值下的值，也不增加
次数。通过对子集作 M\"obius 反演，所得的
多线性多项式具有展开式
\[
 \widetilde p_x(\vec{X})=\sum_{E\subseteq[n]}c_{x,E}\prod_{i\in E}X_i,
 \qquad
 c_{x,E}=\sum_{U\subseteq E}(-1)^{|E|-|U|}p_x(\mathbf1_U).
\]
次数界给出：只要 $|E|>q$，就有 $c_{x,E}=0$。
令 $\vec{X}=x$ 并对 $x$ 求和，即得所述恒等式：
单项式 $\prod_{i\in E}x_i$
恰好选出对所有 $i\in E$ 都有 $x_i=1$ 的赋值。
\end{proof}

对每个 $j<J$，将引理~\ref{lem:boolean-differences} 应用于
$\mathcal A$ 上的 $[z^j]\mathcal G_x$。
所有系数可以使用相同的对 $U\subseteq E$、$|E|\le q$，
因此在
\[
 \mathcal R=\mathcal A[z]/(z^J).
\]
中一起计算它们。
对固定的对 $(E,U)$，令 $\Phi_{f,U}(x;z)$ 表示在令 $\vec X=\mathbf1_U$ 之后
\eqref{eq:auxiliary-polynomial} 中的因子 $f$。
按赋值将这些因子分组：
\begin{equation}\label{eq:group-polynomial}
 A_{g,U}(x;z)=\prod_{f:\sigma_f(t)=g}\Phi_{f,U}(x;z).
\end{equation}
它们的乘积是 $\mathcal G_x(\mathbf1_U,z)$。
对 $h\ne g$，将 $x$ 改为 $\ker C_{t,h}$ 中的向量既不改变
这些因子的原始输入，也不改变它们的固定
修正。因此引理~\ref{lem:measurement-separation}
通过检查至多
\begin{equation}\label{eq:boolean-enumeration}
 2r(2D+1)^{2\tau n}2^{2n/r}
\end{equation}
个布尔向量来计算带 $i\in E$ 时 $x_i=1$ 的限制和。合并所有对可得
\begin{equation}\label{eq:full-taylor-sum}
 H(t+z)\equiv
 \sum_{\substack{E\subseteq[n]\\|E|\le q}}
 \sum_{U\subseteq E}(-1)^{|E|-|U|}
 \sum_{\substack{x\in\{0,1\}^n\\x_i=1\ (i\in E)}}
       \mathcal G_x(\mathbf1_U,z)\pmod{z^J}.
\end{equation}
群组乘积与求和均在 $\mathcal R$ 中形成。结果为
$\mathcal H_t(z)$；在所有 $t$ 处重复即可提供插值数据。
算法从不存储关于全部 $n$ 个不定元的多项式：
它先令 $\vec X=\mathbf1_U$，并且只存储 $\mathcal A$ 中的 $J$ 个系数。

\subsection{三项指数代价}
\label{subsec:three-costs}
对每个 $t\in[Q]$，
满足 $|E|\le q$ 的对 $U\subseteq E\subseteq[n]$ 的数量为
\begin{equation}\label{eq:pair-count}
 C_q=\sum_{s=0}^{\min(q,n)}2^s\binom ns.
\end{equation}
特别地，$C_0=1$，且当 $q\ge n$ 时 $C_q=3^n$。
对每个 $0<\eta\le1$，
\begin{equation}\label{eq:finite-difference-cost}
 \eta^qC_q\le\sum_{s=0}^n\binom ns(2\eta)^s=(1+2\eta)^n,
\end{equation}
因为对每个保留的指标 $s\le q$ 都有 $\eta^q\le\eta^s$。

\begin{proposition}[通用计数过程]\label{prop:general-count}
设 $\mathcal A$ 为交换 $\mathbb Q$-代数。假设在
$\mathcal A$ 上计算的 \eqref{eq:general-product-sum} 中因子乘积
满足次数界 \eqref{eq:general-product-degree}，
且替换行、划分与测量满足
定义~\ref{def:layout}。其中 $q,J$ 如 \eqref{eq:hasse-order} 所示，
该过程通过检查至多
\[
 2Qr\,C_q(2D+1)^{2\tau n}2^{2n/r}
\]
个布尔向量来计算 $Z$。对每个向量，它至多在 $\mathcal A[z]/(z^J)$ 中
对 $|\mathcal F|$ 个因子求值。每个向量的额外运算与最终插值需要多项式多次运算，
分别作用于 $\mathcal A$ 与 $\mathbb Q$ 上，参数为
$n,|\mathcal F|,D,Q,R,J,r$ 以及总作用域长度。
除因子求值器自身的存储外，该过程仅保留
多项式多个代数元素、有理数与索引位。
\end{proposition}
\begin{proof}
共有 $Q$ 个点，每个点 $C_q$ 对，每个对至多
\eqref{eq:boolean-enumeration} 个布尔向量。
形成每个群组乘积至多使用 $|\mathcal F|$ 次因子求值。
模 $z^J$ 的加法与乘法在 $\mathcal A$ 中的系数上花费 $O(J^2)$ 次运算。
前面的泰勒恒等式
与插值证明正确性。
顺序处理点、对、测量元组与局部赋值，
保留它们当前的索引与累加器。
连同 $QJ$ 个插值系数，这给出
所述存储界。位复杂度还依赖于
因子求值的代价与系数大小；这些将在
第~\ref{subsec:one-query-cost} 节的稀疏多项式应用中验证。
\end{proof}

当 $d=O(n)$ 时，有 $d_H=O(nR)$ 且 $q=O(nR/Q)$。
在 \eqref{eq:finite-difference-cost} 中取 $\eta=R/Q$ 可得
\[
 \log C_q\le q\log(Q/R)+n\log(1+2R/Q)
     =O\!\left(n\frac RQ\log\frac{2Q}{R}\right).
\]
因此，若因子求值与系数算术具有
多项式位代价，则运行时间是多项式因子乘以
\begin{equation}\label{eq:three-costs}
 2^{O\left(n\left[
       \frac RQ\log\frac{2Q}{R}
       +\tau\log(2D+1)+\frac1r\right]\right)}.
\end{equation}
这恰好是概述中的三项代价。引理~\ref{lem:incidence-parameters} 中的参数
选择使其和在因子作用域及其
最大关联数都被 $\log n$ 的固定幂次界定时为
$O((\log\log n)^2/\sqrt{\log n})$。
当 $d$ 更大时该过程也适用；此时在
命题~\ref{prop:general-count} 的界中保留
实际值 $q=\lfloor dR/Q\rfloor$。

到目前为止，我们假设设计已给定。下一
节从因子作用域构造它。剩余
两节先使用所得过程计算一个给定多项式系统的零点，
然后为一个公式选择这样的系统。

\section{构造低次替换形式}\label{sec:construction}
计数过程需要为每个因子 $f$ 读取变量 $v$ 的对 $(f,v)$
提供一个替换形式。对于
第~\ref{sec:setup} 节中的应用，$f$ 是一个多项式方程，其作用域是
其各单项式变量集合的并。本节只使用
这些作用域。
我们在定义~\ref{def:layout} 中构造行 $\vec{P}_{f,v}(T)$ 与测量矩阵 $C_{t,g}$，
以及划分 $S_{t,g}$
与因子分配 $\sigma_f$。主要困难在于保持
$\deg \vec{P}_{f,v}\le R$，且 $R$ 远小于 $Q$。$H$ 的次数界
随 $R$ 线性增长，因此较小的 $R/Q$ 会减少
每个点所需的泰勒系数数量。若在 $T=0,1,\ldots,Q$ 处独立选择这些行
并进行插值，通常会得到
次数 $Q$。我们改为同时构造它们在这些点处的值。

构造首先在一个完整的辅助坐标阵列上给出多项式行。然后我们将每个因子--变量对放入该阵列的一个副本中。若干位置可以包含同一个布尔
变量 $x_v$；它们的值始终相等，并且在计数求和中不会成为
独立变量。

\subsection{辅助坐标}\label{subsec:incidence-layout}
我们首先指定更大的坐标集 $\mathcal B$，辅助行
$\vec{P}_u(T)$ 将作用于其上，以及将实际对 $(f,v)$ 放入其副本所需的
关联界。设 $(I_f)_{f\in\mathcal F}$ 是因子作用域，其中 $I_f\subseteq[n]$，
并记
\begin{equation}\label{eq:factor-incidences}
 \begin{gathered}
 \mathcal A=\{(f,v):f\in\mathcal F,\ v\in I_f\},\qquad M=|\mathcal A|,\\
 |I_f|\le w\quad(f\in\mathcal F),\qquad
 \#\{f:v\in I_f\}\le b\quad(v\in[n]).
 \end{gathered}
\end{equation}
因此 $w$ 界定一个因子作用域中变量的数量，$b$
界定包含一个变量的因子数量。作用域界
$w$ 不必界定因子多项式的次数。作用域为空的因子是常数，不需要辅助位置；在每个点将这样的
因子分配给第一组。如有必要，对每个 $v\in[n]$ 添加一个作用域为 $\{v\}$ 的不同的
常一因子。这既不改变乘积也不改变计数，并且至多使最大关联数
增加一。以 $b$ 表示此添加后的界，我们假设
\[
 n\ge1,\qquad b,w\ge1,\qquad n\le M\le bn.
\]
本节中，$r\ge2$ 与 $1\le R<Q$ 是整数，且
$0<\tau\le1/2$。

在每个点 $t\in[Q]$，变量 $x_v$ 将属于一个群组，
而包含 $v$ 的因子可能被分配给另一个群组。我们将选择
\[
 a_v\in[r]^Q\quad(v\in[n]),\qquad
 d_f\in[r]^Q\quad(f\in\mathcal F),
\]
其中 $a_v(t)$ 指定包含 $x_v$ 的群组，$d_f(t)$
指定分配给因子 $f$ 的群组。$(f,v)$ 的位置
将具有标签 $(a_v,d_f)$。特别地，所有包含
$x_v$ 的位置都将具有相同的第一序列 $a_v$。

在选择 $a_v$ 与 $d_f$ 之前，取全部可能的标签：
\begin{equation}\label{eq:auxiliary-labels}
 \mathcal B=([r]^Q)^2,\qquad u=(a_u,d_u),\qquad D=|\mathcal B|=r^{2Q}.
\end{equation}
辅助向量 $\zeta\in\mathbb Q^{\mathcal B}$ 对每个对 $u$ 有一个有理
坐标 $\zeta_u$。标签 $u$ 在 $\zeta$ 变化时保持
固定。对 $t\in[Q]$ 与 $g\in[r]$，定义
\begin{equation}\label{eq:auxiliary-groups}
 S^{\mathcal B}_{t,g}=\{u\in\mathcal B:a_u(t)=g\}.
\end{equation}
对固定的 $t$，这些集合将 $\mathcal B$ 划分为 $r$ 个大小为 $D/r$ 的集合。

对每个 $u\in\mathcal B$，我们寻求一个多项式行
$\vec{P}_u(T)\in\mathbb Q[T]^{1\times\mathcal B}$，其次数至多 $R-1$，
且 $\vec{P}_u(0)=\delta_u$。在 $T=t$ 处，值 $\langle \vec{P}_u(t), \zeta\rangle$ 应由
\[
 \zeta|_{S^{\mathcal B}_{t,d_u(t)}}
 \quad\text{以及}\quad C^{\mathcal B}\zeta
\]
确定，其中 $C^{\mathcal B}$ 是一个整数矩阵，独立于 $u,t$，
元素在 $\{0,1,-1\}$ 中且行数很少。之后，使用
$u_{f,v}=(a_v,d_f)$ 将使因子 $f$ 的每个输入使用群组
$d_f(t)$。$a_v,d_f$ 的选择以及放入副本的安排分别推迟到
第~\ref{subsec:choosing-labels} 节
与第~\ref{subsec:embedding} 节。

\subsection{矩阵 $M_0$ 与 $C^{\mathcal B}$}
为构造行 $\vec{P}_u(T)$，我们寻求一个矩阵
$M_0\in\mathbb Q^{\mathcal B\times\mathcal B}$，它固定
$\ker C^{\mathcal B}$ 并满足
\begin{equation}\label{eq:matrix-support-target}
 M_0[u,y]\ne0\quad\Longrightarrow\quad
 \bigl|\{s\in[Q]:a_y(s)\ne d_u(s)\}\bigr|\le R-1.
\end{equation}
恒等式
\begin{equation}\label{eq:coordinate-decomposition}
 \zeta_u=((I-M_0)\zeta)_u+
          \sum_{y\in\mathcal B}M_0[u,y]\zeta_y
\end{equation}
解释了两个要求。若 $M_0$ 固定 $\ker C^{\mathcal B}$，
则第一项由 $C^{\mathcal B}\zeta$ 决定；我们将使其
独立于 $T$。在求和中，将每个非零
$M_0[u,y]$ 乘以一个在零处为一且在满足 $a_y(t)\ne d_u(t)$ 的每个
$t$ 处为零的多项式。界
\eqref{eq:matrix-support-target} 允许次数至多 $R-1$。
在 $T=t$ 处，求和将只使用坐标
$y\in S^{\mathcal B}_{t,d_u(t)}$。
对 $s\in[Q]$，令 $\vartheta_s:\mathcal B\to\mathcal B$ 交换条目 $a_u(s)$
与 $d_u(s)$，其余所有条目保持不变。
定义其置换矩阵 $\Pi_s$ 为
\[
 (\Pi_s\zeta)_u=\zeta_{\vartheta_s(u)}.
\]
交换后的对 $\vartheta_s(u)$ 属于 $\mathcal B$，因为该集合
包含 $[r]^Q$ 中的每一对序列。矩阵 $\Pi_s$
两两可交换并满足 $\Pi_s^2=I$。它们的乘积
\[
 V_0=\prod_{s=1}^Q\Pi_s
\]
交换完整序列 $a_u$ 与 $d_u$。于是
$V_0[u,y]\ne0$ 蕴含对每个 $s$ 有 $a_y(s)=d_u(s)$；它
无失配地满足
\eqref{eq:matrix-support-target}。我们将修改
$V_0$，使所得矩阵还固定一个短整数测量列表的核。

定义
\begin{equation}\label{eq:swap-weight}
 W_0=\frac12\sum_{s=1}^Q(I-\Pi_s).
\end{equation}
在 $\Pi_s$ 的公共特征向量 $\psi$ 上，令 $j$ 为
满足 $\Pi_s\psi=-\psi$ 的下标 $s$ 的个数。则
\[
 W_0\psi=j\psi,\qquad V_0\psi=(-1)^j\psi,
 \qquad V_0p(W_0)\psi=(-1)^jp(j)\psi
\]
对每个多项式 $p$ 成立。因此，只要
$p(j)=(-1)^j$，$V_0p(W_0)$ 就固定 $\psi$。我们将构造一组
显式的公共特征基，其中具有给定 $j$ 的向量所占比例为
$\Pr[W=j]$，$W\sim\operatorname{Bin}(Q,(r-1)/(2r))$。

对区间
$\mathcal I=\{A_0,\ldots,A_0+R-1\}\subseteq\{0,\ldots,Q\}$，
我们可以用一个次数至多 $R-1$ 的多项式对 $j\in\mathcal I$
强制等式 $f_{\mathcal I}(j)=(-1)^j$ 成立。
我们使用显式插值公式
\begin{equation}\label{eq:swap-matrix-definition}
 f_{\mathcal I}(z)=(-1)^{A_0}\sum_{j=0}^{R-1}(-2)^j
                          \binom{z-A_0}{j},
 \qquad M_0=V_0f_{\mathcal I}(W_0),
\end{equation}
其中
\[
 \binom{z-A_0}{j}=\frac1{j!}\prod_{a=0}^{j-1}(z-A_0-a)
\]
且空乘积为一。对 $0\le q<R$，二项式定理给出
$f_{\mathcal I}(A_0+q)=(-1)^{A_0}(1-2)^q=(-1)^{A_0+q}$。

\begin{lemma}\label{lem:swaps-matrix}
对如上所述的 $\mathcal B$、$\mathcal I$ 与 $M_0$，存在矩阵
$C^{\mathcal B}\in\{0,1,-1\}^{[c_{\mathcal B}]\times\mathcal B}$，使得
\[
 c_{\mathcal B}\le D\Pr[W\notin\mathcal I],\qquad
 W\sim\operatorname{Bin}\!\left(Q,\frac{r-1}{2r}\right),
\]
且对每个 $\xi\in\ker C^{\mathcal B}$ 有 $M_0\xi=\xi$。
\end{lemma}
\begin{proof}
对 $(\alpha,\delta)\in[r]^2$，令 $e_{\alpha\delta}$ 为
该有序对的指示函数。交换其两个条目
具有正交特征基
\[
 e_{\alpha\alpha}\quad(\alpha\in[r]),\qquad
 e_{\alpha\delta}+e_{\delta\alpha},\quad
 e_{\alpha\delta}-e_{\delta\alpha}\quad(1\le\alpha<\delta\le r).
\]
前两类的特征值为 $+1$，第三类的特征值为
$-1$。特征值为 $+1$ 的向量有 $r(r+1)/2$ 个，
特征值为 $-1$ 的有 $r(r-1)/2$ 个。

对每个 $s\in[Q]$，选取这些函数之一
$\psi_s:[r]^2\to\{0,1,-1\}$，并定义
\[
 \psi(u)=\prod_{s=1}^Q\psi_s(a_u(s),d_u(s))
 \qquad(u\in\mathcal B).
\]
$D=r^{2Q}$ 种选择给出非零的正交向量：它们的内积
分解为函数 $\psi_s$ 的内积之积。
它们构成 $\mathbb Q^{\mathcal B}$ 的一组基，元素属于
$\{0,1,-1\}$。若所选的 $\psi_s$ 中恰有 $j$ 个特征值为
$-1$，则 $W_0\psi=j\psi$ 且 $V_0\psi=(-1)^j\psi$。因此 $M_0$
在 $j\in\mathcal I$ 时固定 $\psi$。

为使 $\ker C^{\mathcal B}$ 成为 $M_0$ 所固定的空间，我们
将 $C^{\mathcal B}$ 的行定义为满足 $j\notin\mathcal I$ 的基向量
的转置。
每个位置上的选择中有 $(r-1)/(2r)$ 的比例特征值为 $-1$，且
各位置相互独立。因此保留为行的 $D$ 个基向量所占比例为
$\Pr[W\notin\mathcal I]$，这证明了行数界。若
$C^{\mathcal B}\xi=\vec{0}$，则 $\xi$ 与每个保留的
基向量正交。它是满足 $j\in\mathcal I$ 的基向量的线性组合，
而 $M_0$ 固定所有这些基向量。
\end{proof}

对界 $\epsilon>0$，条件
\begin{equation}\label{eq:incidence-tail}
 \Pr\!\left[\operatorname{Bin}\!\left(Q,\frac{r-1}{2r}\right)
                         \notin\mathcal I\right]\le\epsilon
\end{equation}
于是给出 $c_{\mathcal B}\le\epsilon D$。在放置因子
输入之后，我们将取 $\epsilon=\tau/(rb)$。下面证明对 $M_0$ 的另一项要求。

\begin{lemma}\label{lem:swaps-support}
对所有 $u,y\in\mathcal B$，若 $M_0[u,y]\ne0$，则
\[
 \bigl|\{s\in[Q]:a_y(s)\ne d_u(s)\}\bigr|\le R-1.
\]
\end{lemma}
\begin{proof}
将 $f_{\mathcal I}(W_0)$ 在可交换矩阵 $\Pi_s$ 中展开。
由于 $\deg f_{\mathcal I}\le R-1$ 且 $\Pi_s^2=I$，所得的每个
单项式都是某个 $\prod_{s\in J}\Pi_s$ 的纯量倍，其中
$J\subseteq[Q]$ 且 $|J|\le R-1$。乘以 $V_0$ 得到
\[
 V_0\prod_{s\in J}\Pi_s=\prod_{s\notin J}\Pi_s.
\]
该置换矩阵的 $(u,y)$ 条目非零要求对每个 $s\notin J$ 都有
$a_y(s)=d_u(s)$。其余位置至多有 $R-1$ 个。若多于 $R-1$
个位置不一致，则每个求和项的 $(u,y)$ 条目都为零，故
$M_0[u,y]=0$。
\end{proof}

\subsection{多项式行}
我们现在用 $M_0$ 构造第~\ref{subsec:incidence-layout} 节中所需的
辅助行 $\vec{P}_u(T)$。
引理~\ref{lem:swaps-matrix} 使 $(I-M_0)\zeta$ 成为
$C^{\mathcal B}\zeta$ 的函数。引理~\ref{lem:swaps-support} 界定了
非零项 $M_0[u,y]\zeta_y$ 用到 $S^{\mathcal B}_{t,d_u(t)}$ 之外坐标的
点有多少个。我们现在用在这些点处为零、在零点处为一的
多项式乘该项。

对 $M_0[u,y]\ne0$，定义
\begin{equation}\label{eq:coefficient-correction}
 E_{u,y}=\{s\in[Q]:a_y(s)\ne d_u(s)\},\qquad
 \varphi_{u,y}(T)=\prod_{s\in E_{u,y}}\left(1-\frac Ts\right).
\end{equation}
空乘积为一。于是 $\deg\varphi_{u,y}\le R-1$，
$\varphi_{u,y}(0)=1$，且当
$a_y(t)\ne d_u(t)$ 时 $\varphi_{u,y}(t)=0$。定义
\begin{equation}\label{eq:swaps-row}
 \vec{P}_u(T)=\delta_u(I-M_0)+
 \sum_{\substack{y\in\mathcal B\\M_0[u,y]\ne0}}
 M_0[u,y]\varphi_{u,y}(T)\delta_y.
\end{equation}
该行作为 $\zeta$ 的线性形式作用，其系数是
单参数 $T$ 的多项式。

\begin{lemma}\label{lem:swaps-forms}
\eqref{eq:swaps-row} 中的行次数至多 $R-1$，且对每个 $u\in\mathcal B$
满足 $\vec{P}_u(0)=\delta_u$。对 $t\in[Q]$ 与
$g\in[r]$，令 $C^{\mathcal B}_{t,g}$ 为
$C^{\mathcal B}$ 在列 $S^{\mathcal B}_{t,g}$ 上的限制。对每个
满足 $g\ne d_u(t)$ 的 $u,t,g$ 以及每个
$\xi\in\mathbb Q^{S^{\mathcal B}_{t,g}}$，其零延拓
$\widetilde\xi$ 到 $\mathcal B$ 满足
\[
 C^{\mathcal B}_{t,g}\xi=\vec{0}
 \quad\Longrightarrow\quad \langle \vec{P}_u(t), \widetilde\xi\rangle=0.
\]
更一般地，$\langle \vec{P}_u(t), \zeta\rangle$ 由
$\zeta|_{S^{\mathcal B}_{t,d_u(t)}}$ 与 $C^{\mathcal B}\zeta$ 确定。
\end{lemma}
\begin{proof}
次数界来自 $\varphi_{u,y}$ 的次数界。在零点处，
\[
 \vec{P}_u(0)=\delta_u(I-M_0)+\delta_uM_0=\delta_u.
\]
将 \eqref{eq:swaps-row} 作用于 $\zeta$ 给出
\begin{equation}\label{eq:swaps-form-decomposition}
 \langle \vec{P}_u(T), \zeta\rangle=((I-M_0)\zeta)_u+
 \sum_{\substack{y\in\mathcal B\\M_0[u,y]\ne0}}
 M_0[u,y]\varphi_{u,y}(T)\zeta_y.
\end{equation}
在 $T=t$ 处，所有满足 $a_y(t)\ne d_u(t)$ 的求和项消失。若
$\widetilde\xi$ 支撑在 $S^{\mathcal B}_{t,g}$ 上且
$g\ne d_u(t)$，则每个剩余求和项在 $\widetilde\xi$ 上为零。
假设 $C^{\mathcal B}_{t,g}\xi=\vec{0}$ 还给出
$C^{\mathcal B}\widetilde\xi=\vec{0}$，故由引理~\ref{lem:swaps-matrix} 得
$(I-M_0)\widetilde\xi=\vec{0}$。
这证明了该蕴含。

对最后一项断言，在
$C^{\mathcal B}$ 下具有相同测量的两个向量相差其核中的一个元素。因此
$(I-M_0)\zeta$ 由 $C^{\mathcal B}\zeta$ 确定。在 $T=t$ 处
\eqref{eq:swaps-form-decomposition} 中的其余项只用到
$\zeta|_{S^{\mathcal B}_{t,d_u(t)}}$。
\end{proof}

固定 $t$ 并记 $\zeta_g=\zeta|_{S^{\mathcal B}_{t,g}}$。
$C^{\mathcal B}$ 的限制满足
\begin{equation}\label{eq:auxiliary-measurements}
 C^{\mathcal B}\zeta=\sum_{g=1}^r C^{\mathcal B}_{t,g}\zeta_g.
\end{equation}
因此固定每个向量 $C^{\mathcal B}_{t,g}\zeta_g$ 就固定了
$C^{\mathcal B}\zeta$，同时对
$r$ 个辅助坐标组分别施加条件。当所有向量
$C^{\mathcal B}_{t,g}\zeta_g$ 都被固定时，$\langle \vec{P}_u(t), \zeta\rangle$ 由
$\zeta|_{S^{\mathcal B}_{t,d_u(t)}}$ 确定。
我们将通过在标号为 $(a_v,d_f)$ 的坐标处放置 $(f,v)$，
在实际布尔变量上得到相应的分离。

\begin{remark}
我们说明 \eqref{eq:swaps-row} 如何把指定组之外的依赖限制为
一个线性测量。取 $Q=R=2$，位于运行时间界所用的区域 $R<Q$
之外，并取 $r=2$。以
$\{0,1\}$ 作为组下标，取 $\mathcal I=\{0,1\}$，于是
$f_{\mathcal I}(z)=1-2z$。从标号 $(00,11)$ 出发，
两次交换生成四个坐标：
\[
 \zeta_0=\zeta_{(00,11)},\quad \zeta_1=\zeta_{(10,01)},\quad
 \zeta_2=\zeta_{(01,10)},\quad \zeta_{12}=\zeta_{(11,00)}.
\]
两次交换都保持这四个坐标的集合，因此我们可以
把计算限制在其上。在该限制上，
$M_0=\Pi_1+\Pi_2-\Pi_1\Pi_2$ 固定
\[
 c(\zeta)=\zeta_0-\zeta_1-\zeta_2+\zeta_{12}
\]
的核。
以 $u=(00,11)$ 为指标的行给出
\[
 \langle \vec{P}_u(T), \zeta\rangle=\zeta_0-\frac T2\zeta_1-T\zeta_2.
\]
它在零点处等于 $\zeta_0$，而
\[
 \langle \vec{P}_u(1), \zeta\rangle=c(\zeta)+\frac12\zeta_1-\zeta_{12},\qquad
 \langle \vec{P}_u(2), \zeta\rangle=c(\zeta)-\zeta_2-\zeta_{12}.
\]
在 $t=1$ 处，$\zeta_1,\zeta_{12}$ 的标号满足
$a_y(1)=1=d_u(1)$。在 $t=2$ 处，
$\zeta_2,\zeta_{12}$ 的标号满足 $a_y(2)=1=d_u(2)$。
因此固定 $c(\zeta)$ 后，$\langle \vec{P}_u(t), \zeta\rangle$ 只依赖于
其指定组中的坐标。这里 $c(\zeta)$ 在四坐标限制上扮演
$C^{\mathcal B}\zeta$ 的角色。
完整构造必须对每个 $u\in\mathcal B$
以至多 $\epsilon D$ 个测量行给出这一性质。
\end{remark}

\subsection{选择变量组与因子组}\label{subsec:choosing-labels}
我们现在从第~\ref{subsec:incidence-layout} 节中选取族 $(a_v)_{v\in[n]}$ 与
$(d_f)_{f\in\mathcal F}$，并置
\[
 u_{f,v}=(a_v,d_f)\qquad((f,v)\in\mathcal A).
\]
这些选择有两项要求。第一，若某个标号出现 $c_u$
次，则其关联需要辅助数组副本中的 $c_u$ 个不同位置。因此
所需的副本数为
\[
 c_u=\#\{(f,v)\in\mathcal A:u_{f,v}=u\},\qquad
 h=\max_{u\in\mathcal B}c_u.
\]
第二，实际的变量组 $\{v:a_v(t)=g\}$ 必须很小。
只界定 $h$ 只能界定一个组中的关联数，
而不能给出对其布尔变量个数的所需界。
下面的引理同时控制这两个量。

\begin{lemma}\label{lem:incidence-labels}
设 $(I_f)_{f\in\mathcal F}$ 满足 \eqref{eq:factor-incidences}，
其中 $n\le M\le bn$，并对整数 $r\ge2$ 与
$Q\ge1$ 置 $D=r^{2Q}$。假设
\begin{equation}\label{eq:incidence-balance-hypotheses}
 2(b+w)D^2\le M,\qquad 2Qr^2\le n.
\end{equation}
则可以确定性地选取 $a_v,d_f\in[r]^Q$，其中对空作用域的因子取 $d_f(t)=1$，
使得
\[
 \begin{aligned}
 \max_{u\in\mathcal B}\#\{(f,v)\in\mathcal A:(a_v,d_f)=u\}
 &\le\frac{2M}{D},\\
 \#\{v\in[n]:a_v(t)=g\}&\le\frac{2n}{r}
 \qquad(t\in[Q],\ g\in[r]).
 \end{aligned}
\]
该构造需要
$\poly(n,M,|\mathcal F|,D,Q,r,b,w)$ 次位运算。
\end{lemma}
\begin{proof}
我们先估计随机标号的效果；随后再确定性地
固定它们的取值。为此估计，所有 $a_v$ 以及
所有 $I_f$ 非空的 $d_f$ 都独立且在 $[r]^Q$ 上均匀。
对空的 $I_f$，我们取 $d_f$ 恒为一。
每个关联标号 $u_{f,v}=(a_v,d_f)$ 在 $\mathcal B$ 上均匀。
置 $n_{t,g}=\#\{v:a_v(t)=g\}$，并考虑
\[
 \Phi=\sum_{u\in\mathcal B}(c_u-M/D)^2,
 \qquad
 \Xi=\sum_{t=1}^Q\sum_{g=1}^r(n_{t,g}-n/r)^2.
\]
它们作为指示函数之和的展开为
\[
 \begin{aligned}
 \Phi&=\sum_{(f,v),(f',v')\in\mathcal A}
       \mathds{1}_{\{(a_v,d_f)=(a_{v'},d_{f'})\}}-\frac{M^2}{D},\\
 \Xi&=\sum_{t=1}^Q\sum_{v,v'\in[n]}
       \mathds{1}_{\{a_v(t)=a_{v'}(t)\}}-\frac{Qn^2}{r}.
 \end{aligned}
\]
两个关联标号除非共享一个因子或一个
变量，否则相互独立。对每个 $(f,v)$，共享
$f$ 或 $v$ 的关联至多有 $b+w$ 个，其中包括 $(f,v)$ 自身。独立对的相等
概率为 $1/D$；其余对的概率至多为一。因此
\[
 \mathbb E\Phi\le M(b+w).
\]
对每个 $t$ 与 $v\ne v'$，值 $a_v(t),a_{v'}(t)$ 独立
且在 $[r]$ 上均匀。因此第二个展开给出
\[
 \mathbb E\Xi=Qn(1-1/r)\le Qn.
\]

为确定性地得到标号，我们逐个固定
$a_v$ 以及 $I_f$ 非空时的 $d_f$。每一步我们选取一个
不增大
\[
 \frac{\Phi}{M(b+w)}+\frac{\Xi}{Qn}
\]
的条件期望的值。
这样的选择存在，因为当前期望是下一个 $a_v$ 或 $d_f$
所有可能取值上的平均。最终值
至多为二。由于两个求和项都非负，
\[
 c_u\le\frac MD+\sqrt{2M(b+w)},\qquad
 n_{t,g}\le\frac nr+\sqrt{2Qn}.
\]
\eqref{eq:incidence-balance-hypotheses} 中的两个不等式使
平方根项分别至多 $M/D$ 与 $n/r$。

下面的枚举在
$\poly(n,M,|\mathcal F|,D,Q,r,b,w)$ 次位运算内计算条件期望。
置 $K=r^Q$，则 $D=K^2$。$\Phi$ 的展开中每个指示函数至多依赖于
四个独立选择
$a_v,d_f,a_{v'},d_{f'}$；重复的下标只贡献一个选择。
$\Xi$ 中每个指示函数至多依赖于 $a_v,a_{v'}$。在任意
部分固定之后，对单个指示函数中剩余选择的至多 $K^4=D^2$ 种赋值
进行枚举。其条件概率的
分母整除 $K^4$。对 $M^2+Qn^2$ 个指示函数求和即可计算
两个期望。计入归一化因子 $M(b+w)$ 与 $Qn$，
这些计算中所有有理数的分子与分母都
具有
\[
 O\!\left(\log(2D)+\log(2M)+\log(2n)+\log(2Q)+\log(2+b+w)\right)
\]
位。至多有 $n+|\mathcal F|$ 个选择需要固定，每个都有
$K$ 个候选值。由于 $K^4=D^2$，整个构造需要
$\poly(n,M,|\mathcal F|,D,Q,r,b,w)$ 次位运算。
\end{proof}

\subsection{在辅助数组中放置变量的副本}\label{subsec:embedding}
现在我们把辅助行 $\vec{P}_u(T)$ 与矩阵
$C^{\mathcal B}$ 转移到定义~\ref{def:layout} 所要求的、作用在原始变量上的
行 $\vec{P}_{f,v}(T)$ 与矩阵 $C_{t,g}$。
假设 \eqref{eq:incidence-balance-hypotheses} 成立并应用
引理~\ref{lem:incidence-labels}。取 $h=\max_u c_u$，用 $[h]\times\mathcal B$ 为
辅助坐标的 $h$ 个副本编号。
对每个 $(f,v)\in\mathcal A$，我们指定 $c(f,v)\in[h]$，使得所有
对 $(c(f,v),u_{f,v})$ 互不相同。这是可行的，因为任何
标号出现不超过 $h$ 次，且
\[
 hD\le2M\le2bn.
\]
定义线性映射
$\iota:\mathbb Q^n\longrightarrow\mathbb Q^{[h]\times\mathcal B}$ 为
\[
 (\iota x)_{c,u}=
 \begin{cases}
 x_v,&(c,u)=(c(f,v),u_{f,v})\text{对某个 }(f,v)\in\mathcal A,\\
 0,&\text{其他情形}.
 \end{cases}
\]
被占用的位置互不相同，因此该定义无歧义。
不同位置可以包含同一个 $x_v$，未使用的位置
固定为零。对 $c\in[h]$，定义投影
\[
 \pi_c:\mathbb Q^{[h]\times\mathcal B}\longrightarrow\mathbb Q^{\mathcal B},
 \qquad (\pi_c\eta)_u=\eta_{c,u}.
\]
现在置
\begin{equation}\label{eq:actual-rows}
 \begin{gathered}
 \vec{P}_{f,v}(T)=\vec{P}_{u_{f,v}}(T)\pi_{c(f,v)}\iota,\\
 S_{t,g}=\{v\in[n]:a_v(t)=g\},\qquad \sigma_f(t)=d_f(t).
 \end{gathered}
\end{equation}
行 $\vec{P}_{f,v}(T)$ 作用在实际变量 $x\in\mathbb Q^n$ 上。
在 $T=0$ 处，它取出位置 $(c(f,v),u_{f,v})$ 中的值，
即 $x_v$。因此 $\vec{P}_{f,v}(0)=\delta_v$，其次数至多
为 $R-1$。引理~\ref{lem:incidence-labels} 给出 $|S_{t,g}|\le2n/r$。

用
\begin{equation}\label{eq:actual-measurements}
 \widetilde C=\operatorname{diag}
   \bigl(\underbrace{C^{\mathcal B},\ldots,C^{\mathcal B}}_{h\text{ 个副本}}\bigr),
 \qquad \widehat C=\widetilde C\iota,
 \qquad C_{t,g}=\widehat C|_{S_{t,g}}.
\end{equation}
定义测量矩阵。
该限制保留 $S_{t,g}$ 中的列以及所有行，包括
零行。于是 $\widehat Cx$ 记录每个副本中
在其坐标由 $x$ 填充后的辅助测量。
矩阵 $C_{t,g}$ 是整值的。其元素不必属于
$\{0,1,-1\}$，因为包含
同一个 $x_v$ 的若干位置上的系数会被相加。$\widetilde C$ 的每一行至多
有 $D$ 个非零元素，且都等于 $1$ 或 $-1$。乘以
$\iota$ 会把同一列中的元素相加并舍弃未使用的位置，
因此它不会增大
其绝对值之和。从而 $C_{t,g}$ 的每一行
的绝对行和至多为 $D$，且其在布尔赋值上的取值
落在 $[-D,D]\cap\mathbb Z$ 中。

若 \eqref{eq:incidence-tail} 成立，则每个 $C_{t,g}$ 有
$c_{t,g}=hc_{\mathcal B}$ 行。因此
\begin{equation}\label{eq:incidence-measurement-count}
 \sum_{g=1}^r c_{t,g}=rh c_{\mathcal B}
 \le rh\epsilon D\le2r\epsilon M\le2r\epsilon bn.
\end{equation}
取 $\epsilon=\tau/(rb)$ 使该值至多为 $2\tau n$。

为验证分离性，固定 $t\in[Q]$、$g\in[r]$ 以及
满足 $g\ne\sigma_f(t)$ 的
$(f,v)\in\mathcal A$。设
$\xi\in\mathbb Q^{S_{t,g}}$ 满足 $C_{t,g}\xi=\vec{0}$，并将其
零延拓为 $\widetilde\xi\in\mathbb Q^n$。等式
\eqref{eq:actual-measurements} 给出
$\widetilde C\iota\widetilde\xi=\vec{0}$。于是对每个 $c\in[h]$，
\[
 C^{\mathcal B}\pi_c\iota\widetilde\xi=\vec{0}.
\]
$\pi_c\iota\widetilde\xi$ 的每个非零坐标都包含一个
满足 $a_{v'}(t)=g$ 的变量 $v'$。其辅助标号的第一个
序列为 $a_{v'}$，故该向量支撑在
$S^{\mathcal B}_{t,g}$ 上。在副本
$c(f,v)$ 中应用引理~\ref{lem:swaps-forms}，其中 $d_{u_{f,v}}(t)=d_f(t)=\sigma_f(t)$，得到
\[
 \langle \vec{P}_{f,v}(t), \widetilde\xi\rangle
 =\langle \vec{P}_{u_{f,v}}(t), \pi_{c(f,v)}\iota\widetilde\xi\rangle=0.
\]
这就是定义~\ref{def:layout} 中的核条件。
因子 $f$ 的所有输入都使用同一指定组 $d_f(t)$，且
求和计数仍恰好遍历 $n$ 个布尔位。
\subsection{系数界与完成后的构造}
\eqref{eq:actual-rows} 中的替换行 $\vec{P}_{f,v}(T)\in\mathbb Q[T]^{1\times n}$
具有有理系数的多项式元素。现在我们来界定存储这些系数所需的位数；
这些界将用于控制在第~\ref{subsec:one-query-cost} 节中的精确算术。
我们首先给出 \eqref{eq:swaps-row} 中定义的辅助行 $\vec{P}_u(T)$ 的公分母与系数界。
通过 \eqref{eq:actual-rows} 从 $\vec{P}_u(T)$ 过渡到 $\vec{P}_{f,v}(T)$ 只会
在公共列中相加系数或丢弃未使用的列，
因此相同的分母以及绝对值之和的界仍然适用。

\begin{lemma}\label{lem:row-coefficients}
设 $d_{\mathrm{den}}=(R-1)!Q!$。则对每个 $u\in\mathcal B$ 有
$d_{\mathrm{den}}\vec{P}_u(T)\in\mathbb Z[T]^{1\times D}$。每一行 $d_{\mathrm{den}}\vec{P}_u(T)$ 中
所有系数的绝对值之和至多为 $2^{O(Q\log Q)}$。
对 $d_{\mathrm{den}}\vec{P}_{f,v}(T)$ 也有相同的界。
矩阵 $M_0,C^{\mathcal B}$ 以及行 $\vec{P}_u(T)$ 可以在 $\poly(D,Q,r)$ 次位运算内
构造出来。
\end{lemma}
\begin{proof}
对 $0\le j<R$，$f_{\mathcal I}(W_0)$ 的相应求和项等于
\[
 (-1)^{A_0}(-2)^j\binom{W_0-A_0I}{j}
 =\frac{(-1)^{A_0+j}}{j!}
   \prod_{a=0}^{j-1}\bigl(2W_0-2(A_0+a)I\bigr).
\]
右端的每个矩阵因子都是整值的，因为
$2W_0=QI-\sum_s\Pi_s$。同时 $j!$ 整除 $(R-1)!$。
因此乘以置换矩阵 $V_0$ 表明
$(R-1)!M_0$ 是整值的。$\varphi_{u,y}$ 的系数的
分母整除 $\prod_{s\in E_{u,y}}s$，而后者整除 $Q!$。
方程~\eqref{eq:swaps-row} 证明了所断言的公分母。

对于系数估计，我们使用矩阵范数
$\|A\|_\infty=\max_i\sum_j|A[i,j]|$
以及多项式范数 $\|\sum_j a_jT^j\|_1=\sum_j|a_j|$。
由于 $\|W_0\|_\infty\le Q$ 且上述每个因子满足 $0\le A_0+a\le Q$，每个 $2W_0-2(A_0+a)I$ 的范数至多为 $4Q$。
由此可得
\[
 \|M_0\|_\infty\le R(4Q)^{R-1},\qquad
 \|\varphi_{u,y}\|_1
 \le\prod_{s\in E_{u,y}}(1+1/s)\le2^Q.
\]
因此，行 $d_{\mathrm{den}}\vec{P}_u$ 中所有系数的绝对值之和至多为
\[
 d_{\mathrm{den}}\bigl(1+(1+2^Q)R(4Q)^{R-1}\bigr).
\]
其对数式为 $O(Q\log Q)$，因为 $R<Q$。
对于 $u=u_{f,v}$，将 $d_{\mathrm{den}}\vec{P}_u(T)$ 乘以
$\pi_{c(f,v)}\iota$ 会在公共列中相加系数并丢弃
未使用的列。由 \eqref{eq:actual-rows}，所得的行就是
$d_{\mathrm{den}}\vec{P}_{f,v}(T)$，因此它满足相同的界。

枚举引理~\ref{lem:swaps-matrix} 中的特征基、计算
$f_{\mathcal I}(W_0)$、以及构造多项式 $\varphi_{u,y}$
需要 $\poly(D,Q,r)$ 次算术运算。使用上述公分母，
在这些矩阵与多项式计算中存储的系数的分子与分母
具有 $O(Q\log Q)$ 位。因此该构造需要 $\poly(D,Q,r)$ 次位运算。
\end{proof}

该构造有三个定量要求。\eqref{eq:incidence-balance-hypotheses} 中的两个
不等式界定了辅助副本的数量以及每一组中实际变量的数量。
取 $\epsilon=\tau/(rb)$ 的尾部界通过 \eqref{eq:incidence-measurement-count} 给出总的测量
界。我们汇总该结果以供计数算法使用。

\begin{lemma}\label{lem:layout}
设 $n,b,w\ge1$ 为整数，并设 $(I_f)_{f\in\mathcal F}$ 是 $[n]$ 的大小至多为 $w$ 的子集，
且每个 $v\in[n]$ 至多属于 $b$ 个集合。令
$M=\sum_f|I_f|$ 并假设 $n\le M\le bn$。设 $r\ge2$ 且
$1\le R<Q$ 为整数，设 $0<\tau\le1/2$，并令 $D=r^{2Q}$。
假设
\[
 2(b+w)D^2\le M,\qquad 2Qr^2\le n,
\]
且区间
$\mathcal I=\{A_0,\ldots,A_0+R-1\}\subseteq\{0,\ldots,Q\}$ 满足
\[
 \Pr[W\notin\mathcal I]\le\frac{\tau}{rb},\qquad
 W\sim\operatorname{Bin}\!\left(Q,\frac{r-1}{2r}\right).
\]
则可以确定性地在
$\poly(n,M,|\mathcal F|,D,Q,r,b,w)$ 次位运算内构造出
定义~\ref{def:layout} 中的所有数据。
\end{lemma}
\begin{proof}
引理~\ref{lem:swaps-matrix} 构造出满足
$c_{\mathcal B}\le\tau D/(rb)$ 的 $C^{\mathcal B}$。引理~\ref{lem:swaps-support}
与~\ref{lem:swaps-forms} 给出辅助多项式行。
引理~\ref{lem:incidence-labels} 选取满足 $hD\le2M$ 与 $|S_{t,g}|\le2n/r$ 的变量组与因子组。
然后 \eqref{eq:actual-rows} 与 \eqref{eq:actual-measurements} 中的映射
给出满足 $\vec{P}_{f,v}(0)=\delta_v$、次数至多 $R-1\le R$
以及所需核蕴含关系的行。$C_{t,g}$ 的每一行的
绝对值行和至多为 $D$，且 \eqref{eq:incidence-measurement-count}
对每个 $t\in[Q]$ 给出 $\sum_{g=1}^r c_{t,g}\le2\tau n$。
引理~\ref{lem:row-coefficients} 提供公分母与
系数界。这些构造合起来需要
$\poly(n,M,|\mathcal F|,D,Q,r,b,w)$ 次位运算，并且所得
数组占用 $\poly(n,M,|\mathcal F|,D,Q,r,b,w)$ 位。
\end{proof}

\subsection{选择参数}\label{subsec:incidence-parameters}
在应用中，最大作用域 $w$ 与包含某个变量的因子的最大数量
$b$ 都将被 $\log n$ 的固定幂所界定。我们现在在该情形下选择满足
引理~\ref{lem:layout} 的参数。计数
代价的指数取决于三个量
\[
 \frac RQ\log\frac{2Q}{R},\qquad
 \tau\log(2D+1),\qquad \frac1r.
\]
它们分别对应于恢复泰勒系数、
枚举测量值、以及枚举一组内的赋值。以下选择同时界定这三者。

\begin{lemma}\label{lem:incidence-parameters}
固定整数 $\beta\ge1$。设 $n\ge20$，令
$\lambda=\lceil\log_2 n\rceil$ 且 $\nu=\lceil\log_2\lambda\rceil$，
并假设 $b,w\le\lambda^\beta$ 且 $n\le M\le bn$。
对于所有充分大的 $n$，参数
\begin{equation}\label{eq:parameter-QR}
 r=\left\lfloor\frac{\sqrt\lambda}{\nu^2}\right\rfloor,
 \qquad Q=\left\lfloor\frac{\lambda}{128\nu}\right\rfloor,
 \qquad R=(\beta+3)\lceil\sqrt\lambda\rceil,
\end{equation}
\[
 \tau=\frac1{r\lambda},\qquad
 \epsilon=\frac1{br^2\lambda},\qquad D=r^{2Q},
\]
连同证明中指定的区间，满足引理~\ref{lem:layout} 的所有
数值假设。它们还满足
\[
 R\ge32,\qquad 6R\le Q,\qquad D\le2^{\lambda/128},
\]
以及
\begin{equation}\label{eq:balanced-cost}
 \frac RQ\log\frac{2Q}{R}+\tau\log(2D+1)+\frac1r
 =O_\beta\!\left(\frac{\nu^2}{\sqrt\lambda}\right).
\end{equation}
给定数值 $n,M,b,w$，这些参数与数值
假设可以在 $\poly_\beta(\log n)$ 次位运算内精确计算并检验。对于固定的 $\beta$，只有有限
多个 $n$ 无法通过这些检验。
\end{lemma}
\begin{proof}
对于充分大的 $n$，这些参数具有阶
\[
 r=\Theta(\sqrt\lambda/\nu^2),\qquad
 Q=\Theta(\lambda/\nu),\qquad R=\Theta_\beta(\sqrt\lambda).
\]
特别地，最终有 $r\ge2$、$R\ge32$ 以及 $6R\le Q$。
由于 $\log_2 r\le\nu/2$，
\[
 \log_2 D=2Q\log_2 r
 \le2\frac{\lambda}{128\nu}\frac\nu2=\frac{\lambda}{128}.
\]
对充分大的 $n$，两个平衡不等式成立：
\[
 2(b+w)D^2\le4\lambda^\beta2^{\lambda/64}\le n\le M,
 \qquad 2Qr^2\le n.
\]
对于第二个不等式，$Q,r$ 被 $\lambda$ 的固定幂所界定，
而 $n>2^{\lambda-1}$。第一个不等式由同样的比较可得。

令
\[
 \mu=\frac{Q(r-1)}{2r},\qquad
 A_0=\left\lfloor\mu-\frac{R-1}{2}\right\rfloor,
 \qquad \mathcal I=\{A_0,\ldots,A_0+R-1\}.
\]
这里 $Q/4\le\mu\le Q/2$。由于 $R\le Q/6$，该区间位于
$\{0,\ldots,Q\}$ 中。其中心与 $\mu$ 的差小于一，
且 $R\ge32$ 保证它包含 $\mu$ 的 $R/4$ 邻域内的每个整数。
对于 $W\sim\operatorname{Bin}(Q,(r-1)/(2r))$，Hoeffding 不等式给出
\[
 \Pr[W\notin\mathcal I]
 \le2e^{-R^2/(8Q)}
 \le2e^{-16(\beta+3)^2\nu}
 \le2^{-(\beta+2)\nu}
 \le\lambda^{-(\beta+2)}
 \le\frac1{br^2\lambda}.
\]
最后一个不等式使用了 $b\le\lambda^\beta$ 与 $r^2\le\lambda$。
这证明了带有 $\epsilon=\tau/(rb)$ 的 \eqref{eq:incidence-tail}。同时，一旦 $r\ge2$ 就有 $0<\tau\le1/2$。

最后，
\[
 \frac RQ=O_\beta(\nu/\sqrt\lambda),\qquad
 \log(2Q/R)=O(\nu),\qquad
 \tau\log(2D+1)=O(1/r).
\]
连同 $1/r=O(\nu^2/\sqrt\lambda)$，这些证明了
\eqref{eq:balanced-cost}。

我们最后说明 \eqref{eq:parameter-QR} 中的参数
可以在所述的位界内精确构造出来。
算法将 $\lambda=\lceil\log_2 n\rceil$ 求为满足 $n\le2^\lambda$ 的最小
整数，然后以同样的方式求出
$\nu=\lceil\log_2\lambda\rceil$。它取满足 $r^2\nu^4\le\lambda$ 的最大整数
作为 $r$，通过整数除法计算 $Q$，
并通过整数平方根搜索求出 $\lceil\sqrt\lambda\rceil$ 以得到 $R$。
如果 $r<2$ 或 $Q\le R$，算法在此阶段报告失败。
否则它使用整数或有理算术构造 $\tau,\epsilon,D$ 以及
$\mathcal I$ 的端点。随后它检验
$\mathcal I\subseteq\{0,\ldots,Q\}$、两个平衡不等式、
$R\ge32$ 以及 $6R\le Q$。\eqref{eq:incidence-tail} 中的尾部概率
的分母为 $(2r)^Q$，因此其所需的
界也可以用整数算术检验。这些计算
需要 $\poly_\beta(\log n)$ 次位运算。前面的估计
表明对于充分大的 $n$，所有检验都通过。
\end{proof}

\section{计数稀疏多项式方程的公共零点}\label{sec:setup}
我们现在证明定理~\ref{thm:overview-sparse-polynomials}。输入
是 $\mathbb F_2$ 上 $n$ 个布尔变量中的 $m_{\mathrm{test}}\le cn$ 个多项式方程
$P_i(x)=0$ 构成的系统 $\mathcal P$。每个 $P_i$ 是
次数至多 $k$ 的至多 $(\log n)^c$ 个单项式的列表，其中 $k,c$ 固定。我们的目标是整数
\begin{equation}\label{eq:common-zero-count}
 N(\mathcal P)=\#\{x\in\{0,1\}^n:P_i(x)=0\text{ 对每个 }i\}.
\end{equation}
首先，我们为每个方程构造一个因子，使其在布尔立方体上的取值
记录公共零点指示函数。困难在于
使其总输入次数保持为 $n$ 的线性函数。接着我们为这些因子构造
替换形式，并使用第~\ref{sec:counting} 节
对乘积求和。第~\ref{sec:hashes} 节将所得的
算法用于任意给定的系统，包括那些仍然
接受某个公式的非满足赋值的系统。

\subsection{单个方程的整数多项式}\label{subsec:rational-lift}
在布尔输入上，我们可以将每个正幂 $x_v^j$ 替换为 $x_v$。
将重复的单项式按模二合并。如果 $\mathcal M_i$ 是
$P_i$ 所得的变量集列表，则定义其整数提升为
\[
 L_i(\vec X)=\sum_{M\in\mathcal M_i}\prod_{v\in M}X_v .
\]
空单项式为 $1$，空和为 $0$。值 $L_i(x)$
统计等于一的单项式个数，因此 $P_i(x)=0$ 当且仅当
$L_i(x)$ 为偶数。记 $I_i=\bigcup_{M\in\mathcal M_i}M$。尽管
$L_i$ 可以读取 $k(\log n)^c$ 个变量，它仍然具有至多 $k$ 的次数。

令 $\ell=n+1$。由于 $0\le N(\mathcal P)\le2^n<2^\ell$，只需
求出它模 $2^\ell$ 的余数。我们首先用模该二的幂的整数多项式
表示偶数性指示函数。
直接相乘 $m_{\mathrm{test}}$ 个这样的指示函数只会得到
二次的次数界；下一个引理的形式将让我们
统一控制所有方程上的次数。

令 $h(Y)=Y(Y-1)$。它在每个整数 $Y$ 处的值都是偶数。
将指示函数按 $h$ 的幂展开，会为
乘积中的每个 $h$ 副本显示出一个因子二。

\begin{lemma}\label{lem:rational-idempotent}
对每个整数 $\ell\ge1$，都可以构造多项式
$e_0,\ldots,e_{\ell-1}\in\mathbb Z[Y]$，每个次数至多为一，
使得
\begin{equation}\label{eq:idempotent-expansion}
 p_\ell(Y)=\sum_{j=0}^{\ell-1}e_j(Y)h(Y)^j,
 \qquad
 p_\ell(y)\equiv\mathds{1}_{\{2\mid y\}}\pmod{2^\ell}
 \quad(y\in\mathbb Z).
\end{equation}
$e_j$ 的系数具有 $\poly(\ell)$ 位，且
该构造使用 $\poly(\ell)$ 次位运算与 $\poly(\ell)$
位空间。
\end{lemma}
\begin{proof}
设 $F_\ell$ 为 Beigel 与
Tarui 的模放大多项式~\cite[Section~2.3]{BT94}。Chan 与 Williams 将 $F_\ell$ 应用于
子句违反指示函数之和，以确定性地计数满足赋值~\cite[Theorem~2.5 and proof of Theorem~1.6]{CW21}。
我们需要偶奇偶性的指示函数，因此取其补：
\begin{equation}\label{eq:binary-indicator}
 p_\ell(Y)=1-F_\ell(Y)=(1-Y)^\ell
       \sum_{j=0}^{\ell-1}\binom{\ell+j-1}{j}Y^j.
\end{equation}
该有限和是 $(1-Y)^\ell$ 模 $Y^\ell$ 的逆。
事实上，$(1+Y+Y^2+\cdots)^\ell$ 中 $Y^j$ 的系数
统计将 $j$ 分成 $\ell$ 部分的弱复合数，等于
$\binom{\ell+j-1}{j}$。因此
\[
 p_\ell(Y)\equiv1\pmod{Y^\ell},
 \qquad
 p_\ell(Y)\equiv0\pmod{(1-Y)^\ell}.
\]
对于偶数 $y$，第一个同余给出 $p_\ell(y)\equiv1\pmod{2^\ell}$。
对于奇数 $y$，第二个给出 $p_\ell(y)\equiv0\pmod{2^\ell}$。

多项式 $p_\ell$ 具有整数系数且次数小于
$2\ell$。通过首一二次式
$h=Y^2-Y$ 连续作除法，可将其唯一地写成
$\sum_{j<\ell}e_jh^j$，其中整数余式 $e_j$ 的次数
小于二。附录~\ref{app:arithmetic} 中的引理~\ref{lem:rational-lift-height}
界定了它们的系数，并证明了该构造的 $\poly(\ell)$ 次位运算界。
\end{proof}

\subsection{保持乘积次数为线性}
完整乘积 $\prod_{i=1}^{m_{\mathrm{test}}} p_\ell(L_i(\vec{X}))$ 的次数至多为
$km_{\mathrm{test}}(2\ell-1)$。每个因子 $h(y)$ 在
整数参数处贡献一个因子二，因此含有至少 $\ell$ 个此类因子的项
可以在模 $2^\ell$ 下丢弃。这让我们可以只保留
$h$ 的总幂次小于 $\ell$ 的项，且是在所有方程上统一计算。
引入形式变量 $V$ 来记录该幂次，对每个方程使用同一个
$V$。其指数记录在整数参数处整除相应项的
一个可保证的二的幂。它不是需要求和的
额外布尔变量。
这里我们选择第~\ref{subsec:product-interface} 节中的系数代数：
\[
 \mathcal A_\ell=\mathbb Q[V]/(V^\ell).
\]
在该代数中，每次乘法都会丢弃满足 $s\ge\ell$ 的幂 $V^s$。
形成乘积后，我们使用线性映射
\[
 \mathcal S_\ell\!\left(\sum_{s=0}^{\ell-1}b_sV^s\right)
 =\sum_{s=0}^{\ell-1}b_s .
\]
来相加其保留的系数。该映射不必保持乘积，因此只在
因子相乘之后应用。我们对多项式
系数 $b_s$ 使用相同的记号。

\begin{lemma}\label{lem:total-order}
设 $m_{\mathrm{test}}\ge0$ 与 $\ell\ge1$ 为整数。对整数
$y_1,\ldots,y_{m_{\mathrm{test}}}$，令
\[
 B_i(V)=\sum_{j=0}^{\ell-1}e_j(y_i)h(y_i)^jV^j.
\]
则
\[
 \mathcal S_\ell\!\Bigl(\prod_{i=1}^{m_{\mathrm{test}}} B_i(V)\bmod V^\ell\Bigr)
 \equiv\prod_{i=1}^{m_{\mathrm{test}}}\mathds{1}_{\{2\mid y_i\}}\pmod{2^\ell}.
\]
对任意非负整数 $d_*$，如果 $y_i$ 改为
$\mathbb Q$ 上总次数至多 $d_*$ 的多项式，
那么 $V^s$ 的每个保留系数以及它们的和，
其总输入次数至多为
$d_*(m_{\mathrm{test}}+2(\ell-1))$。
\end{lemma}
\begin{proof}
在截断之前，贡献到 $V^s$ 系数的项具有形式
\[
 \prod_{i=1}^{m_{\mathrm{test}}} e_{j_i}(y_i)h(y_i)^{j_i},
 \qquad \sum_{i=1}^{m_{\mathrm{test}}} j_i=s.
\]
在整数参数处它可被 $2^s$ 整除。因此丢弃所有
满足 $s\ge\ell$ 的系数后，系数之和
在模 $2^\ell$ 下保持不变。若不截断，该和就是
$\prod_i p_\ell(y_i)$，由引理~\ref{lem:rational-idempotent}
它具有所需的同余。

对于多项式参数，每个保留项都有
$\sum_i j_i\le\ell-1$。由于 $\deg e_{j_i}\le1$ 且 $\deg h=2$，
其次数至多为
\[
 \sum_{i=1}^{m_{\mathrm{test}}} d_*(1+2j_i)
 \le d_*(m_{\mathrm{test}}+2(\ell-1)).
\]
该界对每个系数以及它们的和都成立。
\end{proof}

为了理解截断为何使用所有方程上的总阶数，
取 $\ell=2$。则
\[
 p_2(Y)=1-3Y^2+2Y^3=(1-Y)+(2Y-1)h(Y).
\]
对于两个方程，模 $V^2$ 相乘后再应用
$\mathcal S_2$，保留
\[
 (1-y_1)(1-y_2)
 +(2y_1-1)h(y_1)(1-y_2)
 +(1-y_1)(2y_2-1)h(y_2).
\]
其次数至多为四。被省略的项
$(2y_1-1)(2y_2-1)h(y_1)h(y_2)$ 次数为六，但在整数输入处可被
四整除。

\subsection{使用给定的替换形式计数}\label{subsec:count-polynomial}
对于方程 $i$，向第~\ref{sec:counting} 节提供因子
\begin{equation}\label{eq:lifted-factor}
 \mathcal B_i((X_v)_{v\in I_i},V)
 =\sum_{j=0}^{\ell-1}e_j(L_i(\vec X))h(L_i(\vec X))^j V^j
 \quad\text{in }\mathcal A_\ell[(X_v)_{v\in I_i}].
\end{equation}
在一个因子内，$X_v$ 的每次出现都使用相同的替换
形式；不同的因子可以为 $X_v$ 使用不同的形式。即使
每个因子获得各自独立的输入，$L_i$ 在这些输入中的次数
也至多为 $k$。引理~\ref{lem:total-order} 将乘积 \emph{在 $\mathcal A_\ell$ 中} 的输入
次数界定为
$d=k(m_{\mathrm{test}}+2(\ell-1))$。因此该乘积恰好满足
公共例程的次数假设 \eqref{eq:general-product-degree}。
添加作用域为 $\{v\}$ 且等于一的因子可确保
所有声明的变量都出现在构造中。

设 $H(T)\in\mathcal A_\ell[T]$ 为这些因子的 \eqref{eq:general-count-polynomial}。
由于所有替换在零处都恢复其布尔输入，其在 $T=0$ 处的 $V$-系数为整数。
由同一引理，对其值求和得到
\begin{equation}\label{eq:origin-residue}
 \deg_T H\le d_H:=kR(m_{\mathrm{test}}+2(\ell-1)),\qquad
 \mathcal S_\ell(H(0))\equiv N(\mathcal P)\pmod{2^\ell}.
\end{equation}
因此 $\mathcal S_\ell(H(0))$ 的最小非负剩余
就是精确计数。插值在有理数上进行；
该同余仅在 $T=0$ 处的布尔输入上需要。

\subsection{计算因子的代价}\label{subsec:one-query-cost}
\begin{proposition}\label{prop:count}
给定 $n$ 个变量上的次数至多 $k$ 的多项式方程列表
以及为其因子满足定义~\ref{def:layout} 的设计。
设 $S\ge1$ 界定单项式列表与
因子作用域的总长度，并使用任意满足 $2^\ell>2^n$ 的 $\ell$。取
\eqref{eq:origin-residue} 中的 $d_H$ 以及 \eqref{eq:hasse-order} 中的 $q,J$，
则 $N(\mathcal P)$ 可以在
\[
 \poly(n,m_{\mathrm{test}},\ell,D,Q,R,J,r,k,S)\,
 C_q(2D+1)^{2\tau n}2^{2n/r}
\]
次位运算以及 $\poly(n,m_{\mathrm{test}},\ell,D,Q,R,J,r,k,S)$
位工作空间内确定性地计算出来。
\end{proposition}
\begin{proof}
以 $\mathcal A_\ell$ 中的系数应用命题~\ref{prop:general-count}。在求值点 $t$ 处，该例程使用
$\mathcal A_\ell[z]/(z^J)$，因此一个存储元素只有 $J\ell$
个有理系数。为了在 \eqref{eq:jet-substitution} 中于固定的布尔
替换 $\vec X=\mathbf1_U$ 处计算因子 $i$，首先
通过相乘其至多 $k$ 个被替换的输入来计算每个列出的单项式，
然后将结果相加得到 $L_i$。将此值用于
\eqref{eq:lifted-factor}，通过连续乘以 $h(L_i)$
得到其各次幂。这需要显示参数中的多项式次有理运算，
并且从不枚举整个作用域 $I_i$ 的赋值。多项式 $e_j$ 由
引理~\ref{lem:rational-idempotent} 一次性计算。

对于位复杂度，各行具有公分母
$d_{\mathrm{den}}=(R-1)!Q!$。每个被替换的单项式的分母
整除 $\Delta=d_{\mathrm{den}}^k$，每个提升因子的分母
整除 $\Gamma=\Delta^{2\ell-1}$，而它们的乘积的分母
整除 $\Omega=\Gamma^{m_{\mathrm{test}}}$。因此指数多个
求和项可以累积在一个公分母上。
附录~\ref{app:arithmetic} 将这些和中的位长界定为
\begin{equation}\label{eq:sum-bit-bound}
 B_{\mathrm{sum}}=\poly(n,m_{\mathrm{test}},\ell,k,Q,\log(D+2),\log(S+2)).
\end{equation}
对每个 $0\le s<\ell$，我们从 $QJ$ 个泰勒系数
插值有理多项式 $[V^s]H(T)$。附录验证了该插值以及 \eqref{eq:origin-residue} 中
最终整数余式的多项式位代价。$O(QJ\ell)$ 个存储
系数、该设计以及命题~\ref{prop:general-count} 中的顺序枚举
给出了所述的空间界。
\end{proof}
\subsection{构造形式并完成证明}\label{subsec:sparse-queries}
每个稀疏多项式只读取少量变量，但一个变量可能出现在许多方程中。我们首先固定出现次数最多的少数变量的取值。在每个由此得到的分支上，我们可以构造命题~\ref{prop:count} 所需的设计。

\begin{proof}[定理~\ref{thm:overview-sparse-polynomials} 的证明]
设 $\mathcal P$ 为给定的多项式系统。空的方程列表给出 $2^n$ 个公共零点；常数方程 $1=0$ 不给出任何公共零点。丢弃方程 $0=0$。对每个声明的变量 $v$，加入一个作用域为 $\{v\}$、等于 1 的中性因子，并记
\[
 M=\sum_{f\in\mathcal F}|I_f|
\]
为因子—变量关联的总数。输入界蕴含 $M=n(\log n)^{O_{k,c}(1)}$，且每个作用域有 $(\log n)^{O_{k,c}(1)}$ 个变量。然而，某个特定变量可能出现在许多这样的作用域中。

令 $\lambda=\lceil\log_2 n\rceil$ 且
\[
 b_0=\left\lceil\frac Mn\right\rceil\lceil\sqrt\lambda\rceil.
\]
设 $S\subseteq[n]$ 为属于超过 $b_0$ 个作用域的变量。对每个 $b\in\{0,1\}^S$，我们固定 $x|_S=b$ 并计数 $[n]\setminus S$ 上的补全；我们称这个受限问题为一个分支。由于 $|S|\le M/b_0\le n/\sqrt\lambda$，分支至多有 $2^{n/\sqrt\lambda}$ 个。
每个分支有相同数目的自由变量 $n'$，且
\[
 n'\ge n-\frac n{\sqrt\lambda}\ge\frac n2
 \qquad(n\ge20).
\]
在 $\mathbb F_2$ 上把固定比特代入每个 $P_i$，并在构造其整数提升之前将重复单项式按模二合并。次数和单项式个数都不会增加。常数方程仍可立即处理。这些分支划分了布尔立方体，因此它们的公共零点计数之和为 $N(\mathcal P)$。

为应用第~\ref{sec:construction} 节，从作用域中移除固定变量，并用 $[n']$ 重新标记自由变量。记 $I'_f\subseteq[n']$ 为受限作用域，$M'=\sum_f|I'_f|$ 为其总关联数。设 $b$ 为一个自由变量的最大因子关联数，$w=\max_f|I'_f|$。则 $b\le b_0$，且 $b$ 和 $w$ 都由 $\log n'$ 的某个固定幂所界定。剩余的中性因子保证 $M'\ge n'$。令
\[
 \lambda'=\lceil\log_2 n'\rceil,\qquad
 \nu'=\lceil\log_2\lambda'\rceil,
\]
并选择满足 $b,w\le(\lambda')^\beta$ 的最小整数 $\beta\ge1$。对固定的 $k,c$，对所有充分大的 $n$，$\beta$ 的可能取值由一个常数界定。我们对 $[n']$ 上的受限作用域 $(I'_f)$ 应用引理~\ref{lem:incidence-parameters}，关联数为 $M'$，使用
\[
 r=\left\lfloor\frac{\sqrt{\lambda'}}{(\nu')^2}\right\rfloor,
 \qquad Q=\left\lfloor\frac{\lambda'}{128\nu'}\right\rfloor,
 \qquad R=(\beta+3)\lceil\sqrt{\lambda'}\rceil.
\]
若 $r<2$ 或 $Q\le R$，则在构造 $\tau$、$\epsilon$ 或奇偶插值区间 $\mathcal I$ 之前，使用下面的穷举搜索回退方案。
否则令
\[
 \tau=\frac1{r\lambda'},\qquad
 \epsilon=\frac1{br^2\lambda'},\qquad D=r^{2Q},
\]
并执行该引理中剩余的可行性检查。对所有充分大的 $n$，引理~\ref{lem:layout} 提供所需的行和测量矩阵，其中
\[
 r=\Theta\!\left(\frac{\sqrt{\log n'}}{(\log\log n')^2}\right),
 \qquad Q=\Theta\!\left(\frac{\log n'}{\log\log n'}\right),
 \qquad R=O_{k,c}(\sqrt{\log n'}),
 \qquad D\le(2n')^{1/128}.
\]

在每个分支上，我们可以使用原始精度 $\ell=n+1$，因为即使不受限的计数也小于 $2^\ell$。由于 $n'\ge n/2$，$\ell$ 和 $m_{\mathrm{test}}$ 都是 $O_c(n')$。因此 $d_H=O_{k,c}(n'R)$、$J=O_{k,c}(n'R/Q)+1$，并且 \eqref{eq:three-costs} 在维数 $n'$ 下的指数为
\[
 O_{k,c}\!\left(\frac{n'(\log\log n')^2}{\sqrt{\log n'}}\right)
 =O_{k,c}\!\left(\frac{n(\log\log n)^2}{\sqrt{\log n}}\right).
\]
所有替换数据、系数数组和插值系统的总表示长度为 $n^{O_{k,c}(1)}$ 位。将每个分支的运行时间乘以至多 $2^{n/\sqrt\lambda}$ 不改变所显示的指数界，因此对各分支答案求和即得所需的运行时间。

为使构造对每个输入规模都可实现，算法使用整数运算检查其数值假设。除上述预备保护条件外，它还检查来自引理~\ref{lem:incidence-labels} 的 $2(b+w)D^2\le M'$ 和 $2Qr^2\le n'$、$\mathcal I$ 的端点、$R\ge32$、$6R\le Q$，以及 \eqref{eq:incidence-tail} 中的二项式尾界，后者可以精确检验。若任一检查失败，算法放弃分支计算，并直接在全部 $2^n$ 个赋值上检查原始方程。对固定的 $k,c$，刚刚证明的界和引理~\ref{lem:incidence-parameters} 表明失败只能对有限多个 $n$ 发生。因此该回退方案被吸收进定理的常数中。

为在所述空间界内合并各分支计数，我们一次处理一个分支。初始化整数累加器 $Z=0$，并将每个分支的答案加到其上。构造和命题~\ref{prop:count} 在每个分支上使用 $n^{O_{k,c}(1)}$ 位空间。
不相交性保证整个过程 $0\le Z\le2^n$，最终值为 $N(\mathcal P)$，如所断言。
\end{proof}

\section{将 \texorpdfstring{$\#k$-SAT}{\#k-SAT} 归约为计数公共零点}\label{sec:hashes}\label{subsec:sparse-hashes}
给定一个至多 $Bn$ 个子句的稀疏 $k$-CNF $F$，其中 $B$ 关于 $n$ 是多重对数级的，我们在其原来的 $n$ 个变量上构造 $O(n)$ 个稀疏多项式方程，使其公共零点恰好为其满足赋值。每个子句一个方程可能需要 $Bn$ 个方程。相反，每个方程将检验一个短子句列表的违反情况的奇偶性。我们同时使用定理~\ref{thm:overview-sparse-polynomials} 来选择这些方程并得到最终计数。

设 $m$ 为被索引子句的数目，$v_e(x)$ 为子句 $e$ 的违反位。作为 $\mathbb F_2$ 上的多项式，$v_e$ 是正文字对应的因子 $1+x_v$ 与负文字对应的因子 $x_v$ 的乘积。恒真子句有 $v_e=0$，空子句有 $v_e=1$。去除重复文字后，$\deg v_e\le k$，且 $v_e$ 至多展开为 $2^k$ 个单项式。令 $v(x)=(v_1(x),\ldots,v_m(x))$。
对于具有 $m$ 列的二元矩阵 $A$，第 $i$ 行指定方程
\begin{equation}\label{eq:hash-query}
 P_i(x):=\sum_{e:A_{ie}=1} v_e(x)=0\qquad\text{over }\mathbb F_2.
\end{equation}
将 $\mathcal P_A$ 记为该方程组，$N(A)=N(\mathcal P_A)$ 记为其公共零点数。每个满足赋值都通过每一行。
因此 $N(A)=\#F+D(A)$，其中 $D(A)$ 计数仍然通过的未满足赋值。我们不知道 $\#F$，但它对每个候选的下一行都相同：使 $N$ 最小的行也使 $D$ 最小。

我们用一个对 $(e,c)$ 的列表表示检验行，其中 $e\in[m]$、$c\in\{0,1\}$，若子句索引重复则系数按模 $2$ 求和。该列表的长度计数所有对，包括重复索引和系数为零的对。因此长度为 $s$ 的行在消去前至多读取 $ks$ 个变量。

\subsection{排除小违反集的初始检验}

若一个赋值只违反许多子句中的一个，则抽样短子句列表的检验不太可能注意到它。因此我们首先构造一个矩阵 $A^{(0)}$，它拒绝每个具有小支撑的非零违反向量。每个通过这些初始检验的未满足赋值必定违反许多子句。随后抽样短子句列表的进一步检验便能以常数概率拒绝每个这样的赋值。

我们使用标准的基于边的 Tanner 码：比特被分配给一个扩展图的边，每个顶点关联的比特必须构成一个规定的局部线性码的码字
（局部约束见
Tanner~\cite[第~II 节]{Tanner81}，扩展图构造见
Sipser 和 Spielman~\cite[第~IV 和~VI 节]{SS96}）。
通常的局部距离和扩展论证给出非零码字权重所需的下界。下面的证明确定性地构造图和局部码，其参数满足我们对检验数目和长度的界。

\begin{lemma}\label{lem:sparse-checks}
设 $n,m,B$ 为正整数且 $m\le Bn$，令
\[
 \mu=\lfloor\log_2 m\rfloor+1,\qquad
 P=20000B^2\mu^2,\qquad D_0=320000B^2\mu^3.
\]
存在一个具有 $m$ 列、至多 $3n$ 行的二元矩阵 $A^{(0)}$ 的确定性构造，使得每行至多有 $D_0$ 个非零项，且
\[
 z\in\{0,1\}^m\setminus\{\vec{0}\},\quad
 A^{(0)}z=\vec{0}\pmod2
 \quad\Longrightarrow\quad
 |\operatorname{supp}z|\ge m/P.
\]
该构造使用 $\exp(\poly(B,\log(m+2)))$ 次位运算和 $\poly(n,m,B)$ 位工作空间。
\end{lemma}

\begin{proof}
我们构造一个边数介于 $m$ 和 $2m$ 之间的二部图，用其边作为临时二元坐标。每个顶点处的检验将迫使被非零边触及的每个顶点有许多非零关联边；随后一个混合界将迫使总权重很大。最后，我们保留 $m$ 个边坐标并将其余的置零。令
\[
 d_G=\min(m,D_0),
 \qquad
 q_G=\min\{2^a:a\in\mathbb Z_{\ge0},\ 2^a d_G\ge m\}.
\]
则 $m\le q_Gd_G<2m$。写 $q_G=2^{\ell_G}$，并以 $\{0,1\}^{\ell_G}$ 的两个副本作为顶点集。对每个左顶点 $u$ 和每个端口 $j\in[d_G]$，放置一条从 $u$ 到 $u\oplus s_j$ 的边，其中 $\oplus$ 表示按位模 $2$ 加法，移位 $s_1,\ldots,s_{d_G}\in\{0,1\}^{\ell_G}$ 将在下面选择。
边由对 $(u,j)$ 索引；平行边保持不同。两侧的每个顶点在每个端口 $j$ 处都有一条关联边，因此该图是 $d_G$-正则的。

为得到混合估计 \eqref{eq:hash-mixing}，我们选择移位使得非空坐标集合的每种奇偶性都接近平衡。对于非零向量
$b\in\{0,1\}^{\ell_G}$，比特 $b\cdot s_j\bmod2$ 是 $s_j$ 中那些 $b$ 等于 $1$ 的坐标的奇偶性。取 $\varepsilon=1/(200B\mu)$，我们要求
\begin{equation}\label{eq:hash-walsh}
 \Bigl|\sum_{j=1}^{d_G}(-1)^{b\cdot s_j}\Bigr|
 \le\varepsilon d_G
 \qquad
 \bigl(b\in\{0,1\}^{\ell_G}\setminus\{\vec{0}\}\bigr).
\end{equation}
\eqref{eq:hash-walsh} 中的和是偶奇偶性的下标 $j$ 的数目减去奇奇偶性的数目。因此，对每个非零 $b$，每种奇偶性在下标中出现的比例介于 $(1-\varepsilon)/2$ 和 $(1+\varepsilon)/2$ 之间。具有此性质的移位列表称为 $\varepsilon$-偏置的~\cite[定义~2.1--2.2]{NN93}。

若 $q_G=1$，则没有条件需要检查。否则 $d_G=D_0$。为表明合适的移位存在，我们独立且均匀地随机选择它们。对每个固定的非零 $b$，符号 $(-1)^{b\cdot s_j}$ 独立且均匀地取 $\{-1,1\}$ 中的值。Hoeffding 不等式将此 $b$ 违反 \eqref{eq:hash-walsh} 的概率界为 $2\exp(-\varepsilon^2d_G/2)$。对 $b$ 取联合界得到总失败概率至多
\[
 2q_G\exp(-\varepsilon^2d_G/2)
 =2q_G\exp(-4\mu)<1,
\]
其中 $q_G<2m<2^{\mu+1}$。因此穷举枚举移位元组并检查 \eqref{eq:hash-walsh} 可在 $\exp(O(d_G\log(m+2)))$ 时间内找到一个有效元组。

接下来我们从 \eqref{eq:hash-walsh} 推出混合估计。设 $\mathsf A_G$ 为 $q_G\times q_G$ 矩阵，其中 $\mathsf A_G[u,v]=|\{j\in[d_G]:u\oplus s_j=v\}|$；其元素计数从左顶点 $u$ 到右顶点 $v$ 的边。两侧具有相同的索引集 $\{0,1\}^{\ell_G}$，其上函数 $\chi_b(u)=(-1)^{b\cdot u}$ 构成正交基。对每个 $b,u\in\{0,1\}^{\ell_G}$，
\[
 (\mathsf A_G\chi_b)(u)
 =\sum_{j=1}^{d_G}\chi_b(u\oplus s_j)
 =\chi_b(u)\sum_{j=1}^{d_G}(-1)^{b\cdot s_j}.
\]
因此每个 $\chi_b$ 都是特征向量：常函数 $\chi_{\vec{0}}$ 的特征值为 $d_G$，而 \eqref{eq:hash-walsh} 将其余所有特征值的绝对值界为 $\varepsilon d_G$。对于左顶点集合 $U$ 和右顶点集合 $V$，令 $e_G(U,V)$ 计数它们之间的边，计入重数。将 $U,V$ 的指示向量分解为其常数部分和正交部分并应用 Cauchy--Schwarz 得
\begin{equation}\label{eq:hash-mixing}
 e_G(U,V)\le
 \frac{d_G|U||V|}{q_G}
 +\varepsilon d_G\sqrt{|U||V|}.
\end{equation}

为指定局部码，令
\[
 c_G=\left\lceil\frac{d_G}{8B}\right\rceil,
 \qquad \delta=\frac1{100B\mu}.
\]
我们需要一个具有 $c_G$ 行、$d_G$ 列的二元矩阵 $A_{\mathrm{loc}}$，使得其模 $2$ 核中的每个非零字权重至少为 $\delta d_G$。为证明这样的矩阵存在，我们考虑一个元素为独立均匀比特的矩阵。对每个固定的非零字，其某个非零坐标使每一行中的奇偶性均匀；因此各行独立性给出该字位于核中的概率为 $2^{-c_G}$。令 $h=\lceil\delta d_G\rceil-1$。若 $h=0$，则没有权重小于 $\delta d_G$ 的非零字。否则它们的数目至多为
\[
 \sum_{j=1}^{h}\binom{d_G}{j}
 \le(h+1)(2d_G)^h
 \le(2d_G)^{2h}
 <2^{d_G/(25B)}.
\]
对于最后一个不等式，使用 $h<\delta d_G$、$d_G\le m<2^\mu$ 和 $\mu+1\le2\mu$。由于 $c_G\ge d_G/(8B)$，联合界表明合适的 $A_{\mathrm{loc}}$ 存在。枚举 $2^{c_Gd_G}$ 个二元矩阵并测试所有权重小于 $\delta d_G$ 的字，可在 $\exp(\poly(B,\log(m+2)))$ 次位运算中找到这样的矩阵。

在两侧的每个顶点处，按照端口顺序将 $A_{\mathrm{loc}}$ 的 $c_G$ 行施加到其关联边坐标上。这些约束定义一个 Tanner 码，其局部码为 $C_{\mathrm{loc}}=\ker_{\mathbb F_2}A_{\mathrm{loc}}$。其关联图对我们图的每条边有一个比特节点，对每个顶点有一个局部码节点；每个比特参与其两个端点处的约束。考虑一个被所有这些检验接受的非零边字。设 $U,V$ 为其非零边在两侧触及的顶点集合，$h'$ 为其权重。$U\cup V$ 中的每个顶点至少关联 $\delta d_G$ 条非零边，因此 $h'\ge\delta d_G\max(|U|,|V|)$。因此
\[
 \delta d_G\sqrt{|U||V|}
 \le h'\le e_G(U,V)
 \le\frac{d_G|U||V|}{q_G}
       +\varepsilon d_G\sqrt{|U||V|}.
\]
$U$ 和 $V$ 都非空。除以 $d_G\sqrt{|U||V|}$ 得 $\sqrt{|U||V|}\ge(\delta-\varepsilon)q_G$，从而
\[
 h'\ge\delta(\delta-\varepsilon)q_Gd_G
 =\frac{q_Gd_G}{20000B^2\mu^2}\ge\frac mP.
\]

通过保留 $q_Gd_G$ 个边坐标中的任意 $m$ 个，将它们与 $[m]$ 等同，并将其余坐标置零，来缩短这个 Tanner 码。将所有顶点检验限制到保留的坐标上便得到 $A^{(0)}$。任何被 $A^{(0)}$ 接受的字可通过补零扩展为被接受的边字，因此它具有所需的权重。每行至多有 $d_G\le D_0$ 个非零项。共有 $2q_Gc_G$ 行。若 $q_G=1$，则 $d_G=m$ 且
\[
 2c_G\le\frac{m}{4B}+2\le\frac n4+2\le3n.
\]
若 $q_G>1$，则 $d_G=D_0\ge8B$，且
\[
 2q_Gc_G\le\frac{q_Gd_G}{4B}+2q_G
 \le\frac{q_Gd_G}{2B}<\frac mB\le n.
\]
对移位、局部矩阵和检验字的搜索是顺序进行的。存储当前选择和最终行使用 $\poly(n,m,B)$ 位工作空间。
\end{proof}

\subsection{通过计数公共零点选择额外方程}
在初始检验之后，每个仍然通过的未满足赋值至少违反 $m/P$ 个子句。这使得一个短的随机检验能以常数概率抓住每个这样的赋值。为使选择具有确定性，我们使用条件期望背后的平均原理~\cite{Raghavan88}。

\begin{lemma}\label{lem:greedy-hashes}
使用来自引理~\ref{lem:sparse-checks} 的 $P,D_0$ 和 $A^{(0)}$，一个确定性过程追加 $3n$ 个长度为 $P$ 的列表行，并得到一个满足 $N(A)=\#F$ 的矩阵 $A$。它至多使用 $3n(2m)^P$ 次公共零点计数。每个候选系统至多有 $6n$ 个方程，每个方程涉及至多 $2^k\max(D_0,P)$ 个次数至多为 $k$ 的单项式。
\end{lemma}
\begin{proof}
在当前矩阵为 $A$ 的某一阶段，枚举有序列表
\[
 \gamma=((e_1,c_1),\ldots,(e_P,c_P))
 \in\mathcal C:=([m]\times\{0,1\})^P.
\]
列表 $\gamma$ 表示二元行 $h_\gamma$，其中
\[
 (h_\gamma)_e=
 \Bigl(\sum_{j:e_j=e}c_j\Bigr)\bmod2
 \qquad(e\in[m]).
\]
共有 $(2m)^P$ 个列表；当不同列表表示同一行时仍分别计数。用 $A;h_\gamma$ 表示将 $h_\gamma$ 追加到 $A$ 后得到的矩阵。

固定一个被 $A$ 接受的未满足赋值 $x$。初始行仍保留在 $A$ 中，因此
$|\operatorname{supp}v(x)|\ge m/P$。
对于均匀列表 $\gamma$，$e_1,\ldots,e_P$ 均不属于该支撑的概率至多为
\[
 (1-1/P)^P\le e^{-1}<1/2.
\]
在索引命中支撑的条件下，选择一个命中位置 $j$。在固定除 $c_j$ 外的所有系数后，均匀比特 $c_j$ 使 $h_\gamma\cdot v(x)$ 在模 $2$ 下均匀。当某些子句索引重复时这仍然成立。因此，对每个这样的 $x$，至少有四分之一的候选列表给出拒绝 $x$ 的行。

对仍然通过的未满足赋值求和此界，我们得到
\[
 \mathbb E_\gamma D(A;h_\gamma)\le\tfrac34 D(A).
\]
对每个候选，通过计数其多项式系统的公共零点来计算 $N(A;h_\gamma)$。追加达到最小值的候选，按枚举顺序打破平局。由于
$N(A;h_\gamma)=\#F+D(A;h_\gamma)$ 对每个 $\gamma$ 都成立，所选行至多留下 $\tfrac34 D(A)$ 个未满足赋值。经过 $3n$ 轮后至多剩下 $(3/4)^{3n}2^n=(27/32)^n<1$ 个；它们的数目是整数，故为零。该过程运行固定的轮数，且从不需要知道 $D(A)$ 或 $\#F$。

初始矩阵至多有 $3n$ 行，每个候选至多用 $3n$ 个额外行来评估。根据 \eqref{eq:hash-query}，列出 $s$ 个子句的行在原始变量上至多展开为 $2^k s$ 个单项式；重复项在 $\mathbb F_2$ 上相消。这些观察给出所声称的规模和查询界。
\end{proof}

\subsection{计数一个稀疏公式}\label{subsec:count-sparse}
假设 $n\ge20$ 且 $1\le m\le Bn$，其中 $B$ 仅依赖于 $k$ 并由 $\log n$ 的某个固定幂所界定。由引理~\ref{lem:sparse-checks} 构造初始矩阵。这使用
$\exp((\log n)^{O_k(1)})$ 次位运算和多项式空间。
$P$ 和 $D_0$ 都是 $(\log n)^{O_k(1)}$，因此引理~\ref{lem:greedy-hashes} 中的每个候选都是定理~\ref{thm:overview-sparse-polynomials} 的一个实例，带有单个固定的
$c=c(k)$。这既适用于被丢弃的候选，也适用于最终系统：计数器只需要方程的数目、次数和稀疏性。对固定的 $k$，展开每一行并合并单项式耗费多项式时间和空间。

最多需要 $3n(2m)^P+1=\exp((\log n)^{O_k(1)})$ 次调用，其中包括
一次最终计数。每次调用需要
$2^{O_k(n(\log\log n)^2/\sqrt{\log n})}$ 次位运算。由于 $\log n$
的任意固定幂都是 $o(n(\log\log n)^2/\sqrt{\log n})$，总时间
（包括构造与候选比较）为
\begin{equation}\label{eq:sparse-formula-cost}
 (n+m)^{O_k(1)}2^{O_k(n(\log\log n)^2/\sqrt{\log n})}.
\end{equation}
由贪心引理，最终计数为 $\#F$。不含子句或包含空子句的公式
直接处理。候选按顺序检查；除计数器的多项式工作空间外，
只需保留当前行、一个候选、最佳行及其计数。

\subsection{回到一般公式}\label{subsec:general-reduction}

我们使用不相交稀疏化，把满足赋值划分到同一组 $n$ 个声明
变量上的稀疏公式之间。前述算法对其中每个公式计数。
参数 $\sigma$ 控制一项权衡：将其调小会减少公式数量，但会
增大它们的出现次数上界 $B$。我们的求解器允许 $B$ 按
$\log n$ 的固定幂增长，因此可以让 $\sigma$ 趋于零，同时使公式
数量和每个公式的开销都保持在期望的时间界内。

\begin{lemma}\label{lem:sparsify}
固定 $k\ge2$。给定 $n$ 个声明变量上、含 $m$ 个子句的 $k$-CNF
以及有理数 $0<\sigma\le1$，一个确定性算法在相同变量上构造
一棵 $k$-CNF 树。设 $L_\sigma$ 为所给 $\sigma$ 的分子与分母的
二进制编码总长度。
一个赋值满足输入，当且仅当它满足某个叶公式，并且它不能
同时满足两个不同的叶。
该树有
\[
 \operatorname{poly}_k(n,\sigma^{-1})\,2^{\sigma n}
\]
个节点，深度为 $\poly_k(n,\sigma^{-1})$。
在每个叶中，每个变量至多出现 $B$ 次，其中 $B$ 是一个显式
计算出的整数，满足
\[
 B=O_k\!\Bigl(\bigl[\sigma^{-1}\log(2/\sigma)\bigr]^{k-1}\Bigr).
\]
输入规范化使用 $\poly_k(n,m)$ 次位运算。
构造数值参数使用 $\poly_k(L_\sigma)$ 次位运算。此后，处理
每个节点使用 $\poly_k(n,\log(2/\sigma))$ 次位运算。
输入计数为各叶计数之和。
\end{lemma}
附录~\ref{app:sparsify} 证明该引理，包括使叶计数可以相加的
不相交性。

\begin{proposition}\label{prop:interpolation-bound}
对每个固定的 $k\ge2$，插值算法在
$(n+m)^{O_k(1)}2^{O_k(\rho(n))}$ 次位运算和
$(n+m)^{O_k(1)}$ 位空间内，对 $n\ge20$ 个声明变量上、含 $m$
个子句的 $k$-CNF 的满足赋值计数。
\end{proposition}
\begin{proof}
我们首先在 $(n+m)^{O_k(1)}$ 时间内规范化输入：删除重复
文字、重言子句、重复子句以及包含另一子句的子句。我们保留
全部 $n$ 个声明变量。
空子句给出计数零，不含子句的公式给出计数 $2^n$。对于
$n<20$，我们用穷举搜索处理。
对于 $n\ge20$，我们接下来选择稀疏化参数。定义
\[
 \lambda_n=\lceil\log_2 n\rceil,\qquad
 \nu_n=\lceil\log_2\lambda_n\rceil,\qquad
 s=\left\lfloor\frac{\sqrt{\lambda_n}}{\nu_n^2}\right\rfloor.
\]
我们通过整数比较计算 $s$，取满足 $s^2\nu_n^4\le\lambda_n$ 的最大
非负整数。
若 $s<2$，我们使用穷举搜索；
$\sqrt{\lambda_n}/\nu_n^2$ 趋于无穷，因此这样的规模只有
有限多个。否则我们以 $\sigma=1/s$ 应用引理~\ref{lem:sparsify}。
该树至多有
$2^{n/s}\operatorname{poly}_k(n,s)$ 个节点，其叶的出现次数
上界为
\[
 B=O_k\bigl((s\log s)^{k-1}\bigr)=(\log n)^{O_k(1)}.
\]
每个叶使用原来的 $n$ 个变量，并至多有 $Bn$
个非空子句。我们用第~\ref{subsec:count-sparse} 节的稀疏公式
算法对每个叶计数，其运行时间由
\eqref{eq:sparse-formula-cost} 给出。不含子句或包含
空子句的叶直接处理。
该应用中的子句计数是叶中的子句数；最终界中的 $m$ 仍是
输入的子句数。

对于 $s\ge2$，其定义中的下取整给出
\[
 \frac ns=O\!\left(\frac{n\nu_n^2}{\sqrt{\lambda_n}}\right)
 =O(\rho(n)).
\]
因此叶的数量和对每个叶计数的时间的指数均为 $O_k(\rho(n))$。
它们的多项式因子，连同规范化和树节点处的工作，都由
$(n+m)^{O_k(1)}$ 界定。因此位运算总数为
\[
 (n+m)^{O_k(1)}2^{O_k(\rho(n))}.
\]
有限的穷举搜索情形被吸收到常数中。所有参数、初始检查和
附加行都由上述确定性过程计算；每次计数查询都使用
定理~\ref{thm:overview-sparse-polynomials} 的公共零点算法。

为界定完整算法所用的空间，我们深度优先遍历
稀疏化树，并一次对一个叶计数。
当前路径有 $\poly_k(n,s)$ 个节点，且每个规范化
公式有 $\poly_k(n)$ 位的编码。由于
$s\le\sqrt{\lceil\log_2 n\rceil}$，存储该路径并对一个叶
计数使用 $\poly_k(n,m)$ 位。不相交性使已处理的叶计数之和
保持在零与 $2^n$ 之间，因此其累加器
至多使用 $n+1$ 位。于是总空间为 $\poly_k(n,m)$ 位。
\end{proof}

\appendix

\section{划分满足赋值}\label{app:sparsify}

固定 $k\ge2$ 以及有理数 $0<\sigma\le1$。
我们通过反复把公式拆分为两种情形来证明引理~\ref{lem:sparsify}。
两个子公式使用相同的 $n$ 个变量，且父公式的每个
满足赋值恰好满足一个子公式。
我们遵循选择共享文字集的规则以及~\cite[第~6 节]{CIP} 的
插入记账方式。
使两种情形不相交的细化在精确
计数中是标准的~\cite[附录~A]{DHMTW}。

在本附录中，子句是文字集合，表示它们的
析取。其\emph{宽度}是所含文字的数量。
我们删除重言子句、重复子句以及包含另一子句的子句。因此
剩余子句在包含关系下构成反链。我们保留全部 $n$ 个声明
变量，且不传播单位子句。空子句给出计数零，空公式给出
计数 $2^n$。

\subsection{分支规则}

固定整数 $M\ge2$。我们将在每个叶处得到出现次数上界
$B=O_k(M^{k-1})$，然后选择 $M$，使树至多有
$\poly_k(n,\sigma^{-1})2^{\sigma n}$ 个节点。令
\[
 \theta_j=M^j\qquad(0\le j<k).
\]
以 $F$ 表示当前子句族，对每个文字集 $H$ 定义
\[
 \deg_c(H)=\#\{C\in F:|C|=c,\ H\subseteq C\}
\]
只要宽度为 $c$ 的子句的非空真子集 $H$ 满足
\begin{equation}\label{eq:spar-violation}
 \deg_c(H)\ge\theta_{c-|H|}.
\end{equation}
我们就进行分支。
算法选择尽可能小的 $c$，然后选择尽可能大的 $|H|$；
任何剩余的并列都确定性地打破。
因此，一个留下 $j=c-|H|$ 个其他文字的非空共享文字集 $H$，
在我们对其分支之前必须出现在至少 $M^j$ 个子句中。

设 $H\cup P_1,\ldots,H\cup P_z$ 为包含所选 $H$ 的全部宽度为 $c$
的子句，并令 $G$ 由其余子句组成。令
$j=c-|H|$。这里 $P_a$ 是从第 $a$ 个所选子句中删除
$H$ 后剩下的文字集，因此每个 $P_a$ 的宽度为 $j$，
且 $z\ge\theta_j$。
两个子公式为
\begin{equation}\label{eq:spar-children}
 F_A=G\wedge H,\qquad
 F_B=G\wedge\bigwedge_{a=1}^zP_a
          \wedge\bigwedge_{\ell\in H}\{\neg\ell\}.
\end{equation}
我们通过删除因其新子句而变得冗余的子句，在每个子公式中
恢复反链条件。
集合 $P_a$ 可能重叠。
两个子公式的满足赋值不相交：
$F_A$ 要求 $H$ 中某个文字为真，而 $F_B$
要求 $H$ 中每个文字都为假。

\begin{lemma}\label{lem:spar-branch}
\eqref{eq:spar-children} 中的两个子公式划分其父公式的满足
赋值。每个新子句的宽度都小于 $c$。
$A$-子公式插入一个新子句，$B$-子公式插入至少
$\theta_j$ 个宽度为 $j$ 的新子句。所有新插入的子句都能在
立即规范化后保留。任一子公式中删除的每个旧子句都
包含一个新插入的真子集。
\end{lemma}
\begin{proof}
每个赋值恰好满足\emph{$H$ 中某个文字为真}与
\emph{$H$ 中每个文字都为假}这两个条件之一。
在第一种情形中，所有所选子句都被满足，剩下
$G\wedge H$。在第二种情形中，每个所选子句归约为 $P_a$，
而单位子句 $\{\neg\ell\}$ 强制该情形。
因此两个子公式划分父公式的满足赋值。

旧反链不含被 $H$ 或任何 $P_a$ 包含的子句：
这样的子句会是 $H\cup P_a$ 的真子集。
因此 $H$ 与各不相同的等宽子句 $P_a$ 都是新的。
没有 $P_a$ 包含 $H$ 的变量，因为其父公式非重言。
因而，单位子句 $\{\neg\ell\}$（$\ell\in H$）都不是任何 $P_a$
的子集。
任何已存在的单位子句 $\{\neg\ell\}$ 都被保留，且
不计为新插入。
其余新单位子句与子句 $P_a$ 都得以保留。
对子公式 $A$，每个所选旧子句都包含新插入的子句 $H$；
对子公式 $B$，则包含其新插入的子集 $P_a$。
按规范化规则，其他每个被删除的子句都包含一个新插入的
真子集。
\end{proof}

\subsection{计数插入}
为界定树的大小，我们计数沿一条路径引入的每个子句，
包括后来被删除的子句。输入中已有的子句
不计入此数。由引理~\ref{lem:spar-branch}，每一步都
插入一个子句，因此界定插入总数也就
界定了路径长度。

在这种记账中，当某个分支引入一个子句时，该子句被标记。
首先我们界定仍存在且包含给定文字集的被标记子句有多少。
最小宽度优先规则控制分支之前存在的那些；
最大 $|H|$ 优先规则控制该分支引入的那些。

\begin{lemma}\label{lem:spar-marked-degree}
在每个节点处，对于 $1\le i<j<k$ 及每个大小为 $i$ 的文字集 $S$，
仍存在且包含 $S$ 的宽度为 $j$ 的被标记子句少于
$2\theta_{j-i}$ 个。
\end{lemma}
\begin{proof}
初始时没有被标记的子句。只需考虑插入宽度为 $j$ 的子句的分支。
其选定的宽度 $c$ 大于 $j$。在插入之前，包含 $S$ 的当前
宽度为 $j$ 的子句少于 $\theta_{j-i}$ 个：否则最小宽度
优先规则会选择宽度至多为 $j$ 的子句。

选择子公式 $A$ 只插入一个子句。
假设子公式 $B$ 插入宽度为 $j$ 的子句 $P_a$。
若没有 $P_a$ 包含 $S$，则包含 $S$ 的新宽度为 $j$ 的子句
数为零。否则 $S$ 落在某个花瓣中，因而与
$H$ 不相交。在此情形下，$P_a$ 包含 $S$ 当且仅当
所选子句 $H\cup P_a$ 包含 $H\cup S$。由于 $0<i<j$，
$H\cup S$ 大于 $H$，且是某个所选宽度为 $c$ 的子句的
真子集。
因此最大共享集优先规则给出
\[
 \#\{a:S\subseteq P_a\}
 =\deg_c(H\cup S)<\theta_{c-|H\cup S|}
 =\theta_{j-i}.
\]
其他新子句 $\{\neg\ell\}$ 的宽度为一。
因此宽度为 $j$ 的旧子句与新子句合起来满足
所称的界。删除不会使其增大。
\end{proof}

\begin{lemma}\label{lem:spar-insertion-bound}
沿任意路径或路径前缀，宽度为 $j$ 的插入总数 $T_j$
满足 $T_j\le4(3M)^{j-1}n$，其中 $1\le j<k$。
\end{lemma}
\begin{proof}
$n=0$ 的情形是直接的，故假设 $n\ge1$。
对于当前路径前缀，设 $L_j$ 为仍存在的被标记的
宽度为 $j$ 的子句数，并设 $D_j$ 为子句此后已被删除的
宽度为 $j$ 的插入数。
则 $T_j=L_j+D_j$。
单位子句从不被删除，故 $T_1\le2n\le4n$。
对于 $j\ge2$，对 $2n$ 个单文字求和引理~\ref{lem:spar-marked-degree}。
这给出 $L_j\le4n\theta_{j-1}$。

每当一个被标记的宽度为 $j$ 的子句被删除时，把该删除
分配给其中包含的一个新引入真子集；
引理~\ref{lem:spar-branch} 保证存在这样的子集。
若新子句的宽度为 $i<j$，引理~\ref{lem:spar-marked-degree}
表明，在分支之前紧接着，它被包含在少于 $2\theta_{j-i}$ 个
宽度为 $j$ 的被标记子句中。
因此，一个宽度为 $i$ 的子句的插入至多接收这么多
删除。对新子句的宽度 $i$ 求和可得
$D_j\le2\sum_{i=1}^{j-1}\theta_{j-i}T_i$。因此
\begin{equation}\label{eq:spar-insertions}
 T_j=L_j+D_j\le4n\theta_{j-1}
       +2\sum_{i=1}^{j-1}\theta_{j-i}T_i.
\end{equation}
第一项计数仍存在的子句；求和项计数已被删除的子句。除以
$nM^{j-1}$ 并应用归纳法可得
\[
 \frac{T_j}{nM^{j-1}}
 \le4+2\sum_{i=1}^{j-1}4\cdot3^{i-1}
 =4\cdot3^{j-1},
\]
如所需。
\end{proof}

\subsection{树的大小}
我们现在把引理~\ref{lem:spar-insertion-bound} 转化为对
树节点数的界。它同时控制一条路径的长度
以及该路径在
\eqref{eq:spar-children} 中选择 $B$-子公式的次数。将这些界分别写为 $Cn$ 与
$\Gamma n/M$，其中
\[
 C=4\sum_{j=0}^{k-2}(3M)^j,\qquad
 \Gamma=4\sum_{j=0}^{k-2}3^j=2(3^{k-1}-1).
\]
引理~\ref{lem:spar-insertion-bound} 给出深度至多 $Cn$。
选择 $|P_a|=j$ 的子公式 $B$ 至少插入 $M^j$
个宽度为 $j$ 的子句，因此一条路径至多包含
$4\cdot3^{j-1}n/M$ 次这样的选择。
总计，一条路径至多选择 $B$ 共 $\Gamma n/M$ 次。

选择 $B$ 次数少的路径有简短的描述。
由于分支规则是确定性的，一个节点由其深度及其路径选择 $B$ 的
位置确定；其他每个位置都选择 $A$。
对于 $M\ge\Gamma$，这将节点数界定为
\[
 (Cn+1)\sum_{a\le\lfloor\Gamma n/M\rfloor}\binom{Cn}{a}
 \le (Cn+1)^2
       2^{(\Gamma n/M)\log_2(4CM)}.
\]
这里 $\Gamma n/M\le n\le Cn/4$。
对于正整数 $a$，使用
$\binom{Cn}{a}\le(4Cn/a)^a$；该上界在
$a\le Cn/4$ 时递增。
至多有 $Cn+1$ 个求和项。
$\Gamma n/M<1$ 的情形是直接的。

我们选择 $M$，使该界中的指数至多为 $\sigma n$。
按递增顺序检验 $M\ge\Gamma$ 的二的幂。
对每个选择，计算 $C$ 与 $h=\lceil\log_2(4CM)\rceil$，
并在
\[
 \Gamma h\le\sigma M.
\]
时停止。
所有比较都是精确的。由于 $C=O_k(M^{k-2})$ 且 $\Gamma$
只依赖于 $k$，第一个成功的选择满足
\[
 M=O_k\bigl(\sigma^{-1}\log(2/\sigma)\bigr).
\]
因此该树至多有 $(Cn+1)^2\,2^{\sigma n}$ 个节点。

在一个叶处，\eqref{eq:spar-violation} 对子句的每个非空真
子集都不成立。
当 $c\ge2$ 时，一个固定文字属于少于 $\theta_{c-1}$ 个宽度为 $c$
的子句，并且至多属于一个单位子句。因此每个变量
至多出现
\[
 B=2\sum_{j=0}^{k-1}M^j
   \le4M^{k-1}
   =O_k\Bigl(\bigl[\sigma^{-1}\log(2/\sigma)\bigr]^{k-1}\Bigr).
\]
次。于是每个叶都具有所需的每变量出现次数界。
反复的两情形划分还表明，每个原始
满足赋值恰好出现在一个叶处，因此对各叶计数求和即得
原始计数。

为在节点处实现分支规则，我们对其当前公式计算 $\deg_c(H)$。
对每个宽度为 $c$ 的子句，我们枚举
其非空真子集 $H$，并递增以 $(c,H)$ 为索引的计数器。
至多有 $O_k(n^k)$ 个可能的非空规范化子句，
每个至多有 $2^k$ 个子集。子句计数器使用
$O_k(\log(n+2))$ 位，阈值使用 $O_k(\log M)$ 位。
因此，计算计数、选择分支并规范化其子公式，
每个节点使用 $\poly_k(n,\log M)$ 次位运算和位的
工作空间。对 $M$ 的界给出
$\log M=O_k(\log(2/\sigma))$。

取引理~\ref{lem:sparsify} 中的 $L_\sigma$，参数搜索
检查 $O_k(L_\sigma)$ 个二的幂。其整数操作数有
$O_k(L_\sigma)$ 位，因此构造数值参数使用
$\poly_k(L_\sigma)$ 次位运算和位的工作空间。

\section{精确算术与空间}\label{app:arithmetic}
第~\ref{sec:construction} 节界定了替换行 $\vec{P}_{f,v}(T)$ 的
有理系数。我们通过界定提升后的因子、
用于计算 $H$ 的泰勒多项式的和以及最终的
插值，来完成命题~\ref{prop:count} 中的位开销与空间
估计。要点在于每个和都使用固定的
分母。其分子随求和项数增长，因此其
位长只随该数目的对数增长。
对于有限多个不定元中的任意多项式 $G=\sum_\alpha c_\alpha \vec{X}^{\alpha}$，
记 $\|G\|_1=\sum_\alpha|c_\alpha|$ 为其
系数范数。

\subsection{提升指示式的整系数}
\label{subsec:rational-lift-height}
我们首先界定引理~\ref{lem:rational-idempotent} 所提供的
奇偶多项式的系数。它们出现在 \eqref{eq:lifted-factor} 中的每个方程
因子 $\mathcal B_i$ 中，因此其大小
进入评估这些因子的位开销界。
对于整数 $\ell\ge1$，回忆 $h(Y)=Y(Y-1)$ 以及
引理~\ref{lem:rational-idempotent} 中的多项式：
\[
 p_\ell(Y)=(1-Y)^\ell\sum_{j=0}^{\ell-1}
                    \binom{\ell+j-1}{j}Y^j
          =\sum_{j=0}^{\ell-1}e_j(Y)h(Y)^j,
 \qquad \deg e_j\le1.
\]
该引理证明了奇偶同余，并通过用 $h$ 逐次首一除法构造 $e_j$。
这里我们界定此构造中出现的整数。

\begin{lemma}\label{lem:rational-lift-height}
对 $0\le j<\ell$，写 $e_j(Y)=a_j+b_jY$。则 $a_j,b_j\in\mathbb Z$，且
\[
 |a_j|,|b_j|\le2^{C\ell^2}\qquad(0\le j<\ell)
\]
其中 $C$ 为绝对常数。
用 $h$ 对 $p_\ell$ 逐次除法，使用 $\poly(\ell)$ 次位运算和位的工作空间构造这些多项式。
\end{lemma}
\begin{proof}
系数范数不等式
$\|fg\|_1\le\|f\|_1\|g\|_1$ 连同二项式恒等式
\[
 \sum_{j=0}^{\ell-1}\binom{\ell+j-1}{j}
 =\binom{2\ell-1}{\ell-1}
\]
给出
\[
 \|p_\ell\|_1\le2^\ell\binom{2\ell-1}{\ell-1}\le2^{3\ell}.
\]
用 $h=Y^2-Y$ 作首一除法保持整系数。每次
首项消去使当前商与余式的系数范数之和
至多增加三倍。
所有逐次除法中总共有 $O(\ell^2)$ 次消去。
因此每个中间系数的绝对值至多为
$2^{O(\ell^2)}$，这证明了所述关于 $a_j,b_j$ 的界。
其带符号二进制编码因此使用 $O(\ell^2)$ 位。
每个系数向量有 $O(\ell)$
个条目。逐次除法对 $O(\ell^2)$ 位的整数进行 $O(\ell^2)$ 次整数算术
运算，给出所称的
$\poly(\ell)$ 位运算与空间界。
\end{proof}
\subsection{布尔赋值上求和式的位复杂度}
\label{subsec:rational-streaming-height}
我们现在界定由 \eqref{eq:full-taylor-sum} 计算
$\mathcal H_t(z)=H(t+z)\bmod z^J$ 时出现的每个中间值。
从替换行 $\vec{P}_{f,v}(T)$ 出发，我们界定替换后的因子、
它们在每组布尔变量上的和，以及命题~\ref{prop:count} 中
泰勒数据的累加器 $T_t$。
固定一个含 $n$ 个布尔变量的多项式系统 $\mathcal P$，它有
$m_{\mathrm{test}}$ 个次数至多为 $k$ 的方程、精度
$\ell$ 满足 $2^\ell>2^n$，并给其方程因子指定一个替换设计。
假设行 $\vec{P}_{f,v}$、划分 $S_{t,g}$ 与矩阵 $C_{t,g}$ 满足
定义~\ref{def:layout}，参数为 $Q,R,r,D,\tau$。
回忆 $H$ 的次数界以及每个节点保留的泰勒系数个数：
\[
 d_H=kR\bigl(m_{\mathrm{test}}+2(\ell-1)\bigr),\qquad
 J=\left\lfloor\frac{d_H}{Q}\right\rfloor+1.
\]
泰勒变量 $z$ 与公共截断标记 $V$ 中的所有多项式都以其模
$(z^J,V^\ell)$ 的唯一代表元存储。删去这些次数范围之外的项
不可能增大系数范数。

替换行有公分母
\[
 d_{\mathrm{den}}=(R-1)!Q!.
\]
我们取整数 $B_P\ge1$，使得每一行
$d_{\mathrm{den}}\vec{P}_{f,v}(T)$ 的每个整数系数的绝对值
至多为 $2^{B_P}$。定义~\ref{def:layout} 给出
$B_P=O(Q\log Q+\log D)$。
在固定节点 $t\in[Q]$ 与固定修正集 $U\subseteq[n]$ 处，
因子 $f$ 在变量 $v$ 处的输入为
\[
 \theta_{f,v}(z;x)=\langle \vec{P}_{f,v}(t), x\rangle+
       \bigl\langle \vec{P}_{f,v}(t+z)-\vec{P}_{f,v}(t),\mathbf1_U\bigr\rangle,
 \qquad x\in\{0,1\}^n.
\]
设 $L_i(z;x)$ 为方程 $i$ 的所列单项式经这些替换后的和。
每个 $d_{\mathrm{den}}\theta_{i,v}(z;x)$ 都是整数多项式，
其系数范数至多为
\[
 B_0=d_{\mathrm{den}}+3n(R+1)2^{B_P}(Q+1)^R .
\]
事实上，$\theta_{i,v}$ 中的三次行赋值每次在清分母后至多
贡献 $n(R+1)2^{B_P}(Q+1)^R$。该界对群和中的代表元同样成立。

令 $\Delta=d_{\mathrm{den}}^k$。对次数为 $u\le k$ 的单项式，
将其替换值乘以 $\Delta$ 得到范数至多为
$d_{\mathrm{den}}^{k-u}B_0^u\le B_0^k$ 的整数多项式，常数
单项式亦然。由于 $S$ 界定了列表总长度，
\[
 B=\max\{\Delta,SB_0^k\}
 \quad\text{满足}\quad \|\Delta L_i(z;x)\|_1\le B .
\]
因此 $\log d_{\mathrm{den}}=O(Q\log(Q+2))$，
$\log B_0=O(Q\log(Q+2)+\log(nD+2))$，且
$\log B=O(k\log B_0+\log(S+2))$。

方程 $i$ 的替换因子为
\[
 \mathcal B_i(z,V;x)=\sum_{j=0}^{\ell-1}
       e_j\bigl(L_i(z;x)\bigr)h\bigl(L_i(z;x)\bigr)^jV^j
       \pmod{(z^J,V^\ell)}.
\]
由于 $e_j$ 是整系数的且次数至多为 1，其替换的分母整除 $\Delta$；
$h(L_i)^j$ 的分母整除 $\Delta^{2j}$。于是，取
\[
 \Gamma=\Delta^{2\ell-1},
\]
时，$\Gamma\mathcal B_i(z,V;x)$ 的每个系数都是整数。
为界定所得的整数分子，令
\[
 E_\ell=\max\{1,\|e_0\|_1,\ldots,\|e_{\ell-1}\|_1\},\qquad
 B_1=\ell E_\ell(B+\Delta+2)^{2\ell-1}.
\]
对 $A_i=\Delta L_i$，两个替换因子的整数分子分别为
$\Delta e_j(L_i)$ 与 $A_i(A_i-\Delta)$，其范数分别至多
$E_\ell(B+\Delta)$ 与 $B(B+\Delta)$。因此把每一项补齐到
分母 $\Gamma$ 得到
\[
 \|\Gamma\mathcal B_i\|_1\le B_1,\qquad B_1\ge\Gamma,
 \qquad
 \log B_1=O\!\left(\ell^2+
       \ell\log(B+\Delta+2)+\log(\ell+1)\right).
\]

接下来把这些估计用于实际的群和。
记 $S_g=S_{t,g}$、$C_g=C_{t,g}$，并令 $c_g$ 为 $C_g$ 的行数。
于是 $c=\sum_g c_g\le2\tau n$，且布尔测量 $C_gx_g$ 的每个坐标
都落在 $[-D,D]\cap\mathbb Z$ 中。
设 $m_{\mathrm{test},g}$ 为在 $t$ 处分配给群 $g$ 的方程因子个数，
故 $\sum_g m_{\mathrm{test},g}=m_{\mathrm{test}}$。
由 \eqref{eq:group-polynomial} 给出的乘积 $A_{g,U}(x;z)$ 在
$\mathbb Q[z,V]/(z^J,V^\ell)$ 中表示。其分母整除
$\Gamma^{m_{\mathrm{test},g}}$，而相对于该分母的分子范数
至多 $B_1^{m_{\mathrm{test},g}}$。中性的常数一因子不影响此界。

对满足 $|E|<J$ 的 $U\subseteq E\subseteq[n]$ 与可达的测量元组
$\vec y=(y_1,\ldots,y_r)$，选取布尔代表元
$\bar x_h$ 使 $C_h\bar x_h=y_h$。\eqref{eq:group-sum} 中的群和为
\[
 \Sigma_g(\vec y)=
 \sum_{\substack{x_g\in\{0,1\}^{S_g}\\ C_gx_g=y_g\\
                  (x_g)_v=1\ (v\in E\cap S_g)}}
 A_{g,U}\!\Bigl(x_g+\sum_{h\ne g}\bar x_h;z\Bigr).
\]
这里各群上的向量按零延拓。
该和至多有 $2^{|S_g|}$ 项，且各项分母相同，故
\[
 \left\|\Gamma^{m_{\mathrm{test},g}}\Sigma_g(\vec y)\right\|_1
 \le 2^{|S_g|}B_1^{m_{\mathrm{test},g}}.
\]
对 $1\le g\le r$，定义部分乘积
$\Pi_g(\vec y)=\prod_{h=1}^g\Sigma_h(\vec y)$，每次乘法后截断。
其分母整除
\[
 \Gamma^{\sum_{h\le g}m_{\mathrm{test},h}}.
\]
因此每个完整贡献 $\Pi_r(\vec y)$ 都可放在固定分母
\[
 \Omega=\Gamma^{m_{\mathrm{test}}},
\]
之下，并满足
\[
 \|\Omega\Pi_r(\vec y)\|_1
 \le2^nB_1^{m_{\mathrm{test}}}.
\]

设 $T_t(z,V)$ 为累加器，它以符号 $(-1)^{|E|-|U|}$ 在所有
配对与测量元组上累加这些完整贡献。由 \eqref{eq:full-taylor-sum}，
其最终值为 $\mathcal A_\ell$ 值泰勒多项式 $\mathcal H_t(z)$。
一个节点处带符号更新次数 $M_*$ 的界为
\[
 M_*\le
 \biggl(\sum_{s=0}^{\min(J-1,n)}2^s\binom ns\biggr)(2D+1)^c
 \le3^n(2D+1)^c.
\]
因此 $T_t$ 的每个中间值（而不只是最终值）都满足
\[
 \|\Omega T_t\|_1
 \le M_*\,2^nB_1^{m_{\mathrm{test}}}.
\]
整数多项式 $\Omega T_t$ 的每个系数有
\[
 O\!\left(\log(M_*+1)+n+m_{\mathrm{test}}\log B_1\right)
\]
位。该估计也界定了上述固定分母下所有群和与部分乘积的整数分子。
分母满足
\[
 \log\Omega
 =k m_{\mathrm{test}}(2\ell-1)\log d_{\mathrm{den}}
 =O\!\left(k(m_{\mathrm{test}}\ell+1)Q\log(Q+2)\right).
\]

每个贡献至多乘以 $m_{\mathrm{test}}$ 个提升后的方程因子。继续处理更多配对或测量元组只会新增贡献，而不会把它们乘进先前的贡献。
这就是为什么指数多次迭代在位估计中只贡献
$\log(M_*+1)$。每个存储的多项式至多有
$J\ell\le(d_H+1)\ell$ 个系数。枚举只保留当前
索引与布尔代表元、一个群和、一个部分
乘积以及累加器 $T_t$，而不存储所有赋值或所有测量元组的表。
这些估计给出
\eqref{eq:sum-bit-bound} 以及每次有理运算的多项式位代价。
$QJ\ell$ 个泰勒系数与工作状态可用
$\poly(n,m_{\mathrm{test}},\ell,D,Q,R,J,r,k,S)$ 位容纳。

\subsection{插值与整数答案}
\label{subsec:rational-readout}
对 $0\le s<\ell$，令 $H_s(T)=[V^s]H(T)\in\mathbb Q[T]$。
第~\ref{subsec:rational-streaming-height} 节界定了对
$t\in[Q]$、$0\le j<J$ 计算出的数据
$[z^jV^s]\mathcal H_t(z)=[z^j]H_s(t+z)$。
固定 $s$。我们插值该坐标，对所有 $\ell$ 个坐标使用同一个有理
Hermite 矩阵。令 $M_H=QJ>d_H$。
在单项式基 $1,T,\ldots,T^{M_H-1}$ 下，Hermite 插值
矩阵是整数矩阵
\[
 \mathsf V_{(t,j),u}=
 \begin{cases}\binom uj t^{u-j},&u\ge j,\\0,&u<j,\end{cases}
 \qquad t\in[Q],\quad 0\le j<J,\quad 0\le u<M_H.
\]
由第~\ref{sec:counting} 节的 Hermite 插值结论，它非奇异。
其元素有
$O(M_H\log(Q+2))$ 位。
设 $\mathbf y$ 为按 $(t,j)$ 排序的计算系数
$[z^jV^s]\mathcal H_t(z)$ 组成的向量，$\mathbf h$ 为
$H_s$ 的单项式系数向量，用零补齐到长度
$M_H$。于是 $\mathsf V\mathbf h=\mathbf y$。
向量 $\mathbf b=\Omega\mathbf y$ 是整数向量；用
$B_{\mathrm{data}}$ 表示其元素位长的界。
第~\ref{subsec:rational-streaming-height} 节中的界给出
\[
 B_{\mathrm{data}}
 \le\poly(n,m_{\mathrm{test}},\ell,k,Q,\log(D+2),\log(S+2)).
\]

只需要常数系数 $h_0=H_s(0)$。设
$\mathsf V^{(0)}$ 由 $\mathsf V$ 把由 $u=0$ 索引的列
替换为 $\mathbf b$ 得到，并令
\[
 N_H=\det\mathsf V^{(0)},\qquad D_H=\det\mathsf V\ne0.
\]
对 $\mathsf V(\Omega\mathbf h)=\mathbf b$ 应用 Cramer 法则给出
\begin{equation}\label{eq:exact-readout}
 H_s(0)=h_0=\frac{N_H}{\Omega D_H}.
\end{equation}
我们用无分数高斯消元~\cite[第~I--II 节与第~IV--V 节]{Bareiss68} 计算这两个行列式。
对一个 $M_H\times M_H$ 整数矩阵，该算法使用 $O(M_H^3)$ 次
整数算术运算；主元选取使其也能处理奇异矩阵。保留的消元
元素是输入矩阵的带符号子式。每次更新对这些元素作乘积
与差，随后进行一次精确整数除法。因此，$\mathsf V$ 与
$\mathsf V^{(0)}$ 的子式界也界定了计算 $D_H$ 与 $N_H$
所用的中间整数。
$(\mathsf V\mid\mathbf b)$ 的每个子式的绝对值至多为
\[
 2^{O\left(M_H\left(B_{\mathrm{data}}+
                M_H\log(Q+2)+\log(M_H+1)\right)\right)}.
\]
事实上，其行列式是至多 $M_H!$ 个乘积之和，每个乘积
含至多 $M_H$ 个元素。因此无分数消元中的中间整数
在带符号二进制表示下有
\[
 O\!\left(M_H\bigl(B_{\mathrm{data}}+
                   M_H\log(Q+2)+\log(M_H+1)\bigr)\right)
\]
位。于是计算 $N_H,D_H$ 使用
$\poly(M_H,B_{\mathrm{data}})$ 次位运算与相应位的工作空间。

在 $T=0$ 处，由构造每个 $H_s(0)$ 都是整数，故
\eqref{eq:exact-readout} 中的除法对每个 $s$ 都精确。算出
这 $\ell$ 个整数后，令 $u=\sum_{s=0}^{\ell-1}H_s(0)$。
由 \eqref{eq:origin-residue}，最小非负剩余
\[
 u-2^\ell\left\lfloor\frac{u}{2^\ell}\right\rfloor=N(\mathcal P)
\]
就是精确计数。对相继坐标复用同一工作空间，全部 $\ell$
次插值、求和与取余使用
$\poly(M_H,B_{\mathrm{data}},\log(\Omega+1),\ell)$ 次位运算
与多项式多位的工作空间。

\subsection{技术概览中的算术}
\label{subsec:overview-arithmetic}
我们还为第~\ref{sec:technical-overview} 节中有界出现次数算法的
位复杂度与空间界给出依据。
引理~\ref{lem:row-coefficients} 用 $n$ 的多项式界定
替换系数的位长。求解
第~\ref{subsec:preview-separation} 节中多项式大小的有理线性系统，
也得到位长为多项式的测量系数 $\vec\lambda_{e,v,h}$。
令 $\delta$ 为所有这些系数与所有替换系数在所有求值点上的
分母之积。它们只有多项式多个，故 $\delta$ 的位长也是
多项式的。每个子句乘积至多使用 $M\le3n$
个线性输入；因此受限和中所有系数都可用公分母
$\delta^{3n}$ 累加。
多项式多个位长为多项式的输入系数之积，以及指数多个此类项
在该公分母上的和，其位长仍是多项式的。因此上述精确
Hermite 插值对这些泰勒数据同样具有多项式位代价。最后，枚举
可以顺序进行，只保留构造、当前索引、系数数组、运行和，
以及插值所需的多项式多个泰勒系数。这给出
多项式空间与仅多项式级的位运算开销。

\raggedbottom

\section{反驳 ETH 的其他推论}\label{app:consequences}
定理~\ref{thm:main} 给出一个电路下界，以及针对输出局部性
为常数的布尔映射求逆的次指数时间算法。固定密度情形给出
推论~\ref{cor:quantitative-circuits} 中的定量电路
下界。两个电路论证都只使用我们 $3$-SAT 算法的判定
版本。精确原像计数则用到其完整的计数保证。

\subsection{$\mathrm E^{\mathrm{NP}}$ 的电路下界}\label{subsec:circuit-lower-bounds}
对下面的定性界，用 $\mathrm{SIZE}(O(n))$ 表示具有 $O(n)$ 个
有界扇入 AND、OR 与 NOT 门的非一致电路所识别的语言。

\begin{corollary}\label{cor:consequence-circuits}
我们有
\[
 \mathrm E^{\mathrm{NP}}\not\subseteq\mathrm{SIZE}(O(n)).
\]
\end{corollary}
\begin{proof}
Williams 证明了：若 ETH 不成立，则
$\mathrm E^{\mathrm{NP}}$ 没有线性规模
电路~\cite[定理~6.1]{Wil13}。把这一蕴含用于
定理~\ref{thm:main} 提供的确定性次指数时间 $3$-SAT 算法。
\end{proof}

这一联系是从更快的可满足性算法推出电路下界的一个实例。
为编码一个有 $s$ 个门和 $n$ 个输入的布尔电路，为每个门
引入一个变量，由输入强制其取值，并加入一个要求输出位
等于 $1$ 的单位子句。所得的 $3$-CNF 有 $O(n+s)$ 个变量与子句，
且恰在电路可满足时可满足。因此我们的定理在 $2^{o(n)}$
时间内判定线性规模电路的可满足性。
Williams 的论证把这样的加速与假设的见证小电路结合，
以与非确定性时间层次发生矛盾。该推论给出一个
属于 $\mathrm E^{\mathrm{NP}}$ 但没有线性规模电路
族的语言；它并不断言 NP 的超多项式电路下界。

\begin{proof}[推论~\ref{cor:quantitative-circuits} 的证明]
我们使用 Williams 的见证电路论证~\cite[定理~6.1]{Wil13}，
并结合 Jahanjou、Miles 和
Viola 的局部归约~\cite[定理~1 和第~6 节]{JMV13}。

根据 \v{Z}\'ak 的非确定性时间层次定理~\cite{Zak83}，采用
~\cite[完整版中的定理~2.1]{Wil13} 所陈述的形式，固定一个一元语言
$H\in\operatorname{NTIME}(2^n)\setminus\operatorname{NTIME}(2^{n/2})$，
它定义在多带图灵机上。对所有充分大的 $n$，将局部归约应用于其固定机器在 $1^n$ 上的运行，
局部性参数取 $r=2$。时间配额
$T=2^n n^{O(1)}$ 包含了在归约中模拟该机器的多对数级开销，
并给出 $q=n+O(\log n)$ 个子句索引位。该归约在多项式时间内构造一个解码器，
它在每个索引上给出三个变量标签及其符号。令 $\varphi_n$
由全部 $2^q$ 个索引上解码得到的子句组成。它可满足
当且仅当 $1^n\in H$。

解码器的 $O(n)$ 个输出位中的每一个都具有常数深度决策
树，因此至多依赖于固定数目 $w$ 个索引位。
在所有汉明重量至多 $w$ 的串上对该解码器求值
即可恢复每个支撑集。事实上，见证对
某一坐标的依赖的两个串可以在支撑集之外置零，同时保留
它们不同的输出。同样的求值给出每个支撑集上的
真值表。这在多项式时间内构造出一个由 $O(n)$ 个二输入门
构成的解码器。

对固定的正整数 $d$，将每个变量标签填充到 $dn$ 位，
并声明全部 $2^{dn}$ 个变量。若对某个这样的 $n$ 有 $|i|=dn$ 且
$\varphi_n$ 可满足，则令 $W(i)$ 为其
字典序第一个满足赋值的第 $i$ 位，其中变量按其
标签排序。在所有其他情形令 $W(i)=0$。$i$ 的
长度确定 $n$。为判定 $W$，写出这 $2^q$ 个子句，向
NP 谕示询问该公式是否可满足；若可满足，
则通过询问其当前前缀后接零是否具有可满足的扩展
来恢复第一个赋值。至多 $2^{dn}$ 次长度为 $2^{O(n)}$ 的前缀
询问即足够，故 $W\in\mathrm E^{\mathrm{NP}}$。

令 $a(t)=\sqrt{\log_2 t}$。检验定理~\ref{thm:main} 的
常数密度部分中的计数是否为正，可在时间
$V^b2^{A V/a(V)}$ 内判定具有 $V$ 个变量、至多 $5V$ 个子句的 $3$-SAT，其中 $A,b>0$ 固定。我们将
对数底数的变化吸收进 $A$。该求解器使用多项式空间，因此
多带模拟可以在每次位访问时以代价 $V^{O(1)}$ 扫描其工作空间，
该代价被吸收进 $b$。

固定正有理数 $\eta<1/(12A)$，并对 $n\ge2$ 令
$s_n=\lfloor\eta n\lfloor a(n)\rfloor\rfloor$。
该阈值可利用 $n$ 的二进制长度和
$\lfloor\log_2 n\rfloor$ 的整数平方根计算。假设对
所有充分大的 $n$ 有 $C_W(dn)\le s_n$。为判定 $H$，
拒绝非一元输入。在 $1^n$ 上，猜测一个作用于 $dn$ 个索引位、
至多 $s_n$ 个门的电路。将解码器的三个变量标签
输入该电路的三个副本，并检验被索引的子句是否
被违反。若测试电路不可满足则接受。当
$1^n\in H$ 时，计算长度为 $dn$ 的 $W$ 的电路给出一个
可接受的猜测。当 $1^n\notin H$ 时，每个猜测都违反某个子句。

测试电路有 $q$ 个输入和至多 $3s_n+O(n)$ 个门。上述
编码为每个输入和门使用一个变量，为每个门至多使用四个
宽度至多三的子句以强制其真值表，
并使用一个单位子句固定输出。用未使用变量
填充后，对固定整数 $B$ 得到 $U_n=3s_n+Bn$ 个变量和至多
$5U_n$ 个子句。由于 $U_n=(3\eta+o(1))n a(n)$ 且 $a(U_n)\sim a(n)$，
该检验耗时 $2^{(3A\eta+o(1))n}$。猜测和构造
电路花费多项式时间，因此对所有充分大的 $n$，$H$ 可在时间
$2^{n/2}$ 内非确定性判定。有限多个
例外长度可以被硬连线，这与 $H$ 的选取矛盾。
因此对无穷多个 $n$ 有 $C_W(dn)>s_n$。由于
$s_n\sim\eta n\sqrt{\log_2 n}$ 且 $d$ 固定，取 $L=W$ 及任意
$0<c<\eta/(d\sqrt{\log 2})$ 便证明了 \eqref{eq:quantitative-circuit-lower-bound}。
\end{proof}


\subsection{输出局部性为常数的布尔映射求逆}\label{subsec:local-inversion}
固定常数 $d$。考虑布尔映射
$G:\{0,1\}^n\to\{0,1\}^m$，其第 $i$ 个输出位形如
\[
 G(\vec{x})_i=g_i(\vec{x}|_{S_i}),
 \qquad S_i\subseteq[n],\quad |S_i|\le d.
\]
输入描述 $G$ 时列出集合 $S_i$ 与布尔函数 $g_i$ 的
真值表。因此每个输出只查看常数多个输入位，尽管集合
$S_i$ 可以任意重叠。给定目标 $\vec{b}\in\{0,1\}^m$，
求逆问题要求找出任意满足 $G(\vec{x})=\vec{b}$ 的 $\vec{x}$。

\begin{corollary}\label{cor:consequence-local-inversion}
对每个固定 $d$ 与 $n\ge20$，数
$|G^{-1}(\vec{b})|$ 可以精确计算，且只要存在原像就能
找到一个，所用确定性时间为
\[
 (n+m)^{O_d(1)}2^{O_d(\rho(n))}
\]
空间为 $(n+m)^{O_d(1)}$。
\end{corollary}
\begin{proof}
对每个 $i$ 与每个满足 $g_i(\vec{\alpha})\ne b_i$ 的局部赋值
$\vec{\alpha}\in\{0,1\}^{S_i}$，加入子句
\[
 \bigvee_{j\in S_i}\ell_{j,\alpha_j},
 \qquad
 \ell_{j,0}=X_j,\quad \ell_{j,1}=\neg X_j.
\]
该子句恰好禁止局部赋值 $\vec{\alpha}$。
若 $S_i$ 为空，则其约束为常数，可直接检查。其余子句的合取是
原始 $n$ 个声明变量上的 CNF，宽度至多 $d$，子句数至多
$2^d m$。其满足赋值恰为
$G^{-1}(\vec{b})$；无需辅助变量或多重性修正。
取 $k=\max\{2,d\}$ 应用定理~\ref{thm:main}。
通常的逐变量自归约至多再用 $n$ 次计数查询找到一个原像，
这只改变多项式因子。每次查询可在下一次之前丢弃，
从而保持多项式空间。
\end{proof}

当原像集非空时，条件计数还允许以相同界（期望时间至多
多项式因子）精确均匀地采样一个原像~\cite{JVV86}。
特别地，该推论适用于 Goldreich 候选单向函数的常数局部性
版本，它们在输入位的重叠子集上求值局部谓词~\cite{Gol00}。
对多项式多个输出位，它给出对每个像中原像的 $2^{o(n)}$ 时间
求逆算法，从而排除这些族的 $2^{\Omega(n)}$ 求逆困难性。
这并不排除相对于多项式时间算法的单向性。

\section{精确计数的更紧界}\label{app:quantitative}

我们通过收缩布尔张量网络来证明定理~\ref{thm:main}。
这两个运行时间反映了网络中导线的数量。
含有 $m\le Dn$ 个子句的公式给出一个具有 $O_{k,D}(n)$
条导线的网络。对于任意公式，我们使用
第~\ref{sec:hashes} 节的奇偶检验。引理~\ref{lem:total-order}
的次数界使我们能用 $O_k(n\log\log n)$ 条导线来索引其多项式
提升的乘积中保留的一个单项式。我们在证明它们所用的收缩
定理之前给出这两种编码。

\subsection{布尔张量收缩}
\label{subsec:tensor-interface}

布尔张量网络是一个二部多重图 $G$，其中每条边携带一个比特。
每个顶点 $v$ 有一个关联边的有序列表 $\delta(v)$ 和一个张量
$T_v$，后者为这些边的每个布尔赋值指派一个环元素。平行边
具有不同的端口。该网络的\emph{收缩}为
\[
 Z=\sum_{x\in\{0,1\}^{E(G)}}\prod_{v\in V(G)} T_v(x|_{\delta(v)}).
\]
孤立顶点贡献其标量元素。我们用边数来度量网络的
大小。

\begin{theorem}[布尔张量收缩]\label{thm:tensor-contraction}
固定正常数 $C,A,b$。设 $s\ge2$ 为整数，并设
$\mathcal L$ 是一个二部张量网络的图与张量描述的总长度，
该网络至多有 $Cs$ 条边且最大度至多为 $\lceil s^{9/10}\rceil$。
假设顶点数被 $s$ 的某个固定多项式所界，且各元素属于
\[
 \mathcal B_{P_0,\ell}
   =(\mathbb Z/2^{P_0}\mathbb Z)[V]/(V^\ell),
 \qquad 1\le P_0,\ell\le Cs.
\]
假设对每个顶点及其端口的每种布尔取值一致地，其元素可由
该描述在时间 $(s+\mathcal L)^b2^{A s^{9/10}}$ 与空间
$(s+\mathcal L)^b$ 内计算出来。那么该收缩可在
\[
 (s+\mathcal L)^{O(1)}
  2^{O(s/\sqrt{\log s})}
\]
位运算与 $(s+\mathcal L)^{O(1)}$ 比特空间内确定性计算。
常数只依赖于假设中的固定界。特别地，
允许 $\ell=1$。
\end{theorem}

第~\ref{app:tensor} 节中的算法用沿导线的递归取代单次全局
泰勒展开。在每次调用时，它形变导线的等同关系，使得在选定的
插值节点处，只有同一组内的端口发生相互作用。在该节点附近，
组之间的相互作用按称为模式的变量展开。权重为 $j$ 的模式的
系数可被 $2^j$ 整除，因此只有总权重低于工作精度的单项式
才能有贡献。用节点的 2 进赋值渐缩 Hermite 重数，从而
界定该精度。在 $M$ 条导线与 $u$ 个保留的赋值层级下，
模式单项式的数目为 $2^{O(M/\sqrt{u})}$，且每个单项式分解为
至多 $M/8$ 条导线的收缩之积。在递归深度 $d$ 处，
$M=O(s8^{-d})$ 且 $u=\Theta((\log s)/2^d)$。因此这些指数
按几何级数求和为 $O(s/\sqrt{\log s})$。

\subsection{固定子句密度}\label{subsec:sharp-fixed-density}

\begin{proof}[定理~\ref{thm:main} 中固定密度界的证明]
给每个子句一个顶点，该顶点携带其布尔满足表，子句中
每个不同变量对应一个端口。重言式子句可以删除，空子句
使计数为零。对于具有 $d\ge1$ 个端口的变量，要求它们的
比特一致。对 $d=1$ 附加一元张量 $(1,1)$，对 $d=2$ 附加
二元相等张量。对 $d\ge3$，取一棵具有 $d$ 个外部端口、
内部顶点为三价相等张量的树。这样的树有 $d-2$ 个内部
顶点，其端口上的每个公共值唯一地延伸到其所有
内部边。

至多有 $km$ 个子句端口。该图有 $O_k(n+m)$ 条边和顶点，
最大度至多为 $\max\{k,3\}$。用二元相等张量把每条边
细分一次。旧顶点与新顶点构成二部网络的两侧。细分边上
的两个比特必须相等，因此这保持收缩不变。出现变量上的
每个满足赋值贡献一，其他每个赋值贡献零。

取 $s=n+m$，在 $\mathbb Z/2^{n+1}\mathbb Z$ 上应用
定理~\ref{thm:tensor-contraction}，取 $P_0=n+1$ 与 $\ell=1$。
对固定的 $k$，局部表的大小有界。次数与精度的界对
所有充分大的 $s$ 成立，而有界实例可直接枚举。若有 $u$ 个
声明变量未被使用，则把收缩乘以 $2^u$ 并对
$2^{n+1}$ 取模。其最小非负剩余为 $\#F$，因为
$0\le\#F\le2^n$。这包括空公式，其空网络的收缩为一。
时间为 $s^{O_k(1)}2^{O_k(s/\sqrt{\log s})}$，空间为多项式。
当 $m\le Dn$ 时，有 $n\le s\le(1+D)n$，给出指数
$O_{k,D}(n/\sqrt{\log n})$。
\end{proof}

\subsection{任意子句密度}\label{subsec:sharp-dense}

第~\ref{sec:hashes} 节中的稀疏化与贪心行选择把这一情形
归约为计数稀疏多项式方程的公共零点。我们保留
第~\ref{sec:setup} 节的奇偶提升。其次数界使我们能用
$O(n\log\log n)$ 条布尔导线，通过检验—变量关联的
子集来编码一个单项式。

\begin{lemma}[通过收缩计数稀疏多项式系统]
\label{lem:tensor-sparse-query}
固定整数 $k,c\ge1$。设 $P_1,\ldots,P_a$ 是 $\mathbb F_2$ 上
关于 $n\ge20$ 个声明布尔变量的多项式。假设
$a\le cn$，且每个 $P_i$ 被列为至多
$\lceil\log_2 n\rceil^c$ 个次数至多为 $k$ 的单项式。公共零点的
数目可在 $n^{O_{k,c}(1)}2^{O_{k,c}(n\log\log n/\sqrt{\log n})}$
位运算与 $n^{O_{k,c}(1)}$ 比特空间内确定性计算。
\end{lemma}

\begin{proof}
令 $\lambda=\lceil\log_2 n\rceil$ 与
$\nu=\lceil\log_2\lambda\rceil$。
在每个单项式内，把变量的正幂次替换为其一次幂，并按模二
合并重复的单项式。检验 $i$ 的作用域由其剩余单项式中的
变量组成。其大小至多为 $k\lambda^c$。
设 $N_0$ 为各作用域大小之和，其中若变量 $v$ 属于检验 $i$，
则对 $(i,v)$ 只计一次。于是 $N_0\le ak\lambda^c$。
以所有可能的方式固定属于多于
\[
 D_*=\max\{1,\lceil N_0/n\rceil\}\lceil\sqrt{\lambda}\rceil
\]
个作用域的变量。这样的变量至多有 $n/\sqrt{\lambda}$ 个。它们的
赋值把原计数划分为至多 $2^{n/\sqrt{\lambda}}$ 个分支计数。
在每个分支上代入固定比特并再次合并重复单项式。检验 $i$ 的
自由作用域 $I_i$ 的大小仍至多为 $k\lambda^c$，且每个自由
变量出现在至多 $D_*=\lambda^{O_{k,c}(1)}$ 个作用域中。不出现
在任何作用域中的自由变量仍保持声明状态。常值一方程使
分支计数为零。我们可以丢弃常值零方程；若没有剩余方程，
则该分支的答案为二的自由于变量数目的幂。下文 $a$ 表示
剩余方程的数目，$n$ 表示原始变量数。

\paragraph{关联上的多项式。}
给每个满足 $v\in I_i$ 的对 $(i,v)$ 一个自己的未定元
$z_{iv}$。记 $z_i=(z_{iv})_{v\in I_i}$。用整数系数求规范化
分支多项式剩余单项式之和，以 $z_{iv}$ 代替 $x_v$，构成
$\widetilde L_i(z_i)$。我们使用环
\[
 \mathcal B=(\mathbb Z/2^p\mathbb Z)[V]/(V^\ell),
 \qquad p=\ell=n+1 .
\]
取 $h(Y)=Y(Y-1)$ 以及引理~\ref{lem:rational-idempotent} 中的
整系数多项式 $e_j$。将其系数模 $2^p$ 约化，并构成因子
\[
 B_i(z_i,V)=\sum_{j=0}^{\ell-1}
   e_j(\widetilde L_i(z_i))\,h(\widetilde L_i(z_i))^j V^j .
\]
设 $\overline B_i$ 为其在 $z_i$ 中的多线性代表，它通过在 $V$ 的
每个系数中把每个关联未定元的正幂次替换为一次幂而得到。它在
$z_i$ 的每个布尔赋值上与 $B_i$ 一致。在 $\mathcal B[z]$ 中定义
$F(z,V)=\prod_i\overline B_i(z_i,V)$。不同检验的关联未定元
互不相交，因此即使在把同一个布尔 $x_v$ 代入变量 $v$ 的所有
副本之前，$F$ 也是多线性的。

$\overline B_i$ 中 $V^j$ 的系数的关联次数至多为 $k(1+2j)$，
因为 $\deg e_j\le1$、$\deg h=2$，且多线性化只会降低单项式的
次数。$\mathcal B$ 中的乘法只保留满足 $\sum_i j_i\le\ell-1$
的项。从而 $F$ 中每个关联单项式的次数至多为
\begin{equation}\label{eq:sharp-incidence-degree}
 K=k\bigl(a+2(\ell-1)\bigr)=O_{k,c}(n).
\end{equation}
设 $\mathcal S_\ell$ 在 $\mathcal B$ 中求 $1,V,\ldots,V^{\ell-1}$
的系数之和。在 $z_{iv}=x_v$ 处，多线性化因子
$\overline B_i$ 与规范化分支多项式的 $B_i$ 一致。
整数 $\widetilde L_i((x_v)_{v\in I_i})$ 与原 $P_i$ 在 $x$ 与
固定比特处取值具有相同的奇偶性。对这些整数应用
引理~\ref{lem:total-order} 表明，分支计数模 $2^p$ 同余于
\begin{equation}\label{eq:sharp-branch-residue}
 \mathcal S_\ell\left(\sum_x F\bigl((z_{iv}=x_v)_{iv},V\bigr)\right),
\end{equation}
其中 $x$ 遍历所有自由布尔赋值。我们在截断环中乘起各检验
因子之后应用系数求和映射。

\paragraph{索引单项式。}
分别把检验与自由变量划分为 $J=\lceil n^{1/4}\rceil$ 组，
每侧各组的大小至多相差一。用 $\mathcal T_g$ 表示检验组，
$\mathcal X_h$ 表示变量组。一组包含 $O_c(n^{3/4})$ 个对象，
因此其关联数在任一侧至多为
$n^{3/4}\lambda^{O_{k,c}(1)}$。设 $N_{gh}$ 为从 $\mathcal T_g$ 到
$\mathcal X_h$ 的关联数，并设 $N=\sum_{g,h}N_{gh}\le N_0$。

对关联单项式，设 $d_{gh}$ 为其连接组 $g,h$ 的入选关联数。
枚举满足 $0\le d_{gh}\le N_{gh}$ 与 $\sum_{g,h}d_{gh}\le K$ 的
整数数组 $(d_{gh})$。至多有
\[
 (K+1)^{J^2}=2^{O_{k,c}(\sqrt n\log n)}
\]
个数组，可用 $J^2$ 个计数器枚举。对固定数组，这一对存在
\[
 C_{gh}=\binom{N_{gh}}{d_{gh}}
\]
个指定大小的子集。在两个组顶点之间引
$\lceil\log_2 C_{gh}\rceil$ 条布尔导线，以字典序索引它们。
当 $C_{gh}=1$ 时，两个端点使用唯一的子集（可能是整个关联集）
且不使用导线。任何未使用的比特串给出张量元素零。为解码
一个有效索引，按顺序取各关联，并用二项式系数计数包含或
不包含每个关联的可能补全。这些具有多项式比特长度，因此
解码使用多项式时间与空间。

所有对上的子集之并指定了 $N$ 个未定元中一个次数至多 $K$ 的
无平方单项式。计入该次数界的所有单项式，由隔板法给出上界
$\binom{N+K}{K}$。因此该数组对应网络中的导线数 $M$ 满足
\begin{align}
 M &\le J^2+\log_2\prod_{g,h} C_{gh}
      \le J^2+\log_2\binom{N+K}{K} \notag\\
   &\le J^2+K\log_2\!\left(\frac{\mathrm e(N+K)}{K}\right)
      = O_{k,c}(n\nu).
 \label{eq:sharp-wire-budget}
\end{align}
这里 $K=\Theta_{k,c}(n)$ 且 $N\le n\lambda^{O_{k,c}(1)}$。
又 $\lceil\log_2 C_{gh}\rceil\le N_{gh}$，故每个顶点的度至多为
其关联数 $n^{3/4}\lambda^{O_{k,c}(1)}$。

\paragraph{张量元素与收缩。}
导线比特的一个有效赋值在每个检验处确定一个子集
$S_i\subseteq I_i$。记 $z_{i,S_i}=\prod_{v\in S_i}z_{iv}$。检验组
$g$ 处的元素是其所指单项式在
$\prod_{i\in\mathcal T_g}\overline B_i(z_i,V)$ 中的系数。由于各
$z_i$ 互不相交，这等于
\[
 \prod_{i\in\mathcal T_g}[z_{i,S_i}]\overline B_i(z_i,V)
 \qquad\text{在 }\mathcal B\text{ 中}.
\]
对变量组 $h$，对 $v\in\mathcal X_h$ 令 $b_v=|\{i:v\in S_i\}|$。
其元素为
\[
 \sum_{x\in\{0,1\}^{\mathcal X_h}}
        \prod_{v\in\mathcal X_h}x_v^{b_v}
   =2^{|\{v\in\mathcal X_h:b_v=0\}|}\pmod{2^p}.
\]
确实，正指数要求 $x_v=1$，而零指数使 $x_v$ 自由。这包括
所有检验作用域之外的变量。

$F$ 的每个单项式恰好出现在一个数组中，并在每个组对处恰有
一个有效索引。它在 $F$ 中的系数是各检验组系数之积。在把具有
相同 $v$ 的关联设为相等之后，其在布尔赋值上的和是各变量组
元素之积。由 \eqref{eq:sharp-incidence-degree}，对这些数组
求和穷尽了所有单项式。从而 \eqref{eq:sharp-branch-residue}
中 $\mathcal S_\ell$ 内 $\mathcal B$ 的元素是这些二部网络的
收缩之和。

为计算一个检验侧系数，设 $\mathbf1_T$ 为 $I_i$ 上支撑集为
$T\subseteq S_i$ 的布尔向量。在 $\mathbf1_T$ 处，该多线性
多项式的值等于其以下标为 $T$ 的子集所索引的系数之和，且
在该处与 $B_i$ 一致。反解此子集和给出
\[
 [z_{i,S_i}]\overline B_i(z_i,V)=
 \sum_{T\subseteq S_i}(-1)^{|S_i|-|T|}B_i(\mathbf1_T,V)
 \qquad\text{在 }\mathcal B\text{ 中}.
\]
每一项花费多项式多个位运算。由所列单项式计算
$\widetilde L_i(\mathbf1_T)$，并使用引理~\ref{lem:rational-idempotent}
中的多项式 $e_j$，它们可在多项式时间内构造。存储
$B_i(\mathbf1_T,V)$ 的 $\ell$ 个系数模 $2^p$。至多有
$2^{|S_i|}\le 2^{k\lambda^c}$ 个子集需要处理。每次累加一个
检验，并在组内于 $\mathcal B$ 中乘起所得系数。这以
$2^{\lambda^{O_{k,c}(1)}}$ 位运算与多项式空间计算一个检验组
元素。变量侧元素与无效导线索引的检查花费多项式时间。
图、次数数组与元素计算过程的描述具有多项式大小。

以规模参数 $s=n\nu$、根精度 $p=n+1$ 与截断阶 $\ell=n+1$
应用定理~\ref{thm:tensor-contraction}。由
\eqref{eq:sharp-wire-budget}，该图有 $2J$ 个顶点与
$O_{k,c}(s)$ 条导线。对所有充分大的 $n$，其最大度至多为
$\lceil s^{9/10}\rceil$，其精度与截断阶为 $O(s)$，其元素计算
时间至多为 $2^{O_{k,c}(s^{9/10})}$。每次收缩随后花费
$n^{O_{k,c}(1)}2^{O_{k,c}(s/\sqrt{\log s})}$ 位运算与多项式
空间。这里 $s/\sqrt{\log s}=O(n\nu/\sqrt\lambda)$。
高关联分支与次数数组至多把运行时间乘以
\[
 2^{\,n/\sqrt{\lambda}+O_{k,c}(\sqrt n\log n)}.
\]
该指数中的两项均为 $O_{k,c}(n\nu/\sqrt{\lambda})$。
顺序处理各分支与数组。在每个分支上，于 $\mathcal B$ 中
求和各收缩并应用 $\mathcal S_\ell$。模 $2^{n+1}$ 相加各分支
剩余即得原计数，因为它介于 $0$ 与 $2^n$ 之间。这证明了
所声称的界，其中有界例外规模采用直接枚举。
\end{proof}

\begin{proof}[定理~\ref{thm:main} 中任意密度界的证明]
如第~\ref{subsec:general-reduction} 节那样，在其原始 $n$ 个
声明变量上规范化该公式。空子句与不含子句的公式有直接
答案。对 $n\ge20$，取
$\sigma=1/\lceil\sqrt{\lambda}\rceil$ 应用引理~\ref{lem:sparsify}，
其中 $\lambda=\lceil\log_2 n\rceil$。该树有
$n^{O_k(1)}2^{O(n/\sqrt{\lambda})}$ 个节点与多项式深度。
每个叶子的出现次数界为
$B=O_k((\sigma^{-1}\log(2/\sigma))^{k-1})=\lambda^{O_k(1)}$，
且至多有 $Bn$ 个非空子句。

对于具有 $m_G$ 个子句的非平凡叶子，使用
引理~\ref{lem:sparse-checks} 的初始奇偶检验与
引理~\ref{lem:greedy-hashes} 的贪心过程。每个初始行的
支撑大小至多为 $D_0$，每个附加行的列表长度为 $P$。
两个界均为 $\lambda^{O_k(1)}$。由 \eqref{eq:hash-query}，每行
给出一个次数至多 $k$、至多含 $2^k\max\{P,D_0\}$ 个单项式的
方程，且每个被查询的系统至多有 $6n$ 个方程。对充分大的
$n$，单个常数 $c=c(k)$ 同时界定方程与变量之比以及单项式
界中 $\lambda$ 的指数。在每个贪心步骤，使用
引理~\ref{lem:tensor-sparse-query} 对每个候选行 $h_\gamma$ 计算
$N(A;h_\gamma)$，并追加取得最小值的一行。该求值器也适用于
被丢弃的候选，其系统仍接受非满足赋值。最后一次调用给出
所选矩阵的 $N(A)$，由引理~\ref{lem:greedy-hashes} 它等于
叶子计数。

构造初始检验与枚举候选（除求值之外）花费
$2^{\lambda^{O_k(1)}}$ 位运算。每个叶子至多有
$3n(2m_G)^P+1=2^{\lambda^{O_k(1)}}$ 次求值。由于 $\lambda$ 的
每个固定幂次都是 $o(n\log\log n/\sqrt{\log n})$，这些开销
与叶子数目被所声称的时间界
\[
 (n+m)^{O_k(1)}2^{O_k(n\log\log n/\sqrt{\log n})}
\]
吸收。这里 $m$ 计输入子句数。深度优先遍历该树，每次处理
一个查询。各叶子的解集互不相交，因此其精确计数用一个
$(n+1)$ 比特的累加器相加。结合每个查询与每条树路径的
多项式空间，这给出所断言的多项式空间界。
\end{proof}

单个算法同时达到这两个界。取
$\lambda=\lceil\log_2 n\rceil$ 与 $\nu=\lceil\log_2\lambda\rceil$，
当 $m\le n\nu$ 时使用子句网络，否则使用奇偶查询。在第一个
分支中 $n\le s=n+m=O(n\log\log n)$，因此即使在增长的密度下
其指数也是 $O_k(n\log\log n/\sqrt{\log n})$。对每个固定 $D$ 与
所有充分大的 $n$，每个满足 $m\le Dn$ 的输入都走该分支且
$s=O_D(n)$。

\subsection{收缩定理的证明}
\label{app:tensor}

在定理~\ref{thm:tensor-contraction} 的整个证明中，$\log$ 表示以二为底的对数，$\ln$ 表示自然
对数。

\subsubsection{系数环与精确算术}
\label{subsec:tensor-arithmetic}

记 $\lambda=\lceil\log_2 s\rceil$，$e=4\lambda$。
选取首一多项式 $f\in\mathbb Z[U]$，次数为 $e$，其首项系数以下的系数
属于 $\{0,1\}$，且其模二约化不可约。这可以在关于 $s$ 的多项式时间内
确定性地完成：枚举 $2^e$ 个候选，并对每个次数不超过 $e/2$ 的正次数
首一多项式检验 $\mathbb F_2$ 上的整除性。有限域在每个正次数上
都存在不可约多项式~\cite{vzGG13}。
令
\[
 \mathcal O=\mathbb Z_{(2)}[U]/(f),\qquad
 \mathcal K=\mathbb Q[U]/(f),\qquad
 \mathcal A=\mathcal O[V]/(V^\ell),\qquad
 \mathcal A_p=\mathcal A/(2^p).
\]
这里 $\mathbb Z_{(2)}$ 由分母为奇数的有理数组成。模二不可约蕴含
在 $\mathbb Q$ 上不可约，因此 $\mathcal K$ 是域。对 $\mathcal K$ 中的非零
元素，令 $v_2$ 为其在 $1,U,\ldots,U^{e-1}$ 下坐标的通常二进赋值的
最小值。$\mathcal O$ 模二的约化是域 $\mathbb F_{2^e}$。
因此约去最小的坐标赋值即表明
\[
 v_2(xy)=v_2(x)+v_2(y)\qquad(x,y\ne0).
\]
$\mathcal O$ 中任何模二非零的元素都是单位，因为其乘法矩阵的行列式
在 $\mathbb Z_{(2)}$ 中可逆，故逆矩阵的元素属于该环。

$\mathcal A_p$ 中的算术在 $s,p,\ell$ 上花费多项式时间与空间。我们使用
基 $1,U,\ldots,U^{e-1}$ 与 $1,V,\ldots,V^{\ell-1}$，并把每个坐标表示为
零到 $2^p-1$ 之间的整数剩余。标量单位的求逆可以通过模二求逆再作
Newton 提升完成，也可以通过其乘法矩阵完成，所用的是精确的多项式
与整数运算~\cite{vzGG13}。$\mathcal A$ 中元素的可除性与精度始终是
逐坐标的。

$2^e$ 个剩余域元素的二进制坐标提升给出两两差为单位的内插点。
因此在 $\mathcal A_p$ 上，在任意所需数量的这类点上作 Lagrange
插值是合法的。我们将使用少于 $s^4$ 个点。
初始元素属于 $\mathcal B_{P_0,\ell}$ 的收缩在 $\mathcal A_{P_0}$ 中计算后
仍属于该子环，因为它是这些元素的多项式。

当精度提高时，我们逐个提升每个张量元素。对于最初模 $2^p$ 指定的
元素，其\emph{典范提升}是由其所有坐标的整数代表元得到的 $\mathcal A$
中元素。典范提升不必保持加法或乘法。每次递归调用都按其调用的
精度收缩指定的张量。为定义其形变，它首先作这种逐元素的提升
选择。该调用中随后的所有求值都指向这些相同的提升。
更明确地说，每个张量描述都标记有其创建精度 $p_c$（原始张量为 $P_0$）。
在之后每个精度下，其元素都模 $2^{p_c}$ 求值，并使用相同的整数
代表元。精度沿递归路径严格递增，故这些代表元在之后每个精度下
也是典范代表元。特别地，
包张量在整个递归过程中保持其创建时的代表元。

\subsubsection{参数与碰撞打包}
\label{subsec:tensor-packing}

以下所有估计对充分大的 $s$ 成立，其阈值只依赖于固定常数。较小的
有限尺寸范围由直接求和处理。参数不等式可以直接检验。令
\[
 H=\left\lceil\frac{s}{\lfloor\sqrt\lambda\rfloor}\right\rceil,
 \qquad r=64,\qquad \varepsilon=\frac1{128\cdot1024},
 \qquad \Delta=\lceil s^{9/10}\rceil.
\]
边数 $M\le H$ 的调用是叶，直接求值。在深度 $d$ 的活动调用处（根在
深度零），使用
\begin{equation}\label{eq:tensor-parameters}
 u=u_d=\left\lfloor\frac{\varepsilon\lambda}{2^d}\right\rfloor,
 \qquad a=a_d=2\sum_{j=0}^d u_j,\qquad
 a_{-1}=0,\qquad L=2^a,\qquad D_{\rm arr}=r^{2a}.
\end{equation}
我们将安排每个子节点至多有 $M/8$ 条边且最大度数至多 $\Delta$。
这些断言表明，在每条活动路径上，
\[
 H<M\le Cs\,8^{-d},\qquad
 d\le\frac{\log_2\lambda}{6}+O_C(1).
\]
因此
\begin{equation}\label{eq:tensor-parameter-bounds}
 u=\Theta_C(\lambda/2^d),\quad
 \lambda^{4/5}\le u\le M,\quad
 a\le4\varepsilon\lambda,\quad
 L,D_{\rm arr}\le s^{1/1000}
\end{equation}
对所有充分大的 $s$ 成立。特别地，在每个活动调用处
\begin{equation}\label{eq:tensor-packing-conditions}
 M\ge2D_{\rm arr}^{\,2}\Delta,\qquad
 M\ge L D_{\rm arr},\qquad 2L\le\Delta.
\end{equation}
例如，第一个不等式右端为
$O(s^{9/10+2/1000})=o(H)$。其他不等式以同样方式得出。

\begin{lemma}[打包成包]\label{lem:tensor-packets}
在 \eqref{eq:tensor-packing-conditions} 下，用 $[r]^a$ 的元素标记网络每侧的
顶点。存在一个确定性可构造的标记，以及一个把边嵌入到含 $L$ 个槽的
包中的嵌入，具有以下性质。
\begin{enumerate}
\item 至多有 $3M/L$ 个包。每条边在其行端与列端占据相同的槽。
      未使用的槽没有端口。
\item 在每个包和每个坐标中，行符号与列符号的无序对是固定的。
      当这两个符号不同时，槽号的第 $i$ 位指定它们在坐标 $i$ 中的
      取向，其中 $0\le i<a$。
\item 对每个坐标和每个符号，每侧至多有 $3M/r$ 个实端口属于具有
      该符号的顶点。
\end{enumerate}
该构造在 $s+\mathcal L$ 上使用多项式时间与空间。
\end{lemma}
\begin{proof}
对每对标记，令 $X_\gamma$ 为具有该标记对的边数。独立的均匀顶点
标记给出
\[
 \mathbb E\sum_\gamma (X_\gamma-M/D_{\rm arr})^2\le2M\Delta.
\]
事实上，展开把该期望写成有序边对上其标记一致概率减去
$1/D_{\rm arr}$ 的求和。端点不交的边对贡献为零。
至多有 $2M\Delta$ 个有序边对共享一个端点（包括相同的边），
每对至多贡献一。
条件期望给出具有该界的标记。为实现它，对下一个顶点的 $r^a$ 个
标记逐一尝试，并通过枚举边对计算条件期望。每个概率至多涉及
四个顶点，故这是精确的多项式时间计算。这是通常的
条件期望方法~\cite{Raghavan88}。

每个负载至多为 $M/D_{\rm arr}+\sqrt{2M\Delta}$。
把所有标记对类填充到共同的槽数 $H_{\rm pad}$，即最大负载向上取整
到 $L$ 的倍数。假设蕴含
\[
 H_{\rm pad}\le M/D_{\rm arr}+\sqrt{2M\Delta}+L
              \le3M/D_{\rm arr}.
\]
固定每个坐标中符号的无序对。若有 $a'$ 个坐标的符号不相等，
其取向给出 $2^{a'}$ 个标记对。用每个包中来自这些标记的
$2^{a-a'}$ 个槽把它们的填充槽划分为包。
这可行是因为 $H_{\rm pad}$ 可被 $L$ 整除。
把一个包的位置从 $0$ 编号到 $L-1$。其第 $i$ 位在第 $i$ 对不相等时
指定取向，在该对相等时指定副本。在每个包中始终使用一种取向约定。

有 $H_{\rm pad}D_{\rm arr}/L\le3M/L$ 个包。在固定坐标中，给定符号
恰出现在每侧所有填充标记对的 $1/r$ 中。因此其实端口数至多为
$H_{\rm pad}D_{\rm arr}/r\le3M/r$。
\end{proof}

\subsubsection{一个 Cauchy 形变与整局部坐标}
\label{subsec:tensor-deformation}

固定一个活动调用、它的包以及它逐元素提升的张量。使用
\[
 S=\{-L,-L+1,\ldots,L-1\},\qquad
 \mathcal C(T)=\prod_{t\in S}(T-t),\qquad \alpha=2^{a+1}.
\]
在一个包内，对其行槽 $x$ 与列槽 $y$，令
\begin{equation}\label{eq:tensor-cauchy}
 W_{xy}(T)=\frac{\mathcal C(T)}
                 {\mathcal C'(0)(T-(y-x))}.
\end{equation}
分母中的因子与分子中的一个因子相消，因为 $y-x\in S$。
因此这是 $\mathbb Q$ 上次数为 $2L-1$ 的多项式。把不同包之间的
元素置为零，并限制到实端口。在原来的边识别下，
$W(0)=I$，因为对角元素为一，而其他每个元素都有一个未相消的
零因子。

令 $A(Y)$ 与 $B(X)$ 分别为提升后的行张量与列张量在独立端口赋值
$Y,X\in\{0,1\}^M$ 上的乘积。记 $l_0(z)=1-z$，$l_1(z)=z$。定义
\begin{equation}\label{eq:tensor-deformed-contraction}
 F(T)=\sum_{X,Y\in\{0,1\}^M} A(Y)B(X)
                 \prod_{\text{行端口 }x}l_{Y_x}((W(T)X)_x).
\end{equation}
其系数属于 $\mathcal K[V]/(V^\ell)$，次数至多为 $M(2L-1)$，
且 $F(0)$ 是提升后张量的收缩。
特别地，把 $F(0)$ 模该调用的精度约化即给出所需答案。

\begin{lemma}\label{lem:tensor-integrality}
对每个 $t\in S$，$W(t+\alpha z)$ 的元素属于 $\mathcal O[z]$，且
$F(t+\alpha z)$ 属于 $\mathcal A[z]$。
给定 $\zeta$ 模 $2^p$ 的 $e$ 个坐标，每个元素
$W_{xy}(t+\alpha\zeta)\bmod 2^p$ 可用 $s,p$ 上多项式多个位运算算出。
\end{lemma}
\begin{proof}
网格 $S$ 是模 $2^{a+1}$ 的完全剩余系。
对 $S$ 中 $t\ne q$，$v_2(t-q)\le a$，且各赋值上非零差的计数
与 $t$ 无关。因此 $v_2(\mathcal C'(t))=v_2(\mathcal C'(0))$。
在
\[
 \frac{\mathcal C(t+\alpha z)}{\mathcal C'(0)}
  =\frac{\alpha z\prod_{q\in S\setminus\{t\}}
                        (t-q+\alpha z)}{\mathcal C'(0)},
\]
中，每个由 $q\ne t$ 索引的因子都按系数被 $2^{v_2(t-q)}$ 整除。
把这些幂与分母相消后，所示多项式被 $2^{a+1}$ 整除。
为得到 \eqref{eq:tensor-cauchy}，只略去由 $y-x$ 索引的那个因子。
若 $y-x=t$，这去掉 $\alpha z$，留下一个整多项式。
否则 $\alpha z$ 保留，而被略去的因子有二进容度 $2^{v_2(t-y+x)}$，
其中 $v_2(t-y+x)\le a$，故结果仍被二整除。代入
\eqref{eq:tensor-deformed-contraction} 即证得关于 $F$ 的断言。

对求值，先用已知整数 $2^{v_2(t-q)}$ 除每个因子
$t-q+\alpha z$（$q\ne t$）。所得线性多项式的两个系数都是整数。
如上把剩余的显式二的幂与分母相消，并将其奇部模 $2^p$ 求逆。
现在我们可以在给定的 $\zeta$ 剩余处直接计算乘积，包括
$\zeta=0$。输入有 $ep$ 位，乘积有 $O(L)$ 个因子。
所有坐标运算在 $s,p$ 上花费多项式时间。
\end{proof}

\subsubsection{渐缩 Hermite 恢复}
\label{subsec:tensor-hermite}

下一个引理分离出精度估计。其陈述中所有幂级数截断都是 $z$ 中的
通常截断。

\begin{lemma}[渐缩 Hermite 恢复]\label{lem:tapered-hermite}
令整数 $u,M$ 满足 $1\le u\le a$ 与 $M\ge u$，并令
$h=\lceil3M/u\rceil$。
对 $t\in S\setminus\{0\}$ 置
\[
 i=v_2(t),\qquad q_t=q(i)=h\max(0,u-i),\qquad D_t=L/2^i.
\]
假设 $F\in(\mathcal K[V]/(V^\ell))[T]$ 的次数至多为
$M(2L-1)$，$F(0)\in\mathcal A$，且
对每个 $t\in S$ 有 $F(t+\alpha z)\in\mathcal A[z]$。
对 $p\ge1$，其值 $F(0)$ 模 $2^p$ 可由
\[
 F(t+\alpha z)\pmod{2^{p_t}},
 \qquad
 p_t=p+8hD_t+8M(a+1),\qquad q_t>0.
\]
恢复，其中只用到每个多项式的前 $q_t$ 个系数。
恢复在 $M,L,\ell,e$ 与各精度上花费多项式时间与空间。
\end{lemma}
\begin{proof}
$S$ 中有 $L$ 个奇点，故
$\sum_{t\ne0}q_t\ge Lhu\ge3ML>M(2L-1)$。又
$h\le4M/u$，故 $q_t\le4M$。
对每个重数为正的点 $t$ 置
\[
 E_t(T)=\prod_{q\in S\setminus\{0,t\}}
                   \left(\frac{T-q}{t-q}\right)^{q_q}.
\]
Hermite 插值给出
\begin{equation}\label{eq:tensor-hermite-origin}
 F(0)=\sum_{t:q_t>0}E_t(0)
       \left(\left[\frac{F(t+\alpha z)}
                            {E_t(t+\alpha z)}\right]_{<q_t}\right)
                            (-t/\alpha).
\end{equation}
为完整起见，把右端的零换为未定元 $T$，把 $-t/\alpha$ 换为
$(T-t)/\alpha$。每个被加项的次数小于 $\sum q_t$。在其自身节点处
它具有规定的射流，而在其他每个节点处它按所需阶数消失。
因此与 $F$ 的差被
$\prod_{t\ne0}(T-t)^{q_t}$ 整除且次数更小，故为零。
这是 Hermite 插值的通常中国剩余构造 \cite[第~5.6~节]{vzGG13}。

$E_t(t+\alpha z)$ 的每个因子都是 $1+\alpha z/(t-q)$ 的幂。
它是整的且常数项为一。
因此其截断到 $q_t$ 的倒数也是整的，故在
\eqref{eq:tensor-hermite-origin} 中作商不损失精度。

固定 $i=v_2(t)$ 与 $D=D_t$。直接由乘积得
\begin{align}
 v_2(E_t(0))
 &=\sum_{q\ne0,t}q_q\bigl(v_2(q)-v_2(t-q)\bigr)\notag\\
 &=D\sum_{j=1}^{a-i}\frac{j}{2^j}
                          \bigl(q(i+j)-q(i)\bigr)
 \ \ge -6Dh. \label{eq:tensor-hermite-loss}
\end{align}
为说明第二个等式，赋值小于 $i$ 的点贡献为零。对 $1\le j\le a-i$，
有 $D/2^j$ 个点满足 $v_2(q)=i+j$，同样多个点满足
$v_2(t-q)=i+j$。后者满足 $v_2(q)=i$，且这些层穷尽了具有该赋值的
除 $t$ 以外的点。这些计数来自完全剩余网格，其非零整数差的
赋值至多为 $a$。最后，
$q(i+j)-q(i)\ge-hj$ 且 $\sum_{j\ge1}j^2/2^j=6$。

次数低于 $q_t$ 的多项式在 $-t/\alpha$ 处求值至多再损失
$4M(a+1)$ 个二的幂。因此把每个商系数替换为其模
$2^{p_t}$ 剩余的任何提升，会使其在
\eqref{eq:tensor-hermite-origin} 中的贡献改变 $2^p\mathcal A$ 中的
一个元素。把这些误差求和不再损失精度。重构的和是整的，因为它与
整值 $F(0)$ 相差 $2^p\mathcal A$ 中的一个元素。

可以在指定基中用有理算术完成最后的组合。乘子
$E_t(0)(-t/\alpha)^j$（$j<q_t$）是多项式多个 $O(a)$ 位有理因子的
乘积（计重数）。因此其分子与分母长度在 $M,L,a$ 中为多项式。
这样的系数只有多项式多个，故即使取其分母的乘积也构成合适的
多项式位公共分母。相加后，约去其二的幂，并把所得整坐标模 $2^p$
约化。这就证明了算法断言。
\end{proof}

在输入精度为 $p$ 的调用处，用 \eqref{eq:tensor-parameters} 中的参数
应用此引理。对每个活动的 $t$，
在 $M(2L-1)+1$ 个不同的剩余域元素的二进制坐标提升处求值并作 Lagrange
插值，得到整个多项式 $F(t+\alpha z)$ 模 $2^{p_t}$。
引理~\ref{lem:tensor-integrality} 允许在 $\mathcal A_{p_t}$ 上作此
插值。由 \eqref{eq:tensor-parameter-bounds}，点数少于 $s^4$。
下面的求值算法把精度 $p_t$ 传给其子节点。

输入精度与子精度满足
\begin{equation}\label{eq:tensor-precision-bound}
 p\le s\lambda^{O(1)}2^{a_{d-1}},
 \qquad
 p_t=O_C(MD_t/u).
\end{equation}
第一个断言在根处成立，因为 $P_0\le Cs$。沿一条路径，精度的增量为
$8hD_t+8M(a+1)$，其中 $h\le4M/u$、$M\le Cs$、$D_t\le2^{a_d}$。
把它们在 $O(\log\lambda)$ 个深度上求和，即证得对每个调用（包括叶）
的第一个断言。对第二个断言，$q_t>0$ 蕴含 $i<u$，从而
\[
 D_t\ge2^{a_d-u}=2^{a_{d-1}+u}.
\]
由于 $M>H$ 且 $u\ge\lambda^{4/5}$，这个额外因子 $2^u$ 支配 $\lambda$
的每个固定幂。因此关于 $p$ 的第一个界以及 $8M(a+1)$ 都是
$O_C(MD_t/u)$。剩余的增量 $8hD_t$ 直接具有该界。
所有精度在 $s$ 中保持多项式，因为
$a_d\le4\varepsilon\lambda$。

渐缩解释了此估计中的因子 $1/u$。
作为比较，给 $u$ 以下的每个层级指定相同的重数 $q_*=\Theta(M)$，
其余层级指定为零。在最后保留的层级 $i=u-1$ 处，恒等式
\eqref{eq:tensor-hermite-loss} 会给出损失
\[
 -q_*D_t\sum_{j=1}^{a-i}\frac{j}{2^j}=-\Theta(MD_t).
\]
而我们的重数每层只改变 $h=\Theta(M/u)$，
给出损失 $O(MD_t/u)$。$a_d$ 的累积选择提供了吸收继承精度与
求值损失所需的额外因子 $2^u$。因此在每个保留层级与每个递归
深度处，较小的界都成立。

\subsubsection{群之间相互作用的模式}
\label{subsec:tensor-modes}

固定一个活动的 $t$ 与 $z$ 的一个插值 $\zeta\in\mathcal O$。在本小节中
$T=t+\alpha\zeta$，所有计算都模 $2^{p_t}$。置 $i=v_2(t)<u$ 与
$D=L/2^i$。
按顶点在坐标 $i$ 中的符号把两侧的顶点分组。
$W(T)$ 中两个端口属于同一群的元素称为局部的，并保持不变。

这种坐标选择使 $W(t)$ 成为局部的。事实上，
除非 $y-x=t$，$W_{xy}(t)=0$。加上一个赋值为 $i$ 的整数会翻转第 $i$ 位，
故存活的边对具有不同的第 $i$ 位。当包在该坐标中的两个符号不同时，
这意味着行符号与列符号一致。相等的符号对
本来就位于一个群中。从 $t$ 移到 $t+\alpha\zeta$ 会引入
我们现在展开的交叉项。

在坐标 $i$ 中两个符号不同的包里，行位置 $x$ 与列位置 $y$ 位于
不同群中，恰好当它们的第 $i$ 位相等。把列划分为瓦片
\[
 y=y_0+2^{i+1}j,\qquad 0\le y_0<2^{i+1}.
\]
列群在一个瓦片上为常数。对一条交叉到该瓦片的行，
\[
 v_2(t-y_0+x)\le i.
\]
否则 $y_0$ 会与 $x+2^i$ 模 $2^{i+1}$ 一致，与其第 $i$ 位相等矛盾。
加上 $\alpha\zeta$ 不改变该赋值。因此几何展开给出
\begin{equation}\label{eq:tensor-mode-expansion}
 W_{xy}(T)=
 \sum_{\substack{b\ge0\\b+1<p_t}}
 \frac{\mathcal C(T)}
      {\mathcal C'(0)(T-y_0+x)}
 \left(\frac{2^{i+1}}{T-y_0+x}\right)^b j^b
 \pmod {2^{p_t}}.
\end{equation}
事实上第一个分式的赋值至少为 $a+1-i\ge1$，而该比在取 $b$ 次幂
之前的赋值至少为一。几何余项被 $2^{p_t}$ 整除。这些表达式可以
通过分解已知的二的幂并只用单位逆来计算。
对于每个这样的包、瓦片和 $b$，引入一个权为 $w_Z=b+1$ 的模式变量 $Z$，并以
\[
 L_Z(X)=\sum_{\text{瓦片中的实列 }y}j^bX_y.
\]
替换它。按惯例 $j^0=1$，包括 $j=0$ 的情形。其源是该瓦片的列组，其目标是与该包坐标对配对的另一个组。公式
\eqref{eq:tensor-mode-expansion} 将每个行形式表示为其局部部分与其目标模式的线性组合之和。
$Z$ 的系数可被 $2^{w_Z}$ 整除。任一权重的变量至多有
\begin{equation}\label{eq:tensor-mode-multiplicity}
 (3M/L)2^{i+1}=6M/D
\end{equation}
个。

在仿射因子
$l_{Y_x}((W(T)X)_x)$ 的乘积中，每个模式单项式
$Z^\beta$ 的系数都可被 $2^{\sum_Zw_Z\beta_Z}$ 整除。
因此只需考虑满足
\begin{equation}\label{eq:tensor-mode-budget}
 \sum_Zw_Z\beta_Z<p_t
\end{equation}
的指数元组。这些元组的数目为
\begin{equation}\label{eq:tensor-partition-count}
 \exp\bigl(O(\sqrt{Mp_t/D})\bigr)
       =\exp\bigl(O_C(M/\sqrt u)\bigr).
\end{equation}
以下是完整的计数估计。对任意 $\theta>0$，
生成乘积与 \eqref{eq:tensor-mode-multiplicity} 将该数目界为
\[
 \exp\left(\theta p_t+\frac{6M}{D}
                   \sum_{j\ge1}-\ln(1-e^{-\theta j})\right).
\]
展开对数并对几何级数求和得
\[
 \sum_{j\ge1}-\ln(1-e^{-\theta j})
  =\sum_{h\ge1}\frac1{h(e^{\theta h}-1)}
  \le\frac1\theta\sum_{h\ge1}\frac1{h^2}.
\]
取 $\theta=\sqrt{M/(Dp_t)}$ 并使用
\eqref{eq:tensor-precision-bound}。所有模式变量及其
指数都有多项式规模的描述。满足
\eqref{eq:tensor-mode-budget} 的元组可按变量排序，并在剩余预算的约束下依次选取
每个指数，从而以多项式延迟和空间枚举。每个可行的部分
选择都有全零的延续，因此访问的部分选择数至多是完整元组数
乘以变量数，至多相差一个加性常数。

\subsubsection{子网络及其描述}
\label{subsec:tensor-children}

对于 \eqref{eq:tensor-mode-budget} 中的一个元组 $\beta$，为每个组 $g\in[r]$ 构造
一个子网络。把当前网络中来自两侧的所有当前顶点都放在
子网络的一侧。保留它们的端口及其选定的提升项，此时已对 $2^{p_t}$ 取模。
在另一侧，为每个包放置一个顶点，并向该包每个保留下来的端口
连一条边。其指标恰为该包中组 $g$ 的行位 $Y$ 与列位 $X$。

为定义这个新的包张量，取包中组 $g$ 的行因子 $l_{Y_x}$ 的乘积，其中各因子
的自变量用局部列位与形式模式表示。对所有以 $g$ 为目标的模式，提取指数为
$\beta_Z$ 的系数。将其乘以
\[
 \prod_{\substack{Z\text{ 在该包中}\\
                  \text{source}(Z)=g}}L_Z(X)^{\beta_Z}.
\]
空乘积等于一。若没有目标行，则只要某个目标指数为正，提取出的
系数就是零。若 $Z$ 的源瓦片没有实列，则 $L_Z(X)=0$，其
幂在 $\beta_Z=0$ 时贡献一，否则贡献零。

每个模式恰有一个目标和一个源。因此系数提取可沿包-组行乘积分解，
且每个模式的代换只依赖于其源组中的位。将
\eqref{eq:tensor-deformed-contraction} 中的乘积展开，再对组内相互独立的位求和，
即证明精确恒等式
\begin{equation}\label{eq:tensor-recursion-identity}
 F(t+\alpha\zeta)
   =\sum_{\beta:\,\sum w_Z\beta_Z<p_t}
                       \prod_{g=1}^r Z_{g,\beta}
                  \pmod {2^{p_t}},
\end{equation}
其中 $Z_{g,\beta}$ 是刚才定义的子网络的收缩。被省略的单项式由其系数整除性而消失。

由引理~\ref{lem:tensor-packets}，一个子网络至多有
$6M/r=3M/32<M/8$ 条边。旧顶点的度不变，而
新的包顶点至多有 $2L\le\Delta$ 个端口。因此
\eqref{eq:tensor-parameter-bounds} 中使用的不变量归纳成立。

所有新张量都是隐式给出的。要计算某个包张量的一项，展开其至多 $L$ 个仿射行
因子的乘积。每个因子关于 $s$ 的项数都是多项式的，包括其局部项。
枚举这些选择、检查所需指数并对所得各项求和，需要
\[
 \exp(O(L\log s))
\]
时间和多项式空间。因子 $L_Z(X)^{\beta_Z}$ 通过在 $\mathcal A_{p_t}$ 中通常的反复平方
计算。由于 $L\le s^{1/1000}$，这一项的计算代价为
$\exp(O(s^{9/10}))$。

新张量的描述由其包数据、元组 $\beta$、插值取值和创建精度
组成。这些数据连同局部核足以直接计算其项。保留下来的张量
按其原始描述计算；若它是在递归中更早创建的，则按其自身的
创建时描述计算。其在该精度下的剩余类随后被规范地
提升到当前精度。旧描述是不可变的，并通过引用共享。因此项计算
例程的任何复合都不会随深度累积，且所有项计算例程
仍具有所需的时间与空间界。
沿一条路径只添加多项式多个顶点和多项式多的新
描述数据。

\subsubsection{时间与空间界}
\label{subsec:tensor-cost}

\begin{proof}[定理~\ref{thm:tensor-contraction} 的证明]
将原始项嵌入 $\mathcal A_{P_0}$ 并执行上述递归算法。在一个活动调用处，对每个
重数为正的节点和每个所需的剩余类互异插值取值，通过其子收缩计算 \eqref{eq:tensor-recursion-identity}。
对 $F(t+\alpha z)$ 插值并应用
引理~\ref{lem:tapered-hermite}。在叶节点处，直接对所有
布尔边赋值求和。引理~\ref{lem:tensor-integrality}、
递归恒等式与恢复引理从叶节点出发归纳证明正确性。

一个调用处的插值节点数、取值数和组数关于 $s$ 是多项式的，且与 $\mathcal L$ 无关。
由 \eqref{eq:tensor-partition-count}，一个活动深度 $d$ 调用的子节点数至多为
\[
 s^{O(1)}\exp\left(
       O_C\left(\frac{s}{\sqrt\lambda}
                    \left(\frac{\sqrt2}{8}\right)^d\right)\right).
\]
事实上 $M\le Cs\,8^{-d}$ 且
$u_d=\Theta_C(\lambda/2^d)$。将这些界在
$O_C(\log\lambda)$ 个深度上相乘，或以这种方式界定每一层上的节点数，得到总调用数至多为
\[
 \exp\left(O_C(s/\sqrt\lambda)
                   +O_C(\lambda\log\lambda)\right)
       =\exp(O_C(s/\sqrt\lambda))
\]
次。在叶节点处，直接求和至多使用 $2^H$ 个赋值，且每个赋值的项计算次数为多项式。
其代价被 $s+\mathcal L$ 的多项式乘以
\[
 \exp(O_C(H)+O_{C,A}(s^{9/10}))
   =\exp(O_{C,A}(s/\sqrt\lambda)).
\]
所界定。局部算术、打包、描述和恢复，除已界定的项计算之外，每生成一项
只使用多项式时间。依赖于描述的因子界定每次调用的工作量。由于分支界与
$\mathcal L$ 无关，总工作量保持 $s+\mathcal L$ 的固定多项式因子。
这证明了时间界。

关于空间，深度优先遍历递归，每次累加一个关于 $\beta$ 的和与一个关于组的乘积。
存储一个调用恢复所需的多项式多个 Taylor 数据、其当前的插值与模式指标，以及其
图与新张量描述。它们的大小关于 $s+\mathcal L$ 是多项式的，其中精度由
\eqref{eq:tensor-precision-bound} 界定。栈的层数只有
$O_C(\log\lambda)$。引理~\ref{lem:tapered-hermite} 中的公分母
计算也只使用多项式位整数。每个调用在其自身固定精度下累加其收缩项。
原始项计算例程按假设使用多项式空间，
新例程具有同样的性质。
这证明了空间界。最后在 $\mathcal A_{P_0}$ 中提取基
坐标，它们给出 $\mathcal B_{P_0,\ell}$ 中的答案。
\end{proof}


\begin{thebibliography}{CHKX06}

\bibitem[ABHS19]{ABHS19}
Amir Abboud, Karl Bringmann, Danny Hermelin, and Dvir Shabtay.
\href{https://doi.org/10.1137/1.9781611975482.3}
{\emph{SETH-Based Lower Bounds for Subset Sum and Bicriteria Path}}.
30th ACM--SIAM Symposium on Discrete Algorithms, 2019, 41--57.

\bibitem[AW14]{AW14}
Amir Abboud and Virginia Vassilevska Williams.
\href{https://doi.org/10.1109/FOCS.2014.53}
{\emph{Popular Conjectures Imply Strong Lower Bounds for Dynamic Problems}}.
55th IEEE Symposium on Foundations of Computer Science, 2014, 434--443.

\bibitem[AGP23]{AGP23}
Josh Alman, Yunfeng Guan, and Ashwin Padaki.
\href{https://arxiv.org/abs/2211.05217}
{\emph{Smaller Low-Depth Circuits for Kronecker Powers}}.
34th ACM--SIAM Symposium on Discrete Algorithms, 2023, 4159--4187.

\bibitem[AW15]{AW15}
Josh Alman and Ryan Williams.
\href{https://arxiv.org/abs/1507.05106}
{\emph{Probabilistic Polynomials and Hamming Nearest Neighbors}}.
56th IEEE Symposium on Foundations of Computer Science, 2015, 136--150.

\bibitem[AW17]{AW17}
Josh Alman and Ryan Williams.
\href{https://arxiv.org/abs/1611.05558}
{\emph{Probabilistic Rank and Matrix Rigidity}}.
49th ACM Symposium on Theory of Computing, 2017, 641--652.

\bibitem[AR94]{AR94}
Noga Alon and Yuval Roichman.
\href{https://doi.org/10.1002/rsa.3240050203}
{\emph{Random Cayley Graphs and Expanders}}.
Random Structures \& Algorithms \textbf{5}(2), 1994, 271--284.

\bibitem[BI16]{BI16}
Arturs Backurs and Piotr Indyk.
\href{https://doi.org/10.1109/FOCS.2016.56}
{\emph{Which Regular Expression Patterns Are Hard to Match?}}
57th IEEE Symposium on Foundations of Computer Science, 2016, 457--466.
\bibitem[BI18]{BackursIndyk18}
Arturs Backurs and Piotr Indyk.
\href{https://doi.org/10.1137/15M1053128}
{\emph{Edit Distance Cannot Be Computed in Strongly Subquadratic Time
(unless SETH is false)}}.
SIAM Journal on Computing \textbf{47}(3), 2018, 1087--1097.
会议版本见 STOC 2015。

\bibitem[Bar23]{Bar23}
Boaz Barak.
\href{https://introtcs.org/public/lec_13_Cook_Levin.html}
{\emph{Introduction to Theoretical Computer Science}}.
在线教材，2023。第~15 章，引理~15.9--15.10。
\href{https://github.com/boazbk/tcs/releases/tag/v0.95}{存档的 2023 版（v0.95）}。

\bibitem[Bar68]{Bareiss68}
Erwin H. Bareiss.
\href{https://doi.org/10.1090/S0025-5718-1968-0226829-0}
{\emph{Sylvester's Identity and Multistep Integer-Preserving Gaussian Elimination}}.
Mathematics of Computation \textbf{22}(103), 1968, 565--578.

\bibitem[BI13]{BeckImp13}
Chris Beck and Russell Impagliazzo.
\href{https://doi.org/10.1145/2488608.2488669}
{\emph{Strong ETH Holds for Regular Resolution}}.
45th ACM Symposium on Theory of Computing, 2013, 487--494.

\bibitem[BT94]{BT94}
Richard Beigel and Jun Tarui.
\href{https://doi.org/10.1007/BF01263423}
{\emph{On ACC}}.
Computational Complexity \textbf{4}(4), 1994, 350--366.

\bibitem[BGKM23]{BGKM23}
Vishwas Bhargava, Sumanta Ghosh, Mrinal Kumar, and Chandra Kanta Mohapatra.
\href{https://doi.org/10.1145/3625226}
{\emph{Fast, Algebraic Multivariate Multipoint Evaluation in Small
Characteristic and Applications}}.
Journal of the ACM \textbf{70}(6), 2023, 42:1--42:46.
会议版本见 STOC 2022。

\bibitem[Bj\"o16]{Bjo16}
Andreas Bj\"orklund.
\href{https://doi.org/10.4230/LIPIcs.SWAT.2016.17}
{\emph{Below All Subsets for Some Permutational Counting Problems}}.
15th Scandinavian Symposium and Workshops on Algorithm Theory,
LIPIcs \textbf{53}, 2016, 17:1--17:11.

\bibitem[BKW19]{BKW19}
Andreas Bj\"orklund, Petteri Kaski, and Ryan Williams.
\href{https://doi.org/10.1007/s00453-018-0513-7}
{\emph{Generalized Kakeya Sets for Polynomial Evaluation and Faster
Computation of Fermionants}}.
Algorithmica \textbf{81}, 2019, 4010--4028.
会议版本见 IPEC 2017。

\bibitem[Bri14]{Bri14}
Karl Bringmann.
\href{https://doi.org/10.1109/FOCS.2014.76}
{\emph{Why Walking the Dog Takes Time: Fr\'echet Distance Has No Strongly
Subquadratic Algorithms Unless SETH Fails}}.
55th IEEE Symposium on Foundations of Computer Science, 2014, 661--670.

\bibitem[BK15]{BK15}
Karl Bringmann and Marvin K\"unnemann.
\href{https://doi.org/10.1109/FOCS.2015.15}
{\emph{Quadratic Conditional Lower Bounds for String Problems and
Dynamic Time Warping}}.
56th IEEE Symposium on Foundations of Computer Science, 2015, 79--97.

\bibitem[Cal09]{CalabroThesis}
Chris Calabro.
\href{https://cseweb.ucsd.edu/~ccalabro/dissertation_single_spaced.pdf}
{\emph{The Exponential Complexity of Satisfiability Problems}}.
Ph.D. dissertation, University of California, San Diego, 2009.
第~5 章。

\bibitem[CIP06]{CIP}
Chris Calabro, Russell Impagliazzo, and Ramamohan Paturi.
\href{https://cseweb.ucsd.edu/~paturi/myPapers/pubs/CalabroImpagliazzoPaturi_2006_ccc.pdf}
{\emph{A Duality between Clause Width and Clause Density for SAT}}.
21st IEEE Conference on Computational Complexity, 2006, 252--260.

\bibitem[CGI$^{+}$16]{CGIMPS16}
Marco L. Carmosino, Jiawei Gao, Russell Impagliazzo, Ivan Mihajlin,
Ramamohan Paturi, and Stefan Schneider.
\href{https://doi.org/10.1145/2840728.2840746}
{\emph{Nondeterministic Extensions of the Strong Exponential Time
Hypothesis and Consequences for Non-reducibility}}.
7th Innovations in Theoretical Computer Science Conference, 2016, 261--270.

\bibitem[CCK$^{+}$20]{CCK20}
Parinya Chalermsook, Marek Cygan, Guy Kortsarz, Bundit Laekhanukit,
Pasin Manurangsi, Danupon Nanongkai, and Luca Trevisan.
\href{https://doi.org/10.1137/18M1166869}
{\emph{From Gap-Exponential Time Hypothesis to Fixed Parameter Tractable
Inapproximability: Clique, Dominating Set, and More}}.
SIAM Journal on Computing \textbf{49}(4), 2020, 772--810.
会议版本见 FOCS 2017。

\bibitem[CW21]{CW21}
Timothy M. Chan and R. Ryan Williams.
\href{https://doi.org/10.1145/3402926}
{\emph{Deterministic APSP, Orthogonal Vectors, and More:
Quickly Derandomizing Razborov--Smolensky}}.
ACM Transactions on Algorithms \textbf{17}(1), Article~2, 2021, 1--14.
\href{https://people.csail.mit.edu/rrw/derand-apsp-ov-cr-talg.pdf}{作者手稿}，定理~1.6 与~2.5。

\bibitem[CHKX06]{CHKX06}
Jianer Chen, Xiuzhen Huang, Iyad A. Kanj, and Ge Xia.
\href{https://doi.org/10.1016/j.jcss.2006.04.007}
{\emph{Strong Computational Lower Bounds via Parameterized Complexity}}.
Journal of Computer and System Sciences \textbf{72}(8), 2006, 1346--1367.

\bibitem[CLWYY25]{CLWYY25}
Zongchen Chen, Aditya Lonkar, Chunyang Wang, Kuan Yang, and Yitong Yin.
\href{https://doi.org/10.1145/3717823.3718163}
{\emph{Counting Random $k$-SAT near the Satisfiability Threshold}}.
57th ACM Symposium on Theory of Computing, 2025, 867--878.
定理~1.2 的编号依照\href{https://wcysai.com/publication/randomksat/randomkSAT.pdf}{作者的完整版本}。

\bibitem[Coo71]{Cook71}
Stephen A. Cook.
\href{https://www.cs.utoronto.ca/~sacook/homepage/1971.pdf}
{\emph{The Complexity of Theorem-Proving Procedures}}.
3rd ACM Symposium on Theory of Computing, 1971, 151--158.

\bibitem[Cur18]{Cur18}
Radu Curticapean.
\href{https://doi.org/10.1016/j.ic.2018.02.008}
{\emph{Block Interpolation: A Framework for Tight Exponential-Time
Counting Complexity}}.
Information and Computation \textbf{261}, Part~2, 2018, 265--280.
会议版本见 ICALP 2015。

\bibitem[CDL$^{+}$16]{CDL16}
Marek Cygan, Holger Dell, Daniel Lokshtanov, D\'aniel Marx,
Jesper Nederlof, Yoshio Okamoto, Ramamohan Paturi, Saket Saurabh,
and Magnus Wahlstr\"om.
\href{https://doi.org/10.1145/2925416}
{\emph{On Problems as Hard as CNF-SAT}}.
ACM Transactions on Algorithms \textbf{12}(3), Article~41, 2016, 1--24.

\bibitem[DHM$^{+}$14]{DHMTW}
Holger Dell, Thore Husfeldt, D\'aniel Marx, Nina Taslaman, and Martin Wahl\'en.
\href{https://arxiv.org/abs/1206.1775}
{\emph{Exponential Time Complexity of the Permanent and the Tutte Polynomial}}.
ACM Transactions on Algorithms \textbf{10}(4), Article~21, 2014, 1--32.

\bibitem[Din16]{Din16}
Irit Dinur.
\href{https://eccc.weizmann.ac.il/report/2016/128/}
{\emph{Mildly Exponential Reduction from Gap 3SAT to Polynomial-Gap
Label-Cover}}.
Electronic Colloquium on Computational Complexity, Report TR16-128, 2016.

\bibitem[DL20]{DL20}
Zeev Dvir and Allen Liu.
\href{https://theoryofcomputing.org/articles/v016a020/}
{\emph{Fourier and Circulant Matrices are Not Rigid}}.
Theory of Computing \textbf{16}(20), 2020, 1--48.

\bibitem[vzGG13]{vzGG13}
Joachim von zur Gathen and J\"urgen Gerhard.
\href{https://doi.org/10.1017/CBO9781139856065}{\emph{Modern Computer Algebra}}.
3rd edition, Cambridge University Press, 2013.

\bibitem[Gol00]{Gol00}
Oded Goldreich.
\href{https://eccc.weizmann.ac.il/eccc-reports/2000/TR00-090/index.html}
{\emph{Candidate One-Way Functions Based on Expander Graphs}}.
Electronic Colloquium on Computational Complexity, Report TR00-090, 2000.

\bibitem[Gol08]{Gol08}
Oded Goldreich.
\href{https://www.wisdom.weizmann.ac.il/~oded/CC/bc-1.pdf}
{\emph{P, NP, and NP-Completeness: The Basics of Computational Complexity}}.
作者草稿，2008 年 6 月 14 日，命题~4.7。
Cambridge University Press 版于 2010 年出版。

\bibitem[HKZZ19]{HKZZ19}
Thomas Dueholm Hansen, Haim Kaplan, Or Zamir, and Uri Zwick.
\href{https://doi.org/10.1145/3313276.3316359}
{\emph{Faster $k$-SAT Algorithms Using Biased-PPSZ}}.
51st ACM Symposium on Theory of Computing, 2019, 578--589.

\bibitem[Her14]{Hertli14}
Timon Hertli.
\href{https://doi.org/10.1137/120868177}
{\emph{3-SAT Faster and Simpler---Unique-SAT Bounds for PPSZ Hold in General}}.
SIAM Journal on Computing \textbf{43}(2), 2014, 718--729.
会议版本见 FOCS 2011。

\bibitem[Hus11]{Hus11}
Thore Husfeldt.
\href{https://doi.org/10.1007/978-3-642-22012-8_3}
{\emph{Invitation to Algorithmic Uses of Inclusion--Exclusion}}.
38th International Colloquium on Automata, Languages, and Programming,
Part II, Lecture Notes in Computer Science \textbf{6756}, 2011, 42--59.

\bibitem[IMP12]{IMP12}
Russell Impagliazzo, William Matthews, and Ramamohan Paturi.
\href{https://cseweb.ucsd.edu/~paturi/myPapers/pubs/ImpagliazzoMatthewsPaturi_2012_soda.pdf}
{\emph{A Satisfiability Algorithm for $AC^0$}}.
23rd ACM--SIAM Symposium on Discrete Algorithms, 2012, 961--972.
\href{https://doi.org/10.1137/1.9781611973099.77}{出版方版本}。

\bibitem[IP01]{IP01}
Russell Impagliazzo and Ramamohan Paturi.
\href{https://doi.org/10.1006/jcss.2000.1727}
{\emph{On the Complexity of $k$-SAT}}.
Journal of Computer and System Sciences \textbf{62}(2), 2001, 367--375.

\bibitem[IPZ01]{IPZ}
Russell Impagliazzo, Ramamohan Paturi, and Francis Zane.
\href{https://doi.org/10.1006/jcss.2001.1774}
{\emph{Which Problems Have Strongly Exponential Complexity?}}
Journal of Computer and System Sciences \textbf{63}(4), 2001, 512--530.

\bibitem[JMV13]{JMV13}
Hamid Jahanjou, Eric Miles, and Emanuele Viola.
\href{https://eccc.weizmann.ac.il/report/2013/099/revision/2/download/}
{\emph{Local Reductions}}.
Electronic Colloquium on Computational Complexity, Report TR13-099,
revision 2, 2013. 会议版本见 ICALP 2015，749--760。

\bibitem[JVV86]{JVV86}
Mark R. Jerrum, Leslie G. Valiant, and Vijay V. Vazirani.
\href{https://doi.org/10.1016/0304-3975(86)90174-X}
{\emph{Random Generation of Combinatorial Structures from a Uniform Distribution}}.
Theoretical Computer Science \textbf{43}(2--3), 1986, 169--188.

\bibitem[JC26]{JC26}
Tao Jiang and Shaowei Cai.
\href{https://arxiv.org/abs/2607.10697}
{\emph{A Better Analysis For PPSZ For 3-SAT}}.
arXiv:2607.10697, 2026.

\bibitem[Koi06]{Koi06}
Mikko Koivisto.
\href{https://doi.org/10.1016/j.ipl.2005.11.013}
{\emph{Optimal 2-Constraint Satisfaction via Sum--Product Algorithms}}.
Information Processing Letters \textbf{98}(1), 2006, 24--28.

\bibitem[KSY14]{KSY14}
Swastik Kopparty, Shubhangi Saraf, and Sergey Yekhanin.
\href{https://doi.org/10.1145/2629416}
{\emph{High-Rate Codes with Sublinear-Time Decoding}}.
Journal of the ACM \textbf{61}(5), 2014, 28:1--28:20.

\bibitem[KW16]{KW16}
Ioannis Koutis and Ryan Williams.
\href{https://doi.org/10.1145/2742544}
{\emph{Algebraic Fingerprints for Faster Algorithms}}.
Communications of the ACM \textbf{59}(1), 2016, 98--105.

\bibitem[Kut26]{Kutzkov26}
Konstantin Kutzkov.
\href{https://doi.org/10.24963/ijcai.2026/254}
{\emph{Complementary Branching for \#SAT: Bounded-Degree Bounds and Faster \#3-SAT}}.
35th International Joint Conference on Artificial Intelligence, 2026, 2282--2290.

\bibitem[Li21]{Li21}
Ray Li.
\href{https://arxiv.org/abs/2008.05106}
{\emph{Settling SETH vs. Approximate Sparse Directed Unweighted Diameter
(up to (NU)NSETH)}}.
53rd ACM Symposium on Theory of Computing, 2021, 1684--1696.

\bibitem[Liu18]{Liu18}
Sixue Liu.
\href{https://doi.org/10.4230/LIPIcs.ICALP.2018.88}
{\emph{Chain, Generalization of Covering Code, and Deterministic Algorithm for $k$-SAT}}.
45th International Colloquium on Automata, Languages, and Programming,
LIPIcs \textbf{107}, 2018, 88:1--88:13.

\bibitem[LMS18]{LMS18}
Daniel Lokshtanov, D\'aniel Marx, and Saket Saurabh.
\href{https://doi.org/10.1145/3170442}
{\emph{Known Algorithms on Graphs of Bounded Treewidth Are Probably Optimal}}.
ACM Transactions on Algorithms \textbf{14}(2), Article~13, 2018, 1--30.
会议版本见 SODA 2011。

\bibitem[MR17]{MR17}
Pasin Manurangsi and Prasad Raghavendra.
\href{https://doi.org/10.4230/LIPIcs.ICALP.2017.78}
{\emph{A Birthday Repetition Theorem and Complexity of Approximating
Dense CSPs}}.
44th International Colloquium on Automata, Languages, and Programming,
LIPIcs \textbf{80}, 2017, 78:1--78:15.

\bibitem[Mar10]{Marx10}
D\'aniel Marx.
\href{https://doi.org/10.4086/toc.2010.v006a005}
{\emph{Can You Beat Treewidth?}}
Theory of Computing \textbf{6}, 2010, 85--112.

\bibitem[NN93]{NN93}
Joseph Naor and Moni Naor.
\href{https://doi.org/10.1137/0222053}
{\emph{Small-Bias Probability Spaces: Efficient Constructions and Applications}}.
SIAM Journal on Computing \textbf{22}(4), 1993, 838--856.

\bibitem[Ope26]{OpenAIUGC26}
OpenAI, \emph{The Unique Games Theorem},
OpenAI Math Release preprint
\href{https://github.com/openai/math/blob/main/preprints/The-Unique-Games-Theorem-September-23-2026/paper.pdf}{OAI:The-Unique-Games-Theorem-September-23-2026},
2026.

\bibitem[PPSZ05]{PPSZ05}
Ramamohan Paturi, Pavel Pudl\'ak, Michael E. Saks, and Francis Zane.
\href{https://doi.org/10.1145/1066100.1066101}
{\emph{An Improved Exponential-Time Algorithm for $k$-SAT}}.
Journal of the ACM \textbf{52}(3), 2005, 337--364.
会议版本见 FOCS 1998。

\bibitem[Rag88]{Raghavan88}
Prabhakar Raghavan.
\href{https://doi.org/10.1016/0022-0000(88)90003-7}
{\emph{Probabilistic Construction of Deterministic Algorithms:
Approximating Packing Integer Programs}}.
Journal of Computer and System Sciences \textbf{37}(2), 1988, 130--143.

\bibitem[Rag08]{Rag08}
Prasad Raghavendra.
\href{https://doi.org/10.1145/1374376.1374414}
{\emph{Optimal Algorithms and Inapproximability Results for Every CSP?}}
40th ACM Symposium on Theory of Computing, 2008, 245--254.

\bibitem[RV13]{RV13}
Liam Roditty and Virginia Vassilevska Williams.
\href{https://doi.org/10.1145/2488608.2488673}
{\emph{Fast Approximation Algorithms for the Diameter and Radius of Sparse Graphs}}.
45th ACM Symposium on Theory of Computing, 2013, 515--524.

\bibitem[Sch24]{Scheder24}
Dominik Scheder.
\href{https://doi.org/10.46298/theoretics.24.5}
{\emph{PPSZ is better than you think}}.
TheoretiCS \textbf{3}, Article~5, 2024, 1--37.

\bibitem[Sch99]{Schon99}
Uwe Sch\"oning.
\href{https://doi.org/10.1109/SFFCS.1999.814612}
{\emph{A Probabilistic Algorithm for $k$-SAT and Constraint Satisfaction
Problems}}.
40th IEEE Symposium on Foundations of Computer Science, 1999, 410--414.

\bibitem[SS09]{SS09}
Alexander D. Scott and Gregory B. Sorkin.
\href{https://doi.org/10.1145/1597036.1597049}
{\emph{Polynomial Constraint Satisfaction Problems, Graph Bisection,
and the Ising Partition Function}}.
ACM Transactions on Algorithms \textbf{5}(4), 2009, 45:1--45:27.

\bibitem[SKC94]{SKC94}
Bart Selman, Henry A. Kautz, and Bram Cohen.
\href{https://cdn.aaai.org/AAAI/1994/AAAI94-051.pdf}
{\emph{Noise Strategies for Improving Local Search}}.
12th National Conference on Artificial Intelligence, 1994, 337--343.

\bibitem[SS96]{SS96}
Michael Sipser and Daniel A. Spielman.
\href{https://doi.org/10.1109/18.556667}
{\emph{Expander Codes}}.
IEEE Transactions on Information Theory \textbf{42}(6), 1996, 1710--1722.

\bibitem[Tan81]{Tanner81}
R. Michael Tanner.
\href{https://doi.org/10.1109/TIT.1981.1056404}
{\emph{A Recursive Approach to Low Complexity Codes}}.
IEEE Transactions on Information Theory \textbf{27}(5), 1981, 533--547.

\bibitem[Tov84]{Tov84}
Craig A. Tovey.
\href{https://doi.org/10.1016/0166-218X(84)90081-7}
{\emph{A Simplified NP-Complete Satisfiability Problem}}.
Discrete Applied Mathematics \textbf{8}(1), 1984, 85--89.

\bibitem[Val77]{Val77}
Leslie G. Valiant.
\href{https://doi.org/10.1007/3-540-08353-7_135}
{\emph{Graph-Theoretic Arguments in Low-Level Complexity}}.
Mathematical Foundations of Computer Science 1977,
Lecture Notes in Computer Science \textbf{53}, 1977, 162--176.

\bibitem[Val79]{Val79}
Leslie G. Valiant.
\href{https://doi.org/10.1137/0208032}
{\emph{The Complexity of Enumeration and Reliability Problems}}.
SIAM Journal on Computing \textbf{8}(3), 1979, 410--421.

\bibitem[VW19]{VW19}
Virginia Vassilevska Williams.
\href{https://doi.org/10.1142/9789813272880_0188}
{\emph{On Some Fine-Grained Questions in Algorithms and Complexity}}.
Proceedings of the International Congress of Mathematicians (ICM 2018),
World Scientific, 2019, 3447--3487.

\bibitem[VW21]{VW21}
Nikhil Vyas and Ryan Williams.
\href{https://doi.org/10.1613/jair.1.11859}
{\emph{On Super Strong ETH}}.
Journal of Artificial Intelligence Research \textbf{70}, 2021, 473--495.

\bibitem[Wah08]{Wahlstrom08}
Magnus Wahlstr\"om.
\href{https://doi.org/10.1007/978-3-540-79723-4_19}
{\emph{A Tighter Bound for Counting Max-Weight Solutions to 2SAT Instances}}.
3rd International Workshop on Parameterized and Exact Computation,
Lecture Notes in Computer Science \textbf{5018}, 2008, 202--213.

\bibitem[Wil05]{Wil05}
Ryan Williams.
\href{https://doi.org/10.1016/j.tcs.2005.09.023}
{\emph{A New Algorithm for Optimal 2-Constraint Satisfaction and Its Implications}}.
Theoretical Computer Science \textbf{348}(2--3), 2005, 357--365.

\bibitem[Wil13]{Wil13}
Ryan Williams.
\href{https://doi.org/10.1137/10080703X}
{\emph{Improving Exhaustive Search Implies Superpolynomial Lower Bounds}}.
SIAM Journal on Computing \textbf{42}(3), 2013, 1218--1244.
会议版本见 STOC 2010。引理~3.2 与定理~6.1 的编号依照
\href{https://people.csail.mit.edu/rrw/improved-algs-lbs2.pdf}{作者的完整版本}。

\bibitem[Wil14]{Wil14}
Richard Ryan Williams.
\href{https://doi.org/10.4230/LIPIcs.FSTTCS.2014.47}
{\emph{The Polynomial Method in Circuit Complexity Applied to Algorithm
Design (Invited Talk)}}.
34th International Conference on Foundation of Software Technology and
Theoretical Computer Science, LIPIcs \textbf{29}, 2014, 47--60.

\bibitem[Wil16]{Wil16}
Richard Ryan Williams.
\href{https://doi.org/10.4230/LIPIcs.CCC.2016.2}
{\emph{Strong ETH Breaks With Merlin and Arthur: Short Non-Interactive
Proofs of Batch Evaluation}}.
31st Conference on Computational Complexity, LIPIcs \textbf{50}, 2016,
2:1--2:17.

\bibitem[WY07]{WY07}
David Woodruff and Sergey Yekhanin.
\href{https://doi.org/10.1137/06065773X}
{\emph{A Geometric Approach to Information-Theoretic Private Information
Retrieval}}.
SIAM Journal on Computing \textbf{37}(4), 2007, 1046--1056.

\bibitem[\v{Z}\'ak83]{Zak83}
Stanislav \v{Z}\'ak.
\href{https://doi.org/10.1016/0304-3975(83)90015-4}{\emph{A Turing
Machine Time Hierarchy}}.
Theoretical Computer Science \textbf{26}(3), 1983, 327--333.

\end{thebibliography}
\end{document}
